% Applications of Definite Integrals — Further Mathematics Upper Sixth — Cours signé OnBuch+
% Official MINESEC syllabus. Compiler avec : tectonic FMA05-applications-of-definite-integrals.tex
\newcommand{\DOCMATIERE}{Further Mathematics}
\newcommand{\DOCNIVEAU}{Upper Sixth}
\newcommand{\DOCMODULE}{Upper Sixth — Further mathematics}
\newcommand{\DOCLECON}{5}
\newcommand{\DOCTITRE}{Applications of Definite Integrals}
% OnBuch+ — préambule « cours signés » ENRICHI (compilé avec tectonic / XeLaTeX)
% Système visuel « Pop » d'OnBuch+ : fond crème (jamais blanc), bordures encre
% épaisses, ombres « dures » décalées, titres Archivo Black en capitales.
% Version MATHÉMATIQUES : graphes (pgfplots/tikz), boîtes pédagogiques
% (définition, propriété, méthode, attention, le savais-tu, exercice résolu,
% exercice, TP, bilan), page de garde et signature OnBuch+ ancrée en bas.
%
% Macros attendues dans main.tex AVANT \input{preamble} :
%   \DOCMATIERE  \DOCNIVEAU  \DOCMODULE  \DOCLECON (numéro)  \DOCTITRE
\documentclass[11pt]{article}

\usepackage{fontspec}
\usepackage[english]{babel}
\usepackage[a4paper,margin=2cm,top=3.4cm,bottom=2.8cm,headheight=1.4cm,footskip=1.3cm]{geometry}
\usepackage{amsmath,amssymb,amsfonts}
\usepackage{mathtools} % \xmapsto, \coloneqq... souvent utilisés par les modèles
\usepackage{cancel} % simplifications barrées (bilans de forces)
% \abs{z} (module, valeur absolue) et \norm{u} (norme), utilisés par les modèles
\providecommand{\abs}{}\let\abs\relax\DeclarePairedDelimiter{\abs}{\lvert}{\rvert}
\providecommand{\norm}{}\let\norm\relax\DeclarePairedDelimiter{\norm}{\lVert}{\rVert}
\usepackage[table]{xcolor} % [table] : fournit aussi \rowcolors (alternance de lignes)
\usepackage{pagecolor}
\usepackage{tikz}
\usetikzlibrary{arrows.meta,positioning,calc,shapes.geometric,shapes.misc,shapes.arrows,decorations.pathmorphing,decorations.pathreplacing,decorations.markings,patterns,shadows,mindmap,trees,fit,backgrounds,matrix,intersections,angles,quotes}
\usepackage{pgfplots}
\usepackage[version=4]{mhchem} % équations nucléaires \ce{^{235}_{92}U}
\usepackage[european,siunitx]{circuitikz} % schémas électriques
\pgfplotsset{compat=1.18}
\usepgfplotslibrary{fillbetween,groupplots} % aires entre courbes (calcul intégral)
\usepackage{enumitem}
\usepackage{fancyhdr}
% [most] charge toutes les bibliothèques tcolorbox (listings, minted, raster,
% documentation...) et gonfle inutilement la mémoire TeX sur un document long
% (~300 boîtes) : on ne charge que ce qui est réellement utilisé.
\usepackage[skins,breakable]{tcolorbox}
\usepackage{graphicx}
\usepackage{colortbl}
\usepackage{booktabs}
\usepackage{tabularx}
\usepackage{diagbox} % case d'en-tête diagonale des tableaux à double entrée
% Colonnes extensibles centrée / gauche / droite (C, L, R), souvent utilisées par les modèles
\newcolumntype{C}{>{\centering\arraybackslash}X}
\newcolumntype{L}{>{\raggedright\arraybackslash}X}
\newcolumntype{R}{>{\raggedleft\arraybackslash}X}
\usepackage{array}
\usepackage{multicol}
\usepackage{siunitx}
% parse-numbers=false : accepte aussi \SI{2,5 \times 10^{-3}}{...}
\sisetup{locale=UK,output-decimal-marker={.},group-minimum-digits=4,parse-numbers=false}
\DeclareSIUnit\ohm{\ensuremath{\symup{\Omega}}} % U+2126 absent de la police
\DeclareSIUnit\division{div} % réglages d'oscilloscope (ms/div, V/div)
\DeclareSIUnit\tr{tr} % tours (vitesse de rotation, ex. \SI{1500}{\tr\per\minute})
\usepackage{qrcode}
% \degree : commande fréquemment utilisée par les modèles pour un angle ou une
% température en dehors de \SI{}{} (non fournie par les paquets chargés).
\providecommand{\degree}{^{\circ}}
% \qed (fin de démonstration), utilisé par les modèles sans amsthm
\providecommand{\qed}{\hfill$\square$}
% \barre : double barre des valeurs interdites dans un tableau de variations (tkz-tab)
\providecommand{\barre}{\ensuremath{\|}}
% environnement proof (amsthm non chargé) utilisé par les modèles
\ifdefined\proof\else\newenvironment{proof}[1][Proof]{\par\noindent\textit{#1.}\ }{\qed\par}\fi
\usepackage{lastpage}
\usepackage{hyperref}
\hypersetup{hidelinks}

% ============================================================
% Palette « Pop » OnBuch+ — mêmes hex que la classe P (plus_ui.dart)
% ============================================================
\definecolor{poporange}{HTML}{F59321}
\definecolor{popdark}{HTML}{E07A0C}
\definecolor{popcream}{HTML}{FFF6E8}
\definecolor{popink}{HTML}{1C1714}
\definecolor{poppurple}{HTML}{7A5AE0}
\definecolor{popgreen}{HTML}{2E9E6A}
\definecolor{poppink}{HTML}{E0457B}
\definecolor{popblue}{HTML}{2F7DE1}
\definecolor{popgold}{HTML}{C98A12}
\definecolor{popmuted}{HTML}{8A7F76}
\definecolor{popline}{HTML}{EADFD2}
\definecolor{popgray}{HTML}{8A7F76}
\colorlet{popgrey}{popgray}
% Alias tolérants (le modèle nomme parfois les couleurs différemment)
\colorlet{popwhite}{white}
\colorlet{popblack}{popink}
\colorlet{popred}{poppink}
\colorlet{popyellow}{popgold}
% Teintes claires pour fonds de tableaux / zones
\colorlet{popblueL}{popblue!10!white}
\colorlet{poporangeL}{poporange!14!white}
\colorlet{popgreenL}{popgreen!12!white}
\colorlet{poppurpleL}{poppurple!12!white}
\colorlet{poppinkL}{poppink!12!white}
\colorlet{popgoldL}{popgold!14!white}

\pagecolor{popcream}

% ============================================================
% Typographie
% ============================================================
\defaultfontfeatures{Ligatures=TeX}
\setmainfont{PlusJakartaSans-Medium.ttf}[Path=../../../fonts/,BoldFont=PlusJakartaSans-Bold.ttf]
% Maths en sans-serif (Fira Math) avec chiffres et lettres droites de Plus
% Jakarta, pour que les formules se fondent dans le texte.
\usepackage[math-style=ISO,bold-style=ISO]{unicode-math}
\setmathfont{FiraMath-Regular.otf}[Path=../../../fonts/]
\setmathfont{PlusJakartaSans-Medium.ttf}[Path=../../../fonts/,range={up/num,up/{Latin,latin},\mathup}]
% (icomma retiré : décimales avec point en anglais)
% Caractères mathématiques Unicode absents des polices de texte (Plus Jakarta,
% Archivo Black) : rendus par leur équivalent mathématique, en texte comme en
% formule — sinon ils s'affichent en carré vide (ex. « Arithmétique dans ℤ »).
\usepackage{newunicodechar}
\newunicodechar{ℝ}{\ensuremath{\mathbb{R}}}
\newunicodechar{ℤ}{\ensuremath{\mathbb{Z}}}
\newunicodechar{ℂ}{\ensuremath{\mathbb{C}}}
\newunicodechar{ℕ}{\ensuremath{\mathbb{N}}}
\newunicodechar{ℚ}{\ensuremath{\mathbb{Q}}}
\newunicodechar{𝔻}{\ensuremath{\mathbb{D}}}
\newunicodechar{α}{\ensuremath{\alpha}}
\newunicodechar{β}{\ensuremath{\beta}}
\newunicodechar{γ}{\ensuremath{\gamma}}
\newunicodechar{δ}{\ensuremath{\delta}}
\newunicodechar{ε}{\ensuremath{\varepsilon}}
\newunicodechar{θ}{\ensuremath{\theta}}
\newunicodechar{λ}{\ensuremath{\lambda}}
\newunicodechar{μ}{\ensuremath{\mu}}
\newunicodechar{σ}{\ensuremath{\sigma}}
\newunicodechar{τ}{\ensuremath{\tau}}
\newunicodechar{φ}{\ensuremath{\varphi}}
\newunicodechar{ψ}{\ensuremath{\psi}}
\newunicodechar{ω}{\ensuremath{\omega}}
\newunicodechar{Σ}{\ensuremath{\Sigma}}
% majuscules produites par \MakeUppercase dans les titres (α-aminés -> Α)
\newunicodechar{Α}{\ensuremath{\alpha}}
\newunicodechar{Β}{\ensuremath{\beta}}
\newunicodechar{Γ}{\ensuremath{\Gamma}}
\newunicodechar{Δ}{\ensuremath{\Delta}}
\newunicodechar{Ε}{\ensuremath{\varepsilon}}
\newunicodechar{Θ}{\ensuremath{\Theta}}
\newunicodechar{Λ}{\ensuremath{\Lambda}}
\newunicodechar{Μ}{\ensuremath{\mu}}
\newunicodechar{Τ}{\ensuremath{\tau}}
\newunicodechar{Φ}{\ensuremath{\Phi}}
\newunicodechar{Ψ}{\ensuremath{\Psi}}
\newunicodechar{Ω}{\ensuremath{\Omega}}
\newunicodechar{↦}{\ensuremath{\mapsto}}
\newunicodechar{∈}{\ensuremath{\in}}
\newunicodechar{∉}{\ensuremath{\notin}}
\newunicodechar{∧}{\ensuremath{\wedge}}
\newunicodechar{⇔}{\ensuremath{\Leftrightarrow}}
\newunicodechar{⟺}{\ensuremath{\Longleftrightarrow}}
\newunicodechar{⇒}{\ensuremath{\Rightarrow}}
\newunicodechar{⟹}{\ensuremath{\Longrightarrow}}
\newunicodechar{∘}{\ensuremath{\circ}}
\newunicodechar{∪}{\ensuremath{\cup}}
\newunicodechar{∩}{\ensuremath{\cap}}
\newunicodechar{≡}{\ensuremath{\equiv}}
\newunicodechar{⊂}{\ensuremath{\subset}}
\newunicodechar{⊥}{\ensuremath{\perp}}
\newunicodechar{∃}{\ensuremath{\exists}}
\newunicodechar{∀}{\ensuremath{\forall}}
\newunicodechar{∛}{\ensuremath{\sqrt[3]{\,}}}
\newunicodechar{∜}{\ensuremath{\sqrt[4]{\,}}}
\newunicodechar{⃗}{\ensuremath{{}^{\to}}}
\newunicodechar{ⁿ}{\ifmmode^{n}\else\textsuperscript{\ensuremath{n}}\fi}
\newunicodechar{ˣ}{\ifmmode^{x}\else\textsuperscript{\ensuremath{x}}\fi}
\newunicodechar{ᵇ}{\ifmmode^{b}\else\textsuperscript{\ensuremath{b}}\fi}
\newunicodechar{ʳ}{\ifmmode^{r}\else\textsuperscript{\ensuremath{r}}\fi}
\newunicodechar{ᵘ}{\ifmmode^{u}\else\textsuperscript{\ensuremath{u}}\fi}
\newunicodechar{ʸ}{\ifmmode^{y}\else\textsuperscript{\ensuremath{y}}\fi}
\newunicodechar{ᵃ}{\ifmmode^{a}\else\textsuperscript{\ensuremath{a}}\fi}
\newunicodechar{ᵉ}{\ifmmode^{e}\else\textsuperscript{\ensuremath{e}}\fi}
\newunicodechar{ᵏ}{\ifmmode^{k}\else\textsuperscript{\ensuremath{k}}\fi}
\newunicodechar{ᵖ}{\ifmmode^{p}\else\textsuperscript{\ensuremath{p}}\fi}
\newunicodechar{ᴺ}{\ifmmode^{N}\else\textsuperscript{\ensuremath{N}}\fi}
\newunicodechar{ᶜ}{\ifmmode^{c}\else\textsuperscript{\ensuremath{c}}\fi}
\newunicodechar{ʰ}{\ifmmode^{h}\else\textsuperscript{\ensuremath{h}}\fi}
\newunicodechar{ᵠ}{\ifmmode^{q}\else\textsuperscript{\ensuremath{q}}\fi}
\newunicodechar{ₙ}{\ifmmode_{n}\else\textsubscript{\ensuremath{n}}\fi}
\newunicodechar{ₐ}{\ifmmode_{a}\else\textsubscript{\ensuremath{a}}\fi}
\newunicodechar{ᵢ}{\ifmmode_{i}\else\textsubscript{\ensuremath{i}}\fi}
\newunicodechar{ᵣ}{\ifmmode_{r}\else\textsubscript{\ensuremath{r}}\fi}
\newunicodechar{ₖ}{\ifmmode_{k}\else\textsubscript{\ensuremath{k}}\fi}
\newunicodechar{ᵦ}{\ifmmode_{\beta}\else\textsubscript{\ensuremath{\beta}}\fi}
\newunicodechar{ₓ}{\ifmmode_{x}\else\textsubscript{\ensuremath{x}}\fi}
\newfontfamily\popdisplay{ArchivoBlack-Regular.ttf}[Path=../../../fonts/]
\newfontfamily\poplogo{SpaceGrotesk-Bold.ttf}[Path=../../../fonts/]
\newfontfamily\popsemi{PlusJakartaSans-SemiBold.ttf}[Path=../../../fonts/]
\newfontfamily\popxbold{PlusJakartaSans-ExtraBold.ttf}[Path=../../../fonts/]
\newfontfamily\popxboldls{PlusJakartaSans-ExtraBold.ttf}[Path=../../../fonts/,LetterSpace=3.0]

\setlist[enumerate]{leftmargin=1.7em,itemsep=5pt,topsep=4pt}
\setlist[itemize]{leftmargin=1.7em,itemsep=4pt,topsep=4pt,label={\color{poporange}\textbullet}}
\setlength{\parindent}{0pt}
\setlength{\parskip}{5pt}
\renewcommand{\arraystretch}{1.25}

% pgfplots : style Pop par défaut
\pgfplotsset{
  every axis/.append style={axis line style={popink,line width=1pt},tick style={popink},
    grid style={popline,line width=0.5pt},label style={font=\small},tick label style={font=\footnotesize}},
  popcurve/.style={poporange,line width=1.8pt,no markers,smooth},
}
% Même style disponible dans un \draw TikZ simple (hors pgfplots)
\tikzset{popcurve/.style={poporange,line width=1.8pt}}
\tikzset{popgrid/.style={popline,very thin}}
\colorlet{popcurve}{poporange} % le nom du style est parfois utilisé comme couleur

% ============================================================
% Pill / squiggle
% ============================================================
\newcommand{\poppill}[3]{%
  \tikz[baseline=-0.55ex]\node[draw=popink,line width=1.2pt,rounded corners=2.6pt,%
    fill=#2,inner xsep=6.5pt,inner ysep=3pt]{%
    \popxboldls\fontsize{7}{7}\selectfont\color{#3}\MakeUppercase{#1}};%
}
\newcommand{\popsquiggle}[1]{%
  \tikz[baseline=-0.4ex]{
    \draw[#1,line width=2.6pt,line cap=round]
      plot [smooth] coordinates {(0,0.09) (0.34,0.01) (0.68,0.21) (1.02,0.01) (1.36,0.10)};
  }%
}
% Mot-clé surligné (marqueur orange sous le texte)
\newcommand{\cle}[1]{{\popxbold\color{popdark}#1}}

% ============================================================
% En-tête / pied
% ============================================================
\pagestyle{fancy}
\fancyhf{}
\renewcommand{\headrulewidth}{0pt}
\renewcommand{\footrulewidth}{0pt}
\lhead{\raisebox{-0.15ex}{\poplogo\fontsize{13}{13}\selectfont\color{popink}OnBuch\color{poporange}+}}
\rhead{\poppill{\DOCMATIERE\ \textbullet\ \DOCNIVEAU\ \textbullet\ Chapter \DOCLECON}{white}{popink}}
\lfoot{\popxbold\fontsize{7.6}{7.6}\selectfont\color{popmuted}\MakeUppercase{Signed course OnBuch+}\ \textbullet\ {\popsemi\DOCTITRE}}
\rfoot{\tikz[baseline=-0.6ex]\node[rounded corners=8.5pt,fill=popink,minimum height=17pt,minimum width=17pt,inner xsep=5pt,inner sep=0]{\popxbold\fontsize{8}{8}\selectfont\color{popcream}\thepage\,/\,\pageref*{LastPage}};}

% ============================================================
% Titres
% ============================================================
\newcommand{\chaptitre}[2]{%
  \begin{tcolorbox}[enhanced,colback=poporange,colframe=popink,boxrule=2.6pt,arc=11pt,
      drop shadow=popink,left=15pt,right=15pt,top=12pt,bottom=11pt,before skip=3pt,after skip=18pt]
    \poppill{#2}{popink}{popcream}\par\vspace{9pt}
    {\raggedright\hyphenpenalty=10000\exhyphenpenalty=10000\popdisplay\fontsize{22}{24}\selectfont\color{popcream}\MakeUppercase{#1}\par}
  \end{tcolorbox}
}
% Section numérotée : pastille orange avec le numéro + titre Archivo
\newcounter{popsec}
\newcommand{\coursec}[1]{%
  \stepcounter{popsec}\setcounter{popsub}{0}%
  \par\vspace{16pt}\noindent
  \begin{minipage}{\linewidth}
  \tikz[baseline=-0.7ex]\node[draw=popink,line width=1.6pt,fill=poporange,rounded corners=4pt,minimum size=22pt,inner sep=0]{\popdisplay\fontsize{12}{12}\selectfont\color{popcream}\thepopsec};%
  \hspace{8pt}{\popdisplay\fontsize{15}{17}\selectfont\color{popink}\MakeUppercase{#1}}\par
  \vspace{4pt}\noindent{\color{poporange}\rule{36pt}{3.6pt}}
  \end{minipage}\par\nopagebreak\vspace{8pt}
}
\newcounter{popsub}
\newcommand{\courssub}[1]{%
  \stepcounter{popsub}%
  \par\vspace{9pt}\noindent{\popxbold\fontsize{11.5}{13}\selectfont\color{popink}{\color{poporange}\thepopsec.\thepopsub}\hspace{6pt}#1}\par\nopagebreak\vspace{4pt}
}

% ============================================================
% Boîtes pédagogiques — même recette : fond blanc, bordure épaisse colorée,
% ombre dure encre, badge plein. Titre optionnel [#1].
% ============================================================
\tcbset{popbase/.style={breakable,enhanced,colback=white,boxrule=2.3pt,arc=9pt,
  shadow={3pt}{-3pt}{0pt}{popink},left=14pt,right=14pt,top=11pt,bottom=11pt,before skip=11pt,after skip=15pt,
  fontupper=\popsemi\fontsize{10.6}{14.5}\selectfont\color{popink},
  fontlower=\popsemi\fontsize{10.6}{14.5}\selectfont\color{popink}}}
% Coche dessinée (le glyphe ✓ n'existe pas dans les polices du document)
\newcommand{\popcheck}[1]{\tikz[baseline=-0.5ex]\draw[#1,line width=1.6pt,line cap=round,line join=round](-0.13,0.0)--(-0.04,-0.09)--(0.13,0.1);}
\providecommand{\coche}{\popcheck{popgreen}}
\providecommand{\resultat}[1]{\par\smallskip\noindent\fcolorbox{popgreen}{popgreenL}{\parbox{\dimexpr\linewidth-2\fboxsep-2\fboxrule\relax}{\textbf{Result.} #1}}\par\smallskip}
\newcommand{\popboxhead}[3]{\poppill{#1}{#2}{white}\ifx\relax#3\relax\else\hspace{6pt}{\popxbold\fontsize{10.6}{12}\selectfont\color{popink}#3}\fi\par\vspace{7pt}}

\newtcolorbox{aretenir}[1][]{popbase,colframe=popgreen,before upper={\popboxhead{Key points}{popgreen}{#1}}}
\newtcolorbox{exemplebox}[1][]{popbase,colframe=poporange,before upper={\popboxhead{Exemple}{poporange}{#1}}}
\newtcolorbox{definition}[1][]{popbase,colframe=popblue,before upper={\popboxhead{Definition}{popblue}{#1}}}
\newtcolorbox{propriete}[1][]{popbase,colframe=poppurple,before upper={\popboxhead{Property}{poppurple}{#1}}}
\newtcolorbox{methode}[1][]{popbase,colframe=popgold,before upper={\popboxhead{Method}{popgold}{#1}}}
\newtcolorbox{attention}[1][]{popbase,colframe=poppink,before upper={\popboxhead{Warning}{poppink}{#1}}}
\newtcolorbox{experience}[1][]{popbase,colframe=popblue,colback=popblueL,before upper={\popboxhead{Experiment}{popblue}{#1}}}
% « Le savais-tu ? » : carte violette pleine (contexte, culture, Cameroun)
\newtcolorbox{savaistu}[1][]{popbase,colback=poppurple,colframe=popink,
  fontupper=\popsemi\fontsize{10.4}{14.2}\selectfont\color{white},
  before upper={\poppill{Did you know?}{white}{poppurple}\ifx\relax#1\relax\else\hspace{6pt}{\popxbold\color{white}#1}\fi\par\vspace{7pt}}}
% Exercice résolu : énoncé en haut, \tcblower puis la solution sur fond crème
\newcounter{popexo}\newcounter{popexr}
\newtcolorbox{exoresolu}[1][]{popbase,colframe=poporange,colbacklower=popcream,
  segmentation style={popink,line width=1.2pt,dashed},
  before upper={\stepcounter{popexr}\popboxhead{Worked example \thepopexr}{poporange}{#1}},
  before lower={\poppill{Solution}{popink}{popcream}\par\vspace{6pt}}}
% Exercice (non résolu en place) — la correction est dans la partie Corrigés
\newtcolorbox{exercice}[1][]{popbase,colframe=popink,
  before upper={\stepcounter{popexo}\popboxhead{Exercise \thepopexo}{popink}{#1}}}
% Corrigé
\newtcolorbox{corrige}[1][]{popbase,colframe=popgreen,colback=popgreenL,
  before upper={\popboxhead{Solution}{popgreen}{#1}}}
% Objectifs / prérequis / bilan
\newtcolorbox{objectifs}{popbase,colframe=popink,colback=poporangeL,
  before upper={\popboxhead{Chapter objectives}{popink}{}}}
\newtcolorbox{prerequis}{popbase,colframe=popink,colback=popblueL,
  before upper={\popboxhead{Prerequisites}{popblue}{}}}
\newtcolorbox{bilan}{popbase,colframe=popink,colback=white,boxrule=2.8pt,
  before upper={\popboxhead{Fiche bilan}{poporange}{L'essentiel à connaître pour l'examen}}}
% Figure encadrée avec légende
\newtcolorbox{popfigure}[1][]{enhanced,breakable=false,colback=white,colframe=popline,boxrule=1.4pt,arc=8pt,
  left=10pt,right=10pt,top=10pt,bottom=8pt,before skip=10pt,after skip=12pt,halign=center,
  lower separated=false,halign lower=center,
  fontlower=\popsemi\fontsize{9}{11.5}\selectfont\color{popmuted},
  #1}
\newcommand{\legende}[1]{\par\vspace{6pt}{\popsemi\fontsize{9}{11.5}\selectfont\color{popmuted}#1}}

% ============================================================
% Signature OnBuch+
% ============================================================
\newcommand{\poplogoinline}[1]{{\poplogo\fontsize{#1}{#1}\selectfont\color{popink}OnBuch\color{poporange}+}}
\newcommand{\signatureonbuch}{%
  \begin{tcolorbox}[enhanced,colback=popink,colframe=popink,boxrule=2.2pt,arc=10pt,
      drop shadow=poporange,left=14pt,right=14pt,top=10pt,bottom=10pt,before skip=0pt,after skip=0pt]
    \tikz[baseline=-0.7ex]\node[circle,fill=popgreen,draw=popcream,line width=1.3pt,minimum size=22pt,inner sep=0]{\popcheck{white}};%
    \hspace{9pt}\begin{minipage}[c]{0.62\linewidth}
      {\popxbold\fontsize{12}{13}\selectfont\color{popcream}Signed\ \textbullet\ {\poplogo OnBuch\color{poporange}+}}\par\vspace{2pt}
      {\raggedright\popsemi\fontsize{8}{10.5}\selectfont\color{popline}\MakeUppercase{OnBuch+ teaching team}\\\MakeUppercase{Follows the official MINESEC syllabus}\par}
    \end{minipage}\hfill
    {\poplogo\fontsize{22}{22}\selectfont\color{popcream}OnBuch\color{poporange}+}
  \end{tcolorbox}%
}

% ============================================================
% Page de garde
% #1 titre · #2 matière · #3 niveau · #4 module · #5 n° leçon · #6 durée ·
% #7 description / objectifs (texte ou itemize)
% ============================================================
\newcommand{\pagegarde}[7]{%
  \thispagestyle{empty}
  \newgeometry{a4paper,margin=1.7cm,top=1.6cm,bottom=1.4cm}%
  \begin{tikzpicture}[remember picture,overlay]
    \fill[poporange!18] ([xshift=-3.2cm,yshift=-3.6cm]current page.north east) circle (4.6cm);
    \fill[poppurple!14] ([xshift=2.4cm,yshift=7.6cm]current page.south west) circle (3.8cm);
  \end{tikzpicture}%
  {\poplogo\fontsize{30}{30}\selectfont\color{popink}OnBuch\color{poporange}+}\hfill
  \raisebox{0.6ex}{\poppill{Signed course \textbullet\ Premium}{popink}{popcream}}\par
  \vspace{1pt}\popsquiggle{poppurple}\par
  \vspace{8pt}
  \begin{tcolorbox}[enhanced,colback=poporange,colframe=popink,boxrule=2.8pt,arc=13pt,
      drop shadow=popink,left=18pt,right=18pt,top=12pt,bottom=13pt,after skip=10pt]
    \poppill{#2 \textbullet\ #3}{popink}{popcream}\hspace{5pt}\poppill{#4}{white}{popink}\par\vspace{8pt}
    {\popdisplay\fontsize{14}{14}\selectfont\color{popink}CHAPTER #5}\par\vspace{4pt}
    {\raggedright\hyphenpenalty=10000\exhyphenpenalty=10000\popdisplay\fontsize{25}{27}\selectfont\color{popcream}\MakeUppercase{#1}\par}
    \vspace{8pt}
    {\popxbold\fontsize{9.5}{9.5}\selectfont\color{popink}\MakeUppercase{Suggested time: #6}}
  \end{tcolorbox}
  \begin{tcolorbox}[enhanced,colback=white,colframe=popink,boxrule=2.2pt,arc=10pt,
      drop shadow=popink,left=16pt,right=16pt,top=10pt,bottom=10pt,after skip=8pt]
    \poppill{In this chapter}{poppurple}{white}\par\vspace{5pt}
    \popsemi\fontsize{9.3}{12.6}\selectfont\color{popink}#7
  \end{tcolorbox}
  \noindent\begin{minipage}[t]{0.487\linewidth}
    \begin{tcolorbox}[enhanced,colback=popgreen,colframe=popink,boxrule=2.2pt,arc=10pt,
        drop shadow=popink,left=11pt,right=11pt,top=10pt,bottom=10pt,height=3.1cm,valign=center]
      \tcbox[colback=white,colframe=popink,boxrule=1.3pt,arc=5pt,boxsep=0pt,nobeforeafter,
        left=3pt,right=3pt,top=3pt,bottom=3pt,valign=center]{\qrcode[height=1.9cm]{https://play.google.com/store/apps/details?id=com.luvvix.onbuch&hl=en}}%
      \hspace{8pt}\begin{minipage}[c]{0.5\linewidth}
        {\popdisplay\fontsize{11}{12.5}\selectfont\color{white}\MakeUppercase{Download OnBuch}}\par\vspace{4pt}
        {\raggedright\popsemi\fontsize{8.4}{11}\selectfont\color{white}Free \textbullet\ Google Play\\AI tutor, past papers, results\par}
      \end{minipage}
    \end{tcolorbox}
  \end{minipage}\hfill
  \begin{minipage}[t]{0.487\linewidth}
    \begin{tcolorbox}[enhanced,colback=poppurple,colframe=popink,boxrule=2.2pt,arc=10pt,
        drop shadow=popink,left=11pt,right=11pt,top=10pt,bottom=10pt,height=3.1cm,valign=center]
      \tikz[baseline=-0.7ex]\node[circle,fill=white,draw=popink,line width=1.3pt,minimum size=1.6cm,inner sep=0]{\popdisplay\fontsize{24}{24}\selectfont\color{poppurple}+};%
      \hspace{8pt}\begin{minipage}[c]{0.6\linewidth}
        {\popdisplay\fontsize{11}{12.5}\selectfont\color{white}\MakeUppercase{Go OnBuch+}}\par\vspace{4pt}
        {\raggedright\popsemi\fontsize{8.4}{11}\selectfont\color{white}All signed courses, unlimited AI tutor, PDF sheets — OnBuch+ tab in the app\par}
      \end{minipage}
    \end{tcolorbox}
  \end{minipage}
  \vspace{8pt}
  \popcontenu
  \vfill
  \signatureonbuch
  \restoregeometry
  \newpage
}

% Rangée « Ce cours contient » (page de garde)
\newcommand{\poptile}[3]{% #1 couleur #2 gros texte #3 légende
  \begin{tcolorbox}[enhanced,colback=#1,colframe=popink,boxrule=2pt,arc=9pt,shadow={2.5pt}{-2.5pt}{0pt}{popink},
      left=6pt,right=6pt,top=9pt,bottom=9pt,halign=center,valign=center,height=1.75cm,width=\linewidth,nobeforeafter]
    {\popdisplay\fontsize{12.5}{13}\selectfont\color{white}\MakeUppercase{#2}}\par\vspace{3pt}
    {\centering\popsemi\fontsize{7.8}{9.5}\selectfont\color{white}#3\par}
  \end{tcolorbox}}
\newcommand{\popcontenu}{%
  \par\noindent{\popxboldls\fontsize{7.5}{7.5}\selectfont\color{popmuted}THIS SIGNED COURSE CONTAINS}\par\vspace{7pt}
  \noindent\begin{minipage}[t]{0.235\linewidth}\poptile{popblue}{Course}{complete, illustrated, step by step}\end{minipage}\hfill
  \begin{minipage}[t]{0.235\linewidth}\poptile{poporange}{Methods}{and worked examples}\end{minipage}\hfill
  \begin{minipage}[t]{0.235\linewidth}\poptile{poppink}{Exercises}{exam style + solutions}\end{minipage}\hfill
  \begin{minipage}[t]{0.235\linewidth}\poptile{popgreen}{Summary sheet}{the essentials to remember}\end{minipage}\par}

% Sommaire
\newtcolorbox{sommaire}{enhanced,colback=white,colframe=popink,boxrule=2.2pt,arc=10pt,
  drop shadow=popink,left=18pt,right=18pt,top=13pt,bottom=13pt,before skip=6pt,after skip=20pt,
  before upper={\poppill{Contents}{poppurple}{white}\par\vspace{9pt}\popsemi\fontsize{10.6}{14.5}\selectfont\color{popink}}}

% Signature de fin de cours — poussée tout en bas de la dernière page
\newcommand{\findecours}{%
  \par\vfill\nopagebreak
  \begin{center}\popsquiggle{poporange}\end{center}\vspace{4pt}
  \signatureonbuch
}
\AtBeginDocument{\def\llbracket{\mathopen{[\mkern-3.5mu[}}\def\rrbracket{\mathclose{]\mkern-3.5mu]}}} % intervalles d'entiers (glyphe ⟦ absent de la police)
% Glyphes absents de Fira Math : redéfinis à partir de glyphes disponibles
\AtBeginDocument{%
  \def\setminus{\mathbin{\backslash}}\let\smallsetminus\setminus
  \def\vdots{\mathord{\vbox{\baselineskip3.2pt\lineskiplimit0pt\kern4pt\hbox{.}\hbox{.}\hbox{.}}}}%
  \def\ddots{\mathinner{\mkern1mu\raise6.5pt\hbox{.}\mkern2mu\raise3.6pt\hbox{.}\mkern2mu\raise0.7pt\hbox{.}\mkern1mu}}%
  \def\triangle{\mathord{\scalebox{0.9}{$\symup{\Delta}$}}}%
  \def\nabla{\mathord{\raisebox{\depth}{\scalebox{1}[-1]{$\symup{\Delta}$}}}}%
  \def\checkmark{\popcheck{popdark}}%
  \def\star{\mathbin{\ast}}\def\bigstar{\mathord{\ast}}%
  \def\diamond{\mathbin{\scalebox{0.6}{$\Diamond$}}}%
  \def\bigcup{\mathop{\mathchoice{\raisebox{-0.3em}{\scalebox{1.7}{$\cup$}}}{\scalebox{1.25}{$\cup$}}{\cup}{\cup}}\displaylimits}%
  \def\bigcap{\mathop{\mathchoice{\raisebox{-0.3em}{\scalebox{1.7}{$\cap$}}}{\scalebox{1.25}{$\cap$}}{\cap}{\cap}}\displaylimits}%
  \def\lesseqqgtr{\mathrel{\lessgtr}}\def\bigoplus{\mathop{\oplus}\displaylimits}%
}
\providecommand{\ssi}{if and only if} % abréviation fréquente
\AtBeginDocument{\providecommand{\wideparen}{\overparen}} % arc de cercle

%% FIGFIX-BEGIN v4 — correctifs globaux des figures (docs/onbuch-generation/tools/figfix)
\usepackage{adjustbox}\usepackage{etoolbox}
% toute figure TikZ/pgfplots plus large que la ligne est réduite à la largeur disponible
\BeforeBeginEnvironment{tikzpicture}{\begin{adjustbox}{max width=\linewidth}}
\AfterEndEnvironment{tikzpicture}{\end{adjustbox}}
% légendes pgfplots lisibles par-dessus les courbes
\pgfplotsset{every axis/.append style={legend style={fill=white,fill opacity=0.92,text opacity=1,draw=black!15,inner sep=3pt}}}
% étiquettes d'arêtes (syntaxe edge["..."]) avec fond blanc
\tikzset{every edge quotes/.append style={fill=white,inner sep=1.2pt}}
%% FIGFIX-END
\begin{document}

\pagegarde{Applications of Definite Integrals}{Further Mathematics}{Upper Sixth}{Upper Sixth — Further mathematics}{5}{7 periods}{This chapter develops the major geometric and physical applications of definite integrals: mean and RMS values, arc length, areas and volumes of revolution, centroids and centres of gravity, crowned by the elegant theorems of Pappus. It equips Upper Sixth students with the exact toolkit repeatedly examined at the GCE Advanced Level in Further Mathematics.
\vspace{4pt}
\begin{multicols}{2}\small
\begin{itemize}
\item By the end of this chapter I can compute the mean value of a function over an interval and interpret it geometrically and physically.
\item By the end of this chapter I can compute the root mean square (RMS) value of a function and apply it to alternating currents and other periodic quantities.
\item By the end of this chapter I can set up and evaluate the integral giving the arc length of a curve y = f(x), x = g(y), or a parametric curve.
\item By the end of this chapter I can compute the area of a surface of revolution about the x-axis or the y-axis.
\item By the end of this chapter I can compute volumes of solids of revolution by the disc and washer methods, about either axis.
\item By the end of this chapter I can locate the centroid of a plane region (including composite regions) using definite integrals.
\end{itemize}\end{multicols}}

\begin{sommaire}
\begin{multicols}{2}
\begin{enumerate}
\item Mean Value and RMS Value of a Function (Lesson 51)
\item Arc Length of a Curve (Lesson 52)
\item Surfaces and Volumes of Revolution (Lessons 53 \& 54)
\item Centroids of Plane Figures and Centres of Gravity of Solids of Revolution (Lessons 55 \& 56)
\item Theorems of Pappus and Synthesis (Lesson 57)
\item Integration activity
\item Methods and worked examples
\item Exercises
\item Solutions to the exercises
\item Summary sheet
\end{enumerate}
\end{multicols}
\end{sommaire}

\begin{objectifs}
\begin{itemize}
\item By the end of this chapter I can compute the mean value of a function over an interval and interpret it geometrically and physically.
\item By the end of this chapter I can compute the root mean square (RMS) value of a function and apply it to alternating currents and other periodic quantities.
\item By the end of this chapter I can set up and evaluate the integral giving the arc length of a curve y = f(x), x = g(y), or a parametric curve.
\item By the end of this chapter I can compute the area of a surface of revolution about the x-axis or the y-axis.
\item By the end of this chapter I can compute volumes of solids of revolution by the disc and washer methods, about either axis.
\item By the end of this chapter I can locate the centroid of a plane region (including composite regions) using definite integrals.
\item By the end of this chapter I can locate the centre of gravity of a solid of revolution by integration.
\item By the end of this chapter I can state and apply the two theorems of Pappus to find volumes, surface areas, and conversely centroids.
\item By the end of this chapter I can model and solve a concrete Cameroonian problem (tank, pot, cable, current) using several of these tools together.
\end{itemize}
\end{objectifs}

\begin{prerequis}
\begin{itemize}
\item Definite integration techniques from Lower Sixth and Term 1 of Upper Sixth: substitution, integration by parts, standard integrals (including those leading to ln and inverse trigonometric forms).
\item The area under a curve and between two curves as definite integrals (Lower Sixth).
\item Parametric equations of curves and differentiation of parametrically defined functions.
\item Trigonometric identities (especially 1 + tan²θ = sec²θ and double-angle formulae) and hyperbolic functions basics.
\item Moments of forces and the notion of centre of mass from Mechanics / Physics.
\end{itemize}
\end{prerequis}

\newpage

\coursec{Mean Value and RMS Value of a Function (Lesson 51)}

In Bamenda, a potter shapes a clay water pot on a spinning wheel: how much clay does she need, and how much glaze to cover its surface? At the SONARA refinery in Limbe, engineers must know the volume of a curved storage tank and where its centre of gravity lies so it does not topple when half-full. Even the ENEO engineer quoting ``220 volts'' is secretly using a \cle{root mean square} value of an alternating current. All these questions --- mean values, lengths of curves, areas, volumes and centres of gravity --- are answered by one powerful tool: the \cle{definite integral}. This chapter turns the integration techniques you already know into a measuring machine for the real world.

We begin with the simplest measuring question of all: \emph{what is the ``average height'' of a curve?}

\courssub{The Mean Value of a Function over an Interval}

\textbf{From averages of numbers to averages of functions.}\par
If a student scores $12$, $15$ and $9$ out of $20$ in three tests, her average is $\frac{12+15+9}{3}=12$. We add the values and divide by how many there are. But a function $f$ takes \emph{infinitely many} values on $[a,b]$, one for each $x$. We cannot add infinitely many numbers --- but the definite integral does exactly this: $\int_a^b f(x)\,dx$ is a ``continuous sum'' of the values of $f$, weighted by the tiny width $dx$. Dividing by the total length $b-a$ gives the natural continuous average.

\begin{definition}[Mean value of a function]
Let $f$ be continuous on $[a,b]$ with $a<b$. The \cle{mean value} (or average value) of $f$ over $[a,b]$ is
\[
\bar{f} \;=\; \frac{1}{b-a}\int_a^b f(x)\,dx.
\]
\end{definition}

\textbf{Geometric interpretation.}\par
Multiplying the definition by $(b-a)$ gives
\[
\bar{f}\,(b-a) = \int_a^b f(x)\,dx.
\]
The right-hand side is the area under the curve (when $f\ge 0$). The left-hand side is the area of a \emph{rectangle} of base $(b-a)$ and height $\bar{f}$. So $\bar{f}$ is the \cle{height of the rectangle} standing on $[a,b]$ whose area equals the area under the curve. If you levelled the ``hills and valleys'' of the graph into a flat surface without changing the total area, the flat level would be exactly $\bar{f}$.

\begin{popfigure}
\begin{center}
\begin{tikzpicture}
\begin{axis}[
  width=11cm, height=6.5cm,
  axis lines=middle,
  xlabel={$x$}, ylabel={$y$},
  xmin=-0.3, xmax=4.6, ymin=0, ymax=3.4,
  xtick={0.5,4}, xticklabels={$a$,$b$},
  ytick=\empty,
  domain=0.5:4, samples=100,
  every axis x label/.style={at={(ticklabel* cs:1)}, anchor=west},
]
\addplot[fill=popblueL, draw=none, domain=0.5:4] {1.9} \closedcycle;
\addplot[popblue, very thick] {1.6 + 0.9*sin(deg(1.6*x-0.8)) + 0.15*x};
\addplot[poporange, very thick, domain=0.5:4] {1.9};
\draw[dashed, popdark] (axis cs:0.5,0) -- (axis cs:0.5,1.9);
\draw[dashed, popdark] (axis cs:4,0) -- (axis cs:4,1.9);
\node[poporange, anchor=south west] at (axis cs:2.1,1.95) {$y=\bar{f}$};
\node[popblue] at (axis cs:3.4,2.9) {$y=f(x)$};
\end{axis}
\end{tikzpicture}
\end{center}
\legende{The mean value $\bar{f}$ is the height of the rectangle on $[a,b]$ having the same area as the region under the curve.}
\end{popfigure}

\textbf{The Mean Value Theorem for integrals.}\par
If $f$ is continuous on $[a,b]$, it attains every value between its minimum $m$ and maximum $M$. Since $m(b-a)\le \int_a^b f(x)\,dx \le M(b-a)$, dividing by $(b-a)$ gives $m \le \bar{f} \le M$, and the Intermediate Value Theorem then guarantees the following result.

\begin{propriete}[Mean Value Theorem for integrals]
If $f$ is continuous on $[a,b]$, there exists at least one point $c\in[a,b]$ such that
\[
f(c)=\bar{f}=\frac{1}{b-a}\int_a^b f(x)\,dx,
\qquad\text{i.e.}\qquad
\int_a^b f(x)\,dx = f(c)\,(b-a).
\]
\emph{Proof not examinable; the statement and its use are.}
\end{propriete}

In words: a continuous curve actually \emph{passes through} its own average height somewhere on the interval.

\courssub{Mean Values of Standard Functions}

\begin{exemplebox}[{Mean value of $\sin x$ over $[0,\pi]$ and $[0,2\pi]$}]
Over $[0,\pi]$:
\[
\bar{f}=\frac{1}{\pi-0}\int_0^{\pi}\sin x\,dx
=\frac{1}{\pi}\big[-\cos x\big]_0^{\pi}
=\frac{1}{\pi}\big(1+1\big)=\frac{2}{\pi}\approx 0.637.
\]
Over $[0,2\pi]$:
\[
\bar{f}=\frac{1}{2\pi}\int_0^{2\pi}\sin x\,dx
=\frac{1}{2\pi}\big[-\cos x\big]_0^{2\pi}
=\frac{1}{2\pi}( -1+1)=0.
\]
The second answer is $0$ because the positive arch on $[0,\pi]$ is exactly cancelled by the negative arch on $[\pi,2\pi]$: the integral counts \emph{signed} area. The same results hold for $\cos x$ over the corresponding intervals.
\end{exemplebox}

\begin{propriete}[Mean of $\sin^2$ and $\cos^2$ over a full period]
Over any interval of length one period (e.g. $[0,\pi]$ for $\sin^2 x$, or $[0,2\pi]$ for either function),
\[
\frac{1}{b-a}\int_a^b \sin^2 x\,dx=\frac12,
\qquad
\frac{1}{b-a}\int_a^b \cos^2 x\,dx=\frac12.
\]
\emph{Proof (examinable technique).} Using $\cos^2x=\frac{1+\cos 2x}{2}$:
\[
\int_0^{2\pi}\cos^2x\,dx
=\int_0^{2\pi}\frac{1+\cos 2x}{2}\,dx
=\Big[\frac{x}{2}+\frac{\sin 2x}{4}\Big]_0^{2\pi}
=\pi,
\]
so the mean is $\dfrac{\pi}{2\pi}=\dfrac12$. The same works for $\sin^2x=\frac{1-\cos 2x}{2}$. Alternatively, note that $\sin^2x+\cos^2x=1$ forces the two means to add to $1$, and by symmetry they are equal, so each is $\tfrac12$.
\end{propriete}

\begin{attention}[The mean of $f^2$ is NOT the square of the mean of $f$]
A classic GCE trap: $\overline{f^2}\neq (\bar f)^2$ in general. For example, $\sin x$ has mean $0$ over $[0,2\pi]$, but $\sin^2 x$ has mean $\tfrac12$, not $0^2=0$. Squaring \emph{before} averaging and averaging \emph{before} squaring are completely different operations. Always follow the order given by the definition.
\end{attention}

\courssub{The Root Mean Square (RMS) Value}

\textbf{Why square first?}\par
The alternating current in a Cameroonian home changes direction $50$ times per second: its ordinary mean over a period is $0$, which would wrongly suggest ``no current''. Yet the current heats a pressing iron and lights a bulb. The physical effect (power dissipated in a resistor, $P=Ri^2$) depends on $i^2$, which is always $\ge 0$. The right ``effective value'' is therefore: \emph{square, average, then take the square root} to return to the original units.

\begin{definition}[Root mean square value]
Let $f$ be continuous on $[a,b]$. The \cle{root mean square} (RMS) value of $f$ over $[a,b]$ is
\[
f_{\mathrm{rms}}=\sqrt{\frac{1}{b-a}\int_a^b \big[f(x)\big]^2\,dx}.
\]
\end{definition}

\begin{propriete}[RMS of a sinusoid]
For $f(t)=A\sin(\omega t)$ over one period $T=\dfrac{2\pi}{\omega}$,
\[
f_{\mathrm{rms}}=\frac{A}{\sqrt{2}}.
\]
\emph{Proof (examinable).} Using $\sin^2(\omega t)=\frac{1-\cos(2\omega t)}{2}$:
\[
\int_0^{T} A^2\sin^2(\omega t)\,dt
=\frac{A^2}{2}\int_0^{T}\big(1-\cos(2\omega t)\big)\,dt
=\frac{A^2}{2}\Big[t-\frac{\sin(2\omega t)}{2\omega}\Big]_0^{T}
=\frac{A^2 T}{2},
\]
since $\sin(2\omega T)=\sin(4\pi)=0$. Hence
\[
f_{\mathrm{rms}}=\sqrt{\frac{1}{T}\cdot\frac{A^2T}{2}}=\frac{A}{\sqrt{2}}.
\]
The same result holds for $A\cos(\omega t)$, and more generally for $A\sin(\omega t+\varphi)$: the phase shift does not change the RMS value.
\end{propriete}

\begin{popfigure}
\begin{center}
\begin{tikzpicture}
\begin{axis}[
  width=12cm, height=7cm,
  axis lines=middle,
  xlabel={$t$}, ylabel={},
  xmin=-0.15, xmax=6.7, ymin=-1.5, ymax=1.6,
  xtick={6.283}, xticklabels={$T$},
  ytick={1,0.707,0.5}, yticklabels={$I_0$,$I_0/\sqrt{2}$,$I_0^2/2$},
  domain=0:6.283, samples=150,
  every axis x label/.style={at={(ticklabel* cs:1)}, anchor=west},
]
\addplot[popblue, very thick] {sin(deg(x))};
\addplot[poporange, very thick] {(sin(deg(x)))^2};
\addplot[dashed, popdark, domain=0:6.283] {0.5};
\addplot[dashed, popgreen, very thick, domain=0:6.283] {0.707};
\node[popblue] at (axis cs:1.9,1.25) {$i=I_0\sin(\omega t)$};
\node[poporange] at (axis cs:4.4,1.15) {$i^2$};
\node[popdark, anchor=south] at (axis cs:3.1,0.5) {mean of $i^2$};
\node[popgreen, anchor=north west] at (axis cs:3.6,0.72) {$i_{\mathrm{rms}}$};
\end{axis}
\end{tikzpicture}
\end{center}
\legende{Over one period, $i$ has mean $0$ but $i^2$ has mean $I_0^2/2$; the RMS value is $I_0/\sqrt{2}\approx 0.707\,I_0$.}
\end{popfigure}

\begin{savaistu}[Why ENEO says ``220 volts'']
The mains supply in Cameroon is alternating: $u(t)=U_0\sin(\omega t)$ with frequency $50\ \mathrm{Hz}$. The quoted ``$220\ \mathrm{V}$'' is the \emph{RMS} value, not the peak. The peak voltage is $U_0=220\sqrt{2}\approx 311\ \mathrm{V}$. A $220\ \mathrm{V}$ AC supply heats a resistor exactly as a $220\ \mathrm{V}$ DC battery would --- that is precisely what the RMS value is designed to mean.
\end{savaistu}

\courssub{Method and Worked Examples}

\begin{methode}[Computing a mean or RMS value]
\begin{enumerate}
\item \textbf{Identify the interval} $[a,b]$ --- for a periodic function, take exactly one period (use the graph or the equation to find it).
\item \textbf{Quote the formula}: $\bar f=\frac{1}{b-a}\int_a^b f\,dx$ or $f_{\mathrm{rms}}=\sqrt{\frac{1}{b-a}\int_a^b f^2\,dx}$.
\item \textbf{Set up the integral}, splitting at points where the formula for $f$ changes (piecewise functions) or using symmetry to shorten the work.
\item \textbf{Integrate} (linearise $\sin^2$, $\cos^2$ with double-angle identities).
\item \textbf{Divide by $(b-a)$}; for RMS, finish by taking the \textbf{square root}.
\end{enumerate}
\end{methode}

\begin{exoresolu}[Mean and RMS of a rectified current]
A half-wave rectified current is given over one period $[0,2\pi]$ by
\[
i(t)=
\begin{cases}
I_0\sin t, & 0\le t\le \pi,\\[2pt]
0, & \pi<t\le 2\pi.
\end{cases}
\]
Find (a) the mean value and (b) the RMS value of $i$ over $[0,2\pi]$.
\tcblower
\textbf{Solution.} The interval has length $2\pi$; the function is zero on $[\pi,2\pi]$, so only $[0,\pi]$ contributes.

(a) Mean value:
\[
\bar\imath=\frac{1}{2\pi}\int_0^{2\pi} i(t)\,dt
=\frac{1}{2\pi}\int_0^{\pi} I_0\sin t\,dt
=\frac{I_0}{2\pi}\big[-\cos t\big]_0^{\pi}
=\frac{I_0}{2\pi}(1+1)=\frac{I_0}{\pi}.
\]

(b) RMS value:
\[
i_{\mathrm{rms}}^2=\frac{1}{2\pi}\int_0^{\pi} I_0^2\sin^2 t\,dt
=\frac{I_0^2}{2\pi}\int_0^{\pi}\frac{1-\cos 2t}{2}\,dt
=\frac{I_0^2}{4\pi}\Big[t-\frac{\sin 2t}{2}\Big]_0^{\pi}
=\frac{I_0^2}{4\pi}\cdot\pi=\frac{I_0^2}{4}.
\]
Hence $i_{\mathrm{rms}}=\dfrac{I_0}{2}$.

\emph{Remark.} Note the difference with the full sinusoid: the mean over the whole period is halved ($I_0/\pi$ instead of $2I_0/\pi$ over the conducting half-period) and the RMS is $I_0/2$ instead of $I_0/\sqrt2$, because the current ``rests'' for half of each period.
\end{exoresolu}

\begin{exoresolu}[Mean value of a sawtooth function]
A voltage rises linearly from $0$ to $12\ \mathrm{V}$ over $4\ \mathrm{ms}$, then drops instantly to $0$ and repeats: $u(t)=3t$ for $t\in[0,4]$ (in ms, V). Find the mean and RMS values over one period.
\tcblower
\textbf{Solution.} Period $T=4$.

Mean:
\[
\bar u=\frac{1}{4}\int_0^4 3t\,dt=\frac{1}{4}\Big[\frac{3t^2}{2}\Big]_0^4=\frac{1}{4}\times 24=6\ \mathrm{V}.
\]
(Geometrically, the graph is a triangle of height $12$: its average height is half the peak.)

RMS:
\[
u_{\mathrm{rms}}^2=\frac{1}{4}\int_0^4 9t^2\,dt=\frac{1}{4}\Big[3t^3\Big]_0^4=\frac{1}{4}\times 192=48,
\qquad
u_{\mathrm{rms}}=\sqrt{48}=4\sqrt{3}\approx 6.93\ \mathrm{V}.
\]
Observe that $u_{\mathrm{rms}}>\bar u$: for a non-constant positive function the RMS value always exceeds the mean, because squaring amplifies the larger values (equivalently, $\overline{f^2}-(\bar f)^2=\overline{(f-\bar f)^2}\ge 0$).
\end{exoresolu}

\courssub{GCE Examination Technique}

\begin{attention}[Exam tips for mean and RMS questions]
\begin{itemize}
\item \textbf{Quote the formula before substituting.} Examiners award marks for the correct integral set-up $\frac{1}{b-a}\int_a^b f(x)\,dx$ even if a later slip occurs.
\item \textbf{Choose the limits carefully.} For a periodic function, integrate over exactly \emph{one} period: find $T=\frac{2\pi}{\omega}$ first and write it down.
\item \textbf{Use symmetry.} If $f$ is even, integrate over $[0,a]$ and double; if the question concerns $\sin^2$ or $\cos^2$ over a whole number of periods, you may quote the mean $\tfrac12$ directly (with a one-line justification).
\item \textbf{Do not forget the final square root} for RMS --- and do not forget to divide by $(b-a)$ before rooting.
\item \textbf{Check plausibility}: for a positive function, you should obtain $\bar f \le f_{\mathrm{rms}} \le \max f$.
\end{itemize}
\end{attention}

\begin{exercice}[Mean and RMS practice]
\begin{enumerate}
\item Find the mean value of $f(x)=x^2$ over $[0,3]$, and the point $c\in[0,3]$ where $f(c)=\bar f$.
\item Show that the RMS value of $u(t)=U_0\cos(\omega t)$ over one period is $U_0/\sqrt2$.
\item A current is defined on $[0,2]$ by $i(t)=4t$ for $0\le t\le 1$ and $i(t)=4$ for $1<t\le 2$ (in amperes, $t$ in seconds). Find its mean and RMS values over $[0,2]$.
\end{enumerate}
\end{exercice}

\begin{corrige}[Exercise 1]
\textbf{1.} $\bar f=\frac{1}{3}\int_0^3 x^2\,dx=\frac{1}{3}\cdot 9=3$. Then $c^2=3$ gives $c=\sqrt3\in[0,3]$, confirming the Mean Value Theorem for integrals.

\textbf{2.} Over $T=\frac{2\pi}{\omega}$: $\int_0^T U_0^2\cos^2(\omega t)\,dt=\frac{U_0^2}{2}\int_0^T(1+\cos 2\omega t)\,dt=\frac{U_0^2 T}{2}$, so $u_{\mathrm{rms}}=\sqrt{\frac{U_0^2}{2}}=\frac{U_0}{\sqrt2}$.

\textbf{3.} Mean: $\bar\imath=\frac12\left(\int_0^1 4t\,dt+\int_1^2 4\,dt\right)=\frac12(2+4)=3\ \mathrm{A}$.
RMS: $i_{\mathrm{rms}}^2=\frac12\left(\int_0^1 16t^2\,dt+\int_1^2 16\,dt\right)=\frac12\left(\frac{16}{3}+16\right)=\frac{32}{3}$, so $i_{\mathrm{rms}}=\sqrt{\frac{32}{3}}=\frac{4\sqrt6}{3}\approx 3.27\ \mathrm{A}$.
\end{corrige}

\begin{aretenir}[Key points of Lesson 51]
\begin{itemize}
\item Mean value: $\bar f=\dfrac{1}{b-a}\displaystyle\int_a^b f(x)\,dx$ --- the height of the equal-area rectangle.
\item If $f$ is continuous, there exists $c\in[a,b]$ with $f(c)=\bar f$ (Mean Value Theorem for integrals).
\item Over a full period: mean of $\sin$ or $\cos$ is $0$; mean of $\sin^2$ or $\cos^2$ is $\tfrac12$.
\item RMS value: $f_{\mathrm{rms}}=\sqrt{\dfrac{1}{b-a}\displaystyle\int_a^b [f(x)]^2\,dx}$; for $A\sin(\omega t)$, $f_{\mathrm{rms}}=\dfrac{A}{\sqrt2}$.
\item Squaring removes sign problems: that is why RMS, not the mean, measures the ``effective size'' of an alternating quantity (e.g. the $220\ \mathrm{V}$ mains, peak $\approx 311\ \mathrm{V}$).
\item Never confuse $\overline{f^2}$ with $(\bar f)^2$.
\end{itemize}
\end{aretenir}

\coursec{Arc Length of a Curve (Lesson 52)}

In the previous section we used the definite integral to attach a single number to a function over an interval: its mean value and its root mean square value. In both cases the integral was a natural average of the values of $f(x)$. We now turn to a different kind of question, one that is geometric rather than statistical: \emph{how long is a curve?}

Think of a surveyor in Bamenda who must measure the length of a winding road on a map, or an engineer of the electricity company who must know how much cable is needed to hang between two poles: the cable is not straight, it sags. A ruler measures straight distances only. To measure a curved path we need a new idea, and once again the definite integral provides it.

\courssub{From Polygonal Approximations to the Arc Length Element}

\textbf{The intuitive idea.} Suppose we want the length of a curve $\mathcal{C}$ which is the graph of a smooth function $y=f(x)$ between the points $A=(a,f(a))$ and $B=(b,f(b))$. The strategy is the same one Archimedes used to measure the circumference of a circle: replace the curve by a broken line (a polygon) made of many short straight chords, measure the chords, add them up, and then let the chords become shorter and shorter.

Divide $[a,b]$ into $n$ small subintervals by points $a=x_0<x_1<\dots<x_n=b$, and let $P_k=(x_k,f(x_k))$ be the corresponding points on the curve. The polygonal line $P_0P_1\cdots P_n$ has total length
\[
L_n=\sum_{k=1}^{n} |P_{k-1}P_k|
=\sum_{k=1}^{n}\sqrt{(x_k-x_{k-1})^2+\bigl(f(x_k)-f(x_{k-1})\bigr)^2}.
\]
Writing $\Delta x_k=x_k-x_{k-1}$ and $\Delta y_k=f(x_k)-f(x_{k-1})$, each chord has length $\sqrt{(\Delta x_k)^2+(\Delta y_k)^2}$. By the Mean Value Theorem, since $f$ is differentiable, there exists $c_k\in(x_{k-1},x_k)$ with $\Delta y_k=f'(c_k)\,\Delta x_k$. Hence
\[
L_n=\sum_{k=1}^{n}\sqrt{1+\bigl(f'(c_k)\bigr)^2}\;\Delta x_k.
\]
As $n\to\infty$ and every $\Delta x_k\to 0$, this Riemann sum converges to a definite integral, and the polygonal line converges to the curve. The true length of the curve is therefore
\[
s=\int_a^b \sqrt{1+\bigl(f'(x)\bigr)^2}\,dx.
\]

\textbf{The infinitesimal point of view.} It is very useful to think of a tiny piece of the curve as the hypotenuse of a tiny right-angled triangle with horizontal side $dx$ and vertical side $dy$. Pythagoras' theorem then gives the \cle{arc length element}
\[
ds=\sqrt{dx^2+dy^2},
\]
and the total length is the ``sum'' (integral) of all these elements: $s=\displaystyle\int_{\mathcal C} ds$. Depending on how the curve is described, we factor $dx$, $dy$ or $dt$ out of the square root, and this gives the three working formulae below.

\begin{popfigure}
\begin{center}
\begin{tikzpicture}[scale=1.05]
% curve
\draw[very thick,popblue] plot[smooth] coordinates {(0,0.4) (1,0.9) (2,1.1) (3,1.7) (4,2.0) (5,2.9) (6,3.3)};
% chords
\draw[poporange,thick,dashed] (0,0.4) -- (1,0.9) -- (2,1.1) -- (3,1.7) -- (4,2.0) -- (5,2.9) -- (6,3.3);
\foreach \p in {(0,0.4),(1,0.9),(2,1.1),(3,1.7),(4,2.0),(5,2.9),(6,3.3)} \fill[popdark] \p circle (1.6pt);
\node[popblue] at (3.4,2.6) {$\mathcal{C}$};
\node[poporange] at (0.7,0.35) {chords};
% magnified element
\draw[thick,popdark] (7.6,0.6) -- (10.2,0.6) -- (10.2,2.4) -- cycle;
\draw[very thick,popblue] (7.6,0.6) .. controls (8.6,0.75) and (9.4,1.35) .. (10.2,2.4);
\draw[<->,popmuted] (7.6,0.35) -- (10.2,0.35) node[midway,below] {$dx$};
\draw[<->,popmuted] (10.5,0.6) -- (10.5,2.4) node[midway,right] {$dy$};
\node[popblue] at (8.7,2.0) {$ds$};
\node[align=center] at (8.9,-0.5) {$ds=\sqrt{dx^2+dy^2}$};
\end{tikzpicture}
\end{center}
\legende{A curve approximated by chords (left) and one magnified arc element (right): the tiny arc is the hypotenuse of a right triangle with sides $dx$ and $dy$.}
\end{popfigure}

\begin{definition}[Arc length formulae]
Let $\mathcal{C}$ be a smooth curve (continuously differentiable, with no corners) joining two points. Its \cle{arc length} $s$ is obtained by integrating the element $ds=\sqrt{dx^2+dy^2}$:
\begin{enumerate}
\item \textbf{Cartesian form in $x$.} If $\mathcal{C}$ is the graph $y=f(x)$, $a\le x\le b$:
\[
s=\int_a^b\sqrt{1+\left(\frac{dy}{dx}\right)^2}\,dx.
\]
\item \textbf{Cartesian form in $y$.} If $\mathcal{C}$ is given as $x=g(y)$, $c\le y\le d$:
\[
s=\int_c^d\sqrt{1+\left(\frac{dx}{dy}\right)^2}\,dy.
\]
\item \textbf{Parametric form.} If $\mathcal{C}$ is given by $x=x(t)$, $y=y(t)$, $t_1\le t\le t_2$, traced exactly once:
\[
s=\int_{t_1}^{t_2}\sqrt{\left(\frac{dx}{dt}\right)^2+\left(\frac{dy}{dt}\right)^2}\,dt.
\]
\end{enumerate}
In each case the integrand is simply $ds$ with the appropriate variable factored out: $ds=\sqrt{1+(y')^2}\,dx=\sqrt{1+(x')^2}\,dy=\sqrt{\dot x^2+\dot y^2}\,dt$.
\end{definition}

\begin{methode}[The 4-step checklist for arc length]
\begin{enumerate}
\item \textbf{Compute the derivative} of the given equation: $y'=\dfrac{dy}{dx}$, or $x'=\dfrac{dx}{dy}$, or the pair $\dot x,\dot y$.
\item \textbf{Form $1+(\text{derivative})^2$} (or $\dot x^2+\dot y^2$) and simplify it algebraically.
\item \textbf{Look for a perfect square.} In GCE questions the expression almost always becomes a square, so that the square root disappears. If it does not, re-read the question.
\item \textbf{Integrate between the correct limits}, making sure the limits belong to the variable of integration.
\end{enumerate}
\end{methode}

\courssub{Worked Examples in Cartesian Form}

\begin{exoresolu}[Length of $y=x^{3/2}$]
Find the length of the curve $y=x^{3/2}$ from $x=0$ to $x=4$.
\tcblower
\textbf{Solution.} Step 1: $\dfrac{dy}{dx}=\dfrac{3}{2}x^{1/2}$.

Step 2: $1+\left(\dfrac{dy}{dx}\right)^2=1+\dfrac{9}{4}x$.

Step 3: this is not a perfect square, but it integrates directly by substitution (or by inspection):
\[
s=\int_0^4\sqrt{1+\tfrac94x}\,dx
=\left[\frac{2}{3}\cdot\frac{4}{9}\left(1+\tfrac94x\right)^{3/2}\right]_0^4
=\frac{8}{27}\left(10^{3/2}-1\right).
\]
Step 4: numerically, $s=\dfrac{8}{27}\left(10\sqrt{10}-1\right)\approx 9.07$ units.
\end{exoresolu}

\begin{exoresolu}[The catenary: length of a hanging cable]
A cable hanging between two poles takes the shape of the \cle{catenary} $y=\cosh x=\dfrac{e^x+e^{-x}}{2}$. Find the length of the cable from $x=-a$ to $x=a$.
\tcblower
\textbf{Solution.} Step 1: $y'=\sinh x$.

Step 2: $1+(y')^2=1+\sinh^2 x$.

Step 3: the identity $\cosh^2x-\sinh^2x=1$ gives the perfect square $1+\sinh^2x=\cosh^2x$. Since $\cosh x>0$:
\[
s=\int_{-a}^{a}\cosh x\,dx=\bigl[\sinh x\bigr]_{-a}^{a}=2\sinh a.
\]
For example, if the poles are at $x=\pm1$, the cable length is $2\sinh 1\approx 2.35$ units, slightly more than the straight-line distance $2$ — exactly the ``slack'' one observes on a real cable.
\end{exoresolu}

\begin{popfigure}
\begin{center}
\begin{tikzpicture}
\begin{axis}[
  width=11cm, height=6.5cm,
  axis lines=middle,
  xlabel={$x$}, ylabel={$y$},
  xmin=-2.4, xmax=2.4, ymin=0, ymax=4.4,
  xtick={-1.5,1.5}, xticklabels={$-a$,$a$},
  ytick=\empty,
  domain=-2.2:2.2, samples=80,
  clip=false]
\addplot[very thick,popblue] {cosh(x)};
% poles
\draw[very thick,popdark] (axis cs:-1.5,0) -- (axis cs:-1.5,3.09);
\draw[very thick,popdark] (axis cs:1.5,0) -- (axis cs:1.5,3.09);
\fill[poporange] (axis cs:-1.5,2.35) circle (2pt);
\fill[poporange] (axis cs:1.5,2.35) circle (2pt);
\node[popblue] at (axis cs:0,0.65) {$y=\cosh x$};
\node[align=center,popdark] at (axis cs:0,3.9) {length $=2\sinh a$};
\end{axis}
\end{tikzpicture}
\end{center}
\legende{The catenary $y=\cosh x$ hanging between two poles at $x=\pm a$: its length is $2\sinh a$.}
\end{popfigure}

\begin{exoresolu}[Length of $y=\ln(\sec x)$]
Find the length of the curve $y=\ln(\sec x)$ from $x=0$ to $x=\dfrac{\pi}{3}$.
\tcblower
\textbf{Solution.} Step 1: $y'=\dfrac{\sec x\tan x}{\sec x}=\tan x$.

Step 2--3: $1+\tan^2x=\sec^2x$, a perfect square, and $\sec x>0$ on $\left[0,\tfrac{\pi}{3}\right]$:
\[
s=\int_0^{\pi/3}\sec x\,dx=\Bigl[\ln|\sec x+\tan x|\Bigr]_0^{\pi/3}
=\ln(2+\sqrt3)-\ln 1=\ln(2+\sqrt3)\approx 1.32.
\]
\end{exoresolu}

\textbf{The key algebraic trick: the perfect square.} The three examples above succeed because $1+(y')^2$ collapses to a square. A classic family of examination curves is built precisely so that this happens. Consider
\[
y=\frac{x^2}{4}-\frac{1}{2}\ln x \qquad (x>0).
\]
Then $y'=\dfrac{x}{2}-\dfrac{1}{2x}$, and
\[
1+(y')^2=1+\left(\frac{x}{2}-\frac{1}{2x}\right)^2
=1+\frac{x^2}{4}-\frac12+\frac{1}{4x^2}
=\frac{x^2}{4}+\frac12+\frac{1}{4x^2}
=\left(\frac{x}{2}+\frac{1}{2x}\right)^2.
\]
Since $\dfrac{x}{2}+\dfrac{1}{2x}>0$, the square root is immediate:
\[
s=\int_1^e\left(\frac{x}{2}+\frac{1}{2x}\right)dx
=\left[\frac{x^2}{4}+\frac12\ln x\right]_1^e
=\frac{e^2}{4}+\frac12-\frac14=\frac{e^2+1}{4}.
\]

\begin{attention}[Exam tip: the perfect square is by design]
If your expression $1+(y')^2$ does \emph{not} simplify to a perfect square (or to a simple linear expression under the root), stop and re-read the question. GCE questions are designed so that it simplifies. A very common pattern is
\[
1+\left(\frac{u}{2}-\frac{1}{2u}\right)^2=\left(\frac{u}{2}+\frac{1}{2u}\right)^2,
\]
because the cross term $-2\cdot\frac{u}{2}\cdot\frac{1}{2u}=-\frac12$ combines with the $+1$ to give $+\frac12$. Watch for it whenever $y'$ is a difference of reciprocal-type terms.
\end{attention}

\courssub{Curves Given as $x=g(y)$}

Sometimes differentiating with respect to $y$ is far easier. For instance the curve $y^2=4x$ from $(0,0)$ to $(1,2)$ is awkward as $y=2\sqrt{x}$ near $x=0$ (the derivative blows up), but as $x=\dfrac{y^2}{4}$ it is perfectly smooth:
\[
\frac{dx}{dy}=\frac{y}{2},\qquad
s=\int_0^2\sqrt{1+\frac{y^2}{4}}\;dy.
\]
This integral is handled by the substitution $y=2\sinh u$ (or quoted from the standard $\sqrt{1+t^2}$ formula), giving
\[
s=\sqrt{2}+\ln(1+\sqrt2)\approx 2.30.
\]
The lesson: \emph{choose the form of the formula that makes the derivative and the integral simplest.}

\courssub{Parametric Curves}

For a curve given parametrically, $x=x(t)$, $y=y(t)$, the element is
\[
ds=\sqrt{dx^2+dy^2}=\sqrt{\dot x^2+\dot y^2}\;dt,
\]
so
\[
s=\int_{t_1}^{t_2}\sqrt{\dot x^2+\dot y^2}\;dt,
\]
where $t_1,t_2$ are the \emph{parameter} values at the endpoints.

\begin{exoresolu}[The circle checked parametrically]
Verify that the circumference of a circle of radius $R$ is $2\pi R$ using the parametrisation $x=R\cos t$, $y=R\sin t$.
\tcblower
\textbf{Solution.} $\dot x=-R\sin t$, $\dot y=R\cos t$, so
\[
\dot x^2+\dot y^2=R^2\sin^2t+R^2\cos^2t=R^2.
\]
One full turn corresponds to $t$ from $0$ to $2\pi$:
\[
s=\int_0^{2\pi}\sqrt{R^2}\;dt=\int_0^{2\pi}R\,dt=2\pi R. \qquad\checkmark
\]
Note how clean this is compared with the Cartesian attempt $y=\sqrt{R^2-x^2}$, which leads to the improper integral $\displaystyle\int_{-R}^{R}\frac{R}{\sqrt{R^2-x^2}}dx$.
\end{exoresolu}

\begin{exoresolu}[Length of an astroid arch]
Find the length of the curve $x=\cos^3 t$, $y=\sin^3 t$ for $0\le t\le \dfrac{\pi}{2}$.
\tcblower
\textbf{Solution.} $\dot x=-3\cos^2t\sin t$, $\dot y=3\sin^2 t\cos t$. Then
\[
\dot x^2+\dot y^2=9\cos^4t\sin^2t+9\sin^4t\cos^2t
=9\sin^2t\cos^2t(\cos^2t+\sin^2t)=9\sin^2t\cos^2t.
\]
On $\left[0,\tfrac{\pi}{2}\right]$, $\sin t\cos t\ge0$, so $\sqrt{\dot x^2+\dot y^2}=3\sin t\cos t$:
\[
s=\int_0^{\pi/2}3\sin t\cos t\,dt
=\frac32\int_0^{\pi/2}\sin 2t\,dt
=\frac32\left[-\frac{\cos 2t}{2}\right]_0^{\pi/2}
=\frac32.
\]
\end{exoresolu}

\begin{attention}[Limits must match the variable of integration]
In a parametric problem the limits of the integral are values of $t$, not of $x$. If a question says ``from $(1,0)$ to $(0,1)$'' for $x=\cos^3t$, $y=\sin^3t$, you must first convert: $(1,0)$ corresponds to $t=0$ and $(0,1)$ to $t=\dfrac{\pi}{2}$. Mixing $x$-limits with a $t$-integral is one of the most frequent errors at the GCE.
\end{attention}

\begin{attention}[The square root covers the whole sum]
A very common slip is to write $\sqrt{\dot x^2+\dot y^2}$ as $\dot x+\dot y$, or to forget the square root altogether. Remember: $\sqrt{a^2+b^2}\ne a+b$ in general. The integrand is the \emph{single} square root of the \emph{sum} of the two squares — it comes from Pythagoras' theorem, $ds=\sqrt{dx^2+dy^2}$.
\end{attention}

\courssub{When the Integral Is Not Elementary: a Numerical Estimate (Extension)}

Not every curve yields a perfect square. For example the parabola $y=x^2$ gives $s=\displaystyle\int\sqrt{1+4x^2}\,dx$ (doable with a hyperbolic substitution), while the ellipse gives rise to the famous \emph{elliptic integrals}, which have no elementary antiderivative. In such cases we estimate the integral numerically. The simplest idea is the \cle{trapezium rule}: split $[a,b]$ into $n$ strips of width $h=\dfrac{b-a}{n}$, evaluate the integrand $g(x)=\sqrt{1+(y')^2}$ at the nodes, and use
\[
s\approx \frac{h}{2}\Bigl[g(x_0)+2g(x_1)+\dots+2g(x_{n-1})+g(x_n)\Bigr].
\]
Notice the beautiful consistency: approximating the \emph{curve} by chords and approximating the \emph{integral} by trapezia are two faces of the same idea. This is an extension remark only — at the GCE the integrals you meet are designed to be exact.

\courssub{Applications}

\textbf{Length of a sagging electric cable.} As we saw, a uniform cable hanging under its own weight between two poles at the same height follows a catenary $y=c\cosh\!\left(\dfrac{x}{c}\right)$, where $c$ depends on the tension and the weight per unit length. Then $y'=\sinh\!\left(\dfrac{x}{c}\right)$ and $1+(y')^2=\cosh^2\!\left(\dfrac{x}{c}\right)$, so for poles at $x=\pm a$:
\[
s=\int_{-a}^{a}\cosh\!\left(\tfrac{x}{c}\right)dx=2c\sinh\!\left(\tfrac{a}{c}\right).
\]
The engineer orders slightly more cable than the pole spacing $2a$; the formula tells exactly how much more.

\textbf{Length of a road bend on a map.} A cartographer models a bend of the Yaoundé--Douala road by a curve $y=f(x)$ between two junctions. The true driving distance along the bend is not the straight-line (``as the crow flies'') distance but $\displaystyle\int_a^b\sqrt{1+(f'(x))^2}\,dx$, scaled by the map ratio. The same computation gives the length of a railway curve, a pipeline, or the perimeter of a piece of land bounded by a curved river.

\begin{aretenir}[Key points of Lesson 52]
\begin{itemize}
\item The arc length element is $ds=\sqrt{dx^2+dy^2}$ (Pythagoras on an infinitesimal triangle).
\item $y=f(x)$: $s=\displaystyle\int_a^b\sqrt{1+(y')^2}\,dx$;\quad $x=g(y)$: $s=\displaystyle\int_c^d\sqrt{1+(x')^2}\,dy$.
\item Parametric: $s=\displaystyle\int_{t_1}^{t_2}\sqrt{\dot x^2+\dot y^2}\,dt$, with limits in $t$.
\item In examination questions $1+(y')^2$ almost always becomes a perfect square — hunt for it.
\item If no exact form exists, estimate the integral numerically (trapezium rule).
\end{itemize}
\end{aretenir}

\begin{exercice}[Practice]
\begin{enumerate}
\item Find the length of $y=\dfrac{x^2}{2}-\dfrac14\ln x$ from $x=1$ to $x=2$.
\item Find the length of $y=\ln(\cos x)$ from $x=0$ to $x=\dfrac{\pi}{4}$.
\item A curve is given by $x=e^t\cos t$, $y=e^t\sin t$. Show that $\sqrt{\dot x^2+\dot y^2}=\sqrt2\,e^t$ and find the length from $t=0$ to $t=\pi$.
\item The curve $y=\dfrac{2}{3}(x-1)^{3/2}$ runs from $x=1$ to $x=5$. Find its length.
\end{enumerate}
\end{exercice}

\begin{corrige}[Exercise 1]
$y'=\dfrac{x}{2}-\dfrac{1}{4x}$? No — differentiate carefully: $y'=x-\dfrac{1}{4x}$. Then
$1+(y')^2=1+x^2-\dfrac12+\dfrac{1}{16x^2}=\left(x+\dfrac{1}{4x}\right)^2$.
Hence $s=\displaystyle\int_1^2\left(x+\frac{1}{4x}\right)dx=\left[\frac{x^2}{2}+\frac14\ln x\right]_1^2=\frac32+\frac14\ln 2$.
\end{corrige}

\begin{corrige}[Exercise 2]
$y'=-\tan x$, so $1+(y')^2=1+\tan^2x=\sec^2x$ and $\sec x>0$ on $\left[0,\tfrac{\pi}{4}\right]$.
$s=\displaystyle\int_0^{\pi/4}\sec x\,dx=\bigl[\ln|\sec x+\tan x|\bigr]_0^{\pi/4}=\ln(\sqrt2+1)\approx0.881$.
\end{corrige}

\begin{corrige}[Exercise 3]
$\dot x=e^t(\cos t-\sin t)$, $\dot y=e^t(\sin t+\cos t)$. Squaring and adding:
$\dot x^2+\dot y^2=e^{2t}\bigl[(\cos t-\sin t)^2+(\sin t+\cos t)^2\bigr]=e^{2t}\cdot 2(\cos^2t+\sin^2t)=2e^{2t}$.
So $\sqrt{\dot x^2+\dot y^2}=\sqrt2\,e^t$ and
$s=\displaystyle\int_0^\pi\sqrt2\,e^t\,dt=\sqrt2\left(e^\pi-1\right)\approx 31.3$.
\end{corrige}

\begin{corrige}[Exercise 4]
$y'=(x-1)^{1/2}$, so $1+(y')^2=x$. Then
$s=\displaystyle\int_1^5\sqrt{x}\,dx=\left[\frac23 x^{3/2}\right]_1^5=\frac23\left(5\sqrt5-1\right)\approx 6.79$.
\end{corrige}

\begin{savaistu}[Why the hanging chain is not a parabola]
Galileo believed a hanging chain took the shape of a parabola. In the 17th century, Huygens showed this was wrong, and Johann Bernoulli found the true shape, the catenary $y=\cosh x$, using the brand-new calculus. The word \emph{catenary} comes from the Latin \emph{catena}, ``chain''. The same curve, turned upside down, is the ideal shape for an arch: the Gateway Arch in Saint Louis (USA) is a weighted catenary. Every time you compute $2\sinh a$ for a sagging cable, you are repeating Bernoulli's discovery.
\end{savaistu}

\coursec{Surfaces and Volumes of Revolution (Lessons 53 \& 54)}

In the previous section we learnt how to measure the \emph{length} of a curve by adding up tiny arc elements $ds = \sqrt{1+(y')^{2}}\,dx$. We now put that same philosophy — \emph{slice, approximate, sum, integrate} — to work in three dimensions. Take the region under a curve $y=f(x)$ and spin it once around an axis: the region sweeps out a solid (a \cle{solid of revolution}) and the curve itself sweeps out a surface (a \cle{surface of revolution}). A water pot on a potter's wheel in Bamenda, a football, a drinking glass, a church bell: all are solids or surfaces of revolution. Our task is to compute their \textbf{volumes} and their \textbf{surface areas} using definite integrals.

\courssub{Volume of Revolution about the $x$-axis: the Disc Method}

\textbf{Intuition.} Imagine slicing a loaf of bread into very thin slices perpendicular to its axis. Each slice is almost a thin cylinder (a \cle{disc}). If we can find the volume of one slice and add all slices together, we obtain the volume of the loaf. Integration performs exactly this addition.

\textbf{Derivation.} Let $f$ be continuous and non-negative on $[a,b]$, and rotate the region under $y=f(x)$ about the $x$-axis. A thin slice at position $x$, of thickness $\delta x$, is approximately a cylinder of radius $y=f(x)$ and height $\delta x$, hence of volume
\[ \delta V \approx \pi y^{2}\,\delta x = \pi [f(x)]^{2}\,\delta x. \]
Summing over all slices and letting $\delta x \to 0$ turns the sum into a definite integral.

\begin{definition}[Volume of revolution about the $x$-axis (disc method)]
The volume of the solid generated by rotating about the $x$-axis the region bounded by $y=f(x)$, the $x$-axis, and the lines $x=a$ and $x=b$ is
\[ \boxed{V = \pi \int_{a}^{b} y^{2}\,dx = \pi \int_{a}^{b} [f(x)]^{2}\,dx.} \]
\end{definition}

\begin{popfigure}
\begin{center}
\begin{tikzpicture}[scale=1.05]
% axis
\draw[->, thick, popink] (-0.4,0) -- (6.6,0) node[below] {$x$};
\draw[->, thick, popink] (0,-0.4) -- (0,2.6) node[left] {$y$};
% curve
\draw[very thick, popblue, domain=0.5:5.6, samples=60] plot (\x,{0.35+1.5*exp(-0.25*(\x-2.6)^2)+0.12*\x});
% region fill
\fill[popblueL, opacity=0.55] (0.5,0) -- plot[domain=0.5:5.6, samples=60] (\x,{0.35+1.5*exp(-0.25*(\x-2.6)^2)+0.12*\x}) -- (5.6,0) -- cycle;
% disc at x = 3.4, radius y = 2.036
\fill[poporange, opacity=0.85] (3.4,0) ellipse (0.16 and 2.036);
\draw[popdark, thick] (3.4,0) ellipse (0.16 and 2.036);
\draw[popdark] (3.4,0) -- (3.4,2.036) node[midway, right] {$y$};
\draw[<->, popdark] (3.24,-0.35) -- (3.56,-0.35) node[midway, below] {$dx$};
% ghost of solid
\draw[popmuted, dashed, domain=0.5:5.6, samples=60] plot (\x,{-(0.35+1.5*exp(-0.25*(\x-2.6)^2)+0.12*\x)});
\draw[popmuted, dashed] (5.6,0) ellipse (0.16 and 1.18);
\node[popink] at (1.1,2.3) {$y=f(x)$};
\node[popink, align=center] at (5.9,-1.6) {axis of\\ revolution};
\end{tikzpicture}
\end{center}
\legende{Rotating the region under $y=f(x)$ about the $x$-axis. The slice at $x$ is a disc of radius $y$ and thickness $dx$, of volume $\pi y^{2}\,dx$.}
\end{popfigure}

\begin{exemplebox}[Volume of a cone recovered by integration]
Rotate the line $y=\dfrac{r}{h}x$ between $x=0$ and $x=h$ about the $x$-axis. This generates a cone of base radius $r$ and height $h$:
\[ V=\pi\int_{0}^{h}\left(\frac{r}{h}x\right)^{2}dx=\pi\frac{r^{2}}{h^{2}}\left[\frac{x^{3}}{3}\right]_{0}^{h}=\frac{1}{3}\pi r^{2}h, \]
exactly the mensuration formula you met in lower classes. Integration \emph{proves} it.
\end{exemplebox}

\begin{exoresolu}[Volume of a sphere]
Show, by rotating a semicircle about the $x$-axis, that the volume of a sphere of radius $R$ is $V=\dfrac{4}{3}\pi R^{3}$.
\tcblower
The upper semicircle of radius $R$ centred at the origin has equation $y=\sqrt{R^{2}-x^{2}}$, so $y^{2}=R^{2}-x^{2}$ for $-R\le x\le R$. Rotating it about the $x$-axis gives the sphere:
\begin{align*}
V&=\pi\int_{-R}^{R}\left(R^{2}-x^{2}\right)dx
=\pi\left[R^{2}x-\frac{x^{3}}{3}\right]_{-R}^{R}\\
&=\pi\left[\left(R^{3}-\frac{R^{3}}{3}\right)-\left(-R^{3}+\frac{R^{3}}{3}\right)\right]
=\pi\cdot\frac{4R^{3}}{3}
=\frac{4}{3}\pi R^{3}.
\end{align*}
\end{exoresolu}

\courssub{Volume about the $y$-axis and the Washer Method}

If the region is described more naturally by $x$ as a function of $y$, say $x=g(y)$ for $c\le y\le d$, and is rotated about the $y$-axis, the same slicing argument (horizontal slices this time) gives
\[ V=\pi\int_{c}^{d} x^{2}\,dy=\pi\int_{c}^{d}[g(y)]^{2}\,dy. \]

Now suppose the region rotated about the $x$-axis lies \emph{between two curves} $y=y_{2}(x)$ (outer, farther from the axis) and $y=y_{1}(x)$ (inner), with $0\le y_{1}(x)\le y_{2}(x)$ on $[a,b]$. The solid is hollow: each slice is a \cle{washer} (a disc with a hole), of area $\pi y_{2}^{2}-\pi y_{1}^{2}$.

\begin{propriete}[Washer method]
The volume generated by rotating about the $x$-axis the region between $y=y_{2}(x)$ (outer radius) and $y=y_{1}(x)$ (inner radius) is
\[ \boxed{V=\pi\int_{a}^{b}\left(y_{2}^{\,2}-y_{1}^{\,2}\right)dx.} \]
About the $y$-axis, with outer radius $x_{2}(y)$ and inner radius $x_{1}(y)$: $V=\pi\displaystyle\int_{c}^{d}\left(x_{2}^{\,2}-x_{1}^{\,2}\right)dy$.
\end{propriete}

\begin{attention}[A very costly squaring error]
In the washer method the integrand is $\left(\text{outer radius}\right)^{2}-\left(\text{inner radius}\right)^{2}$, \textbf{not} $\left(\text{outer}-\text{inner}\right)^{2}$. Remember: $(a^{2}-b^{2})\neq(a-b)^{2}$. Examiners report this mistake every session.
\end{attention}

\begin{exoresolu}[A truncated cone (bucket)]
The region between $y=x$ and $y=\tfrac{x}{2}$ for $0\le x\le 4$ is rotated about the $x$-axis. Find the volume generated.
\tcblower
Outer radius $y_{2}=x$, inner radius $y_{1}=\tfrac{x}{2}$:
\[ V=\pi\int_{0}^{4}\left(x^{2}-\frac{x^{2}}{4}\right)dx=\pi\int_{0}^{4}\frac{3x^{2}}{4}\,dx=\pi\left[\frac{x^{3}}{4}\right]_{0}^{4}=16\pi\ \text{cubic units}. \]
\end{exoresolu}

\begin{exoresolu}[Volume of a paraboloid ``pot'']
A pot is modelled by rotating the curve $y=\dfrac{x^{2}}{10}$ (in cm) about the $y$-axis, from $y=0$ to $y=10$. Find its capacity in litres.
\tcblower
Solve for $x$: $x^{2}=10y$. Rotating about the $y$-axis:
\[ V=\pi\int_{0}^{10}x^{2}\,dy=\pi\int_{0}^{10}10y\,dy=\pi\left[5y^{2}\right]_{0}^{10}=500\pi\ \mathrm{cm^{3}}. \]
Since $1$ litre $=1000\ \mathrm{cm^{3}}$, the capacity is $\dfrac{500\pi}{1000}=\dfrac{\pi}{2}\approx 1.57$ litres.
\end{exoresolu}

\courssub{Surface Area of Revolution}

\textbf{Intuition.} How much sheet metal is needed to build a funnel? How much paint covers a dome? These are questions about the \emph{area} of a surface of revolution. Slice the surface into thin circular \textbf{bands}. A band at position $x$ is a piece of a cone (a \cle{frustum}) of radius $y$ and \emph{slant height} $ds$ — the arc element from Lesson 52. Unrolled, such a band has area approximately
\[ \delta S \approx 2\pi y\,ds \qquad\text{(circumference } \times \text{ slant height).} \]
Summing and passing to the limit gives the formula.

\begin{popfigure}
\begin{center}
\begin{tikzpicture}[scale=1.1]
% frustum band
\draw[thick, popdark, fill=poporange, opacity=0.9] (0,0.9) -- (2.6,1.5) -- (2.6,-1.5) -- (0,-0.9) -- cycle;
\draw[popdark, thick] (0,0) ellipse (0.22 and 0.9);
\draw[popdark, thick] (2.6,0) ellipse (0.3 and 1.5);
% axis
\draw[->, thick, popink] (-0.8,0) -- (4.6,0) node[below] {$x$};
% radius
\draw[popblue, thick] (2.6,0) -- (2.6,1.5) node[midway, right] {$y$};
% slant
\draw[popgreen, very thick] (0,0.9) -- (2.6,1.5) node[midway, above left] {$ds$};
\node[popink, align=center] at (3.4,-1.9) {frustum band:\\ area $\approx 2\pi y\,ds$};
\end{tikzpicture}
\end{center}
\legende{A thin band of the surface is a frustum of slant height $ds$ and radius $y$; its area is approximately $2\pi y\,ds$.}
\end{popfigure}

\begin{propriete}[Surface area of revolution]
If $y=f(x)$ has a continuous derivative on $[a,b]$, the area of the surface generated by rotating the curve about the $x$-axis is
\[ \boxed{S=2\pi\int_{a}^{b} y\,ds=2\pi\int_{a}^{b} y\sqrt{1+\left(\frac{dy}{dx}\right)^{2}}\,dx.} \]
Rotating about the $y$-axis instead, the radius becomes $x$:
\[ S=2\pi\int_{a}^{b} x\,ds=2\pi\int_{a}^{b} x\sqrt{1+\left(\frac{dy}{dx}\right)^{2}}\,dx. \]
\textbf{Parametric versions.} If the curve is given by $x=x(t)$, $y=y(t)$, $t\in[t_{1},t_{2}]$, then $ds=\sqrt{\dot{x}^{2}+\dot{y}^{2}}\,dt$ and
\[ S_{x\text{-axis}}=2\pi\int_{t_{1}}^{t_{2}} y(t)\sqrt{\dot{x}^{2}+\dot{y}^{2}}\,dt,\qquad
S_{y\text{-axis}}=2\pi\int_{t_{1}}^{t_{2}} x(t)\sqrt{\dot{x}^{2}+\dot{y}^{2}}\,dt. \]
\end{propriete}

\begin{attention}[The most common error of the chapter]
For surfaces, the integrand is $2\pi\times\text{radius}\times ds$, \textbf{never} $2\pi\times\text{radius}\times dx$. The slant element $ds=\sqrt{1+(y')^{2}}\,dx$ is essential: a steep curve sweeps far more surface than its horizontal projection suggests. Forgetting the square-root factor is the single most frequent mistake in GCE scripts on this topic.
\end{attention}

\begin{exoresolu}[Surface area of a sphere]
Prove that the surface area of a sphere of radius $R$ is $S=4\pi R^{2}$ by rotating a semicircle about the $x$-axis.
\tcblower
Take $y=\sqrt{R^{2}-x^{2}}$, $-R\le x\le R$. Then
\[ \frac{dy}{dx}=\frac{-x}{\sqrt{R^{2}-x^{2}}},\qquad
1+\left(\frac{dy}{dx}\right)^{2}=1+\frac{x^{2}}{R^{2}-x^{2}}=\frac{R^{2}}{R^{2}-x^{2}}. \]
Hence
\[ y\sqrt{1+(y')^{2}}=\sqrt{R^{2}-x^{2}}\cdot\frac{R}{\sqrt{R^{2}-x^{2}}}=R, \]
a constant! Therefore
\[ S=2\pi\int_{-R}^{R}R\,dx=2\pi R\,(2R)=4\pi R^{2}. \]
The miraculous cancellation is why Archimedes loved this result.
\end{exoresolu}

\begin{exoresolu}[Sheet metal for a paraboloid dish]
A satellite dish is formed by rotating $y=\dfrac{x^{2}}{10}$ (in cm) about the $y$-axis for $0\le x\le 10$. Find the area of sheet metal required, to the nearest $\mathrm{cm^{2}}$.
\tcblower
Rotation about the $y$-axis, radius $x$, with $y'=\dfrac{x}{5}$:
\[ S=2\pi\int_{0}^{10}x\sqrt{1+\frac{x^{2}}{25}}\,dx. \]
Substitute $u=1+\dfrac{x^{2}}{25}$, $du=\dfrac{2x}{25}dx$, so $x\,dx=\dfrac{25}{2}du$; when $x=0$, $u=1$; when $x=10$, $u=5$:
\[ S=2\pi\cdot\frac{25}{2}\int_{1}^{5}\sqrt{u}\,du=25\pi\left[\frac{2}{3}u^{3/2}\right]_{1}^{5}
=\frac{50\pi}{3}\left(5\sqrt{5}-1\right). \]
Numerically $S\approx \dfrac{50\pi}{3}(11.180-1)\approx 533.0\ \mathrm{cm^{2}}$. About $533\ \mathrm{cm^{2}}$ of sheet metal is needed.
\end{exoresolu}

\begin{popfigure}
\begin{center}
\begin{tikzpicture}
\begin{axis}[
  width=10.5cm, height=7cm,
  axis lines=middle,
  xlabel={$x$}, ylabel={$y$},
  xmin=-0.3, xmax=2.6, ymin=-0.3, ymax=2.6,
  xtick={1,2}, ytick={1,2},
  domain=0:2, samples=80,
  legend style={at={(0.98,0.95)}, anchor=north east, font=\small}
]
% semicircle y = sqrt(4-x^2)
\addplot[very thick, popblue] {sqrt(max(0,4-x^2))};
\addlegendentry{$y=\sqrt{4-x^{2}}$ (outer)}
% inner curve
\addplot[very thick, poporange] {0.5*x+0.4};
\addlegendentry{$y=0.5x+0.4$ (inner)}
% washer slice at x=1.2
\draw[popdark, thick] (axis cs:1.2,0) -- (axis cs:1.2,1.6);
\draw[popgreen, very thick] (axis cs:1.2,1.0) -- (axis cs:1.2,1.6);
\node[popgreen, font=\small, anchor=west] at (axis cs:1.25,1.35) {washer};
\draw[<->, popmuted] (axis cs:1.35,0) -- (axis cs:1.35,1.6) node[midway, right, font=\small] {$y_{2}$};
\draw[<->, popmuted] (axis cs:1.05,0) -- (axis cs:1.05,1.0) node[midway, left, font=\small] {$y_{1}$};
\end{axis}
\end{tikzpicture}
\end{center}
\legende{The semicircle $y=\sqrt{4-x^{2}}$ rotated about the $x$-axis gives a sphere; the region between the two curves generates a hollow solid whose slice is a washer of area $\pi(y_{2}^{2}-y_{1}^{2})$.}
\end{popfigure}

\courssub{Strategy, Units and an Astonishing Paradox}

\begin{methode}[Five-step strategy for volumes and surfaces of revolution]
\begin{enumerate}
\item \textbf{Sketch} the region or curve and the axis of rotation. This is a compulsory habit, never optional.
\item \textbf{Identify the axis} of rotation ($x$-axis or $y$-axis).
\item \textbf{Choose the variable} of integration: $dx$ for slicing perpendicular to the $x$-axis, $dy$ for the $y$-axis. Write $y^{2}$ (or $x^{2}$) in terms of that variable.
\item \textbf{Write the radius} of a slice (and both radii for a washer); for a surface, write also the slant element $ds=\sqrt{1+(y')^{2}}\,dx$.
\item \textbf{Integrate} between the correct limits, simplify, and give the answer with sensible units.
\end{enumerate}
\end{methode}

\begin{attention}[Units and realistic answers]
\begin{itemize}
\item Volumes come in cubic units: $\mathrm{cm^{3}}$, $\mathrm{m^{3}}$; convert with $1$ litre $=1000\ \mathrm{cm^{3}}$ when a capacity is requested.
\item Surface areas come in square units: $\mathrm{cm^{2}}$, $\mathrm{m^{2}}$ — think of sheet metal or paint.
\item Always check plausibility: a drinking pot of $1.6$ litres is reasonable; one of $1600$ litres is not. A quick estimate ($V\approx$ base area $\times$ height) catches most blunders.
\end{itemize}
\end{attention}

\begin{savaistu}[Gabriel's horn: finite volume, infinite area (enrichment, off-syllabus)]
Rotate the curve $y=\dfrac{1}{x}$, $x\ge 1$, about the $x$-axis. The resulting infinitely long ``horn'' has volume
\[ V=\pi\int_{1}^{\infty}\frac{1}{x^{2}}\,dx=\pi\left[-\frac{1}{x}\right]_{1}^{\infty}=\pi, \]
which is \emph{finite}. But its surface area satisfies
\[ S=2\pi\int_{1}^{\infty}\frac{1}{x}\sqrt{1+\frac{1}{x^{4}}}\,dx>2\pi\int_{1}^{\infty}\frac{dx}{x}=\infty, \]
which is \emph{infinite}! You could fill the horn with $\pi$ cubic units of paint, yet no amount of paint could cover its inside. This paradox, discussed since the seventeenth century, shows how subtle infinite processes can be — and why the rigorous machinery of limits and improper integrals matters.
\end{savaistu}

\begin{exercice}[Practice]
\begin{enumerate}
\item The region under $y=\sqrt{x}$ from $x=0$ to $x=4$ is rotated about the $x$-axis. Find the volume of the solid.
\item The region enclosed by $y=x^{2}$ and $y=2x$ is rotated about the $x$-axis. Using the washer method, find the volume generated.
\item The curve $y=\dfrac{x^{3}}{3}+\dfrac{1}{4x}$ from $x=1$ to $x=2$ is rotated about the $x$-axis. Show that the surface area is $S=2\pi\displaystyle\int_{1}^{2}\left(\frac{x^{3}}{3}+\frac{1}{4x}\right)\left(x^{2}+\frac{1}{4x^{2}}\right)dx$ and evaluate it. \emph{Hint: check that }$1+(y')^{2}=\left(x^{2}+\frac{1}{4x^{2}}\right)^{2}$.
\item A bowl is formed by rotating $y=x^{2}$ (in cm) about the $y$-axis for $0\le y\le 9$. Find its volume in $\mathrm{cm^{3}}$ and its capacity in litres.
\end{enumerate}
\end{exercice}

\begin{corrige}[Exercise 1]
\begin{enumerate}
\item $V=\pi\displaystyle\int_{0}^{4}x\,dx=\pi\left[\frac{x^{2}}{2}\right]_{0}^{4}=8\pi$ cubic units.
\item Curves meet where $x^{2}=2x$, i.e. $x=0,2$. On $[0,2]$, $y_{2}=2x$ (outer), $y_{1}=x^{2}$ (inner):
\[ V=\pi\int_{0}^{2}\left(4x^{2}-x^{4}\right)dx=\pi\left[\frac{4x^{3}}{3}-\frac{x^{5}}{5}\right]_{0}^{2}=\pi\left(\frac{32}{3}-\frac{32}{5}\right)=\frac{64\pi}{15}. \]
\item $y'=x^{2}-\dfrac{1}{4x^{2}}$, so $1+(y')^{2}=1+x^{4}-\dfrac{1}{2}+\dfrac{1}{16x^{4}}=\left(x^{2}+\dfrac{1}{4x^{2}}\right)^{2}$, giving the stated form. Multiplying out:
\[ \left(\frac{x^{3}}{3}+\frac{1}{4x}\right)\left(x^{2}+\frac{1}{4x^{2}}\right)=\frac{x^{5}}{3}+\frac{x}{12}+\frac{x}{4}+\frac{1}{16x^{3}}=\frac{x^{5}}{3}+\frac{x}{3}+\frac{1}{16x^{3}}. \]
Hence
\[ S=2\pi\left[\frac{x^{6}}{18}+\frac{x^{2}}{6}-\frac{1}{32x^{2}}\right]_{1}^{2}
=2\pi\left[\left(\frac{64}{18}+\frac{4}{6}-\frac{1}{128}\right)-\left(\frac{1}{18}+\frac{1}{6}-\frac{1}{32}\right)\right]
=\frac{515\pi}{64}. \]
\item $x^{2}=y$, so $V=\pi\displaystyle\int_{0}^{9}y\,dy=\pi\left[\frac{y^{2}}{2}\right]_{0}^{9}=\frac{81\pi}{2}\approx 127.2\ \mathrm{cm^{3}}$, i.e. about $0.127$ litre.
\end{enumerate}
\end{corrige}

\begin{aretenir}[Key points of Lessons 53 \& 54]
\begin{itemize}
\item Disc method about the $x$-axis: $V=\pi\displaystyle\int_{a}^{b}y^{2}\,dx$; about the $y$-axis: $V=\pi\displaystyle\int_{c}^{d}x^{2}\,dy$.
\item Washer method: $V=\pi\displaystyle\int_{a}^{b}\left(y_{2}^{2}-y_{1}^{2}\right)dx$ — squares \emph{first}, subtract \emph{after}.
\item Surface of revolution: $S=2\pi\displaystyle\int(\text{radius})\,ds$, with $ds=\sqrt{1+(y')^{2}}\,dx$ (Cartesian) or $ds=\sqrt{\dot{x}^{2}+\dot{y}^{2}}\,dt$ (parametric). Never replace $ds$ by $dx$.
\item Integration recovers and \emph{proves} the classical formulae: cone $\tfrac13\pi r^{2}h$, sphere $\tfrac43\pi R^{3}$ and $4\pi R^{2}$.
\item Always sketch first, choose the variable to match the axis, and finish with realistic units.
\end{itemize}
\end{aretenir}

\coursec{Centroids of Plane Figures and Centres of Gravity of Solids of Revolution (Lessons 55 \& 56)}

In the previous section we learnt how to compute the \emph{volume} and the \emph{surface area} of the solid obtained by spinning a region about an axis. But knowing how much material a solid contains is not the whole story. Imagine a water tank shaped like a cone, manufactured in a workshop in Douala: at which point should it be supported so that it balances perfectly horizontally? At which point will a spinning top, a funnel or a metal bowl balance on a finger? The answer is the \cle{centre of gravity}, and definite integrals give us the exact tool to locate it. For flat objects (thin metal plates, called \emph{laminas}), the same question leads to the \cle{centroid}. This section develops both ideas rigorously.

\courssub{Moments and the Centroid of a Plane Region}

\textbf{Intuition: the see-saw principle.}\par
Two children on a see-saw balance when the products ``weight $\times$ distance from the pivot'' are equal on both sides. This product is called a \cle{moment}. A plane lamina balances on a knife edge along a line when the total moment of its area about that line is the same on both sides. The \cle{centroid} is the point where the whole area may be considered concentrated for balance purposes.

\begin{definition}[Centroid and centre of gravity]
The \cle{centroid} (or geometric centre) of a plane lamina of uniform thickness and density is the point $(\bar{x},\bar{y})$ at which the lamina would balance if supported there. For a solid of uniform density, the corresponding balance point is called the \cle{centre of gravity}. Since density is uniform, it cancels out of every formula, so we work purely with areas and volumes.
\end{definition}

\textbf{Setting up the moment integrals.}\par
Consider the region under the curve $y=f(x)$, with $f(x)\geq 0$, between $x=a$ and $x=b$. Slice it into thin vertical strips of width $\delta x$ at position $x$; each strip is approximately a rectangle of height $y$ and area $y\,\delta x$, whose own centre is at $\left(x,\ \tfrac{y}{2}\right)$.

\begin{itemize}
\item \textbf{Moment about the $y$-axis.} The strip's area is $y\,\delta x$ and its centre is at distance $x$ from the $y$-axis, so its moment is $x\,y\,\delta x$. Summing and passing to the limit:
\[ M_y=\int_a^b x\,y\,dx. \]
\item \textbf{Moment about the $x$-axis.} The strip's centre is at height $\tfrac{y}{2}$, so its moment is $\tfrac{y}{2}\cdot y\,\delta x=\tfrac12 y^2\,\delta x$. Hence:
\[ M_x=\frac12\int_a^b y^2\,dx. \]
\end{itemize}

\begin{popfigure}
\begin{center}
\begin{tikzpicture}[scale=1.1]
\draw[->,thick] (-0.3,0) -- (5.3,0) node[below]{$x$};
\draw[->,thick] (0,-0.3) -- (0,3.2) node[left]{$y$};
\draw[domain=0.4:4.6,smooth,thick,popblue] plot (\x,{0.35*(\x-0.4)*(4.6-\x)*0.55+0.6});
\fill[popblueL,opacity=0.7] (0.4,0) -- plot[domain=0.4:4.6] (\x,{0.35*(\x-0.4)*(4.6-\x)*0.55+0.6}) -- (4.6,0) -- cycle;
\fill[poporange,opacity=0.75] (2.3,0) rectangle (2.65,1.66);
\draw[popdark] (2.3,0) rectangle (2.65,1.66);
\fill[popdark] (2.475,0.83) circle (0.045);
\node[popdark,right] at (2.65,0.83) {\small centre $\left(x,\frac{y}{2}\right)$};
\draw[<->,popmuted] (2.3,-0.35) -- (2.65,-0.35) node[midway,below]{\small $\delta x$};
\draw[<->,popmuted] (0.15,0) -- (0.15,1.66) node[midway,left]{\small $y$};
\draw[dashed,popmuted] (2.475,0) -- (2.475,1.66);
\node[below] at (2.475,-0.05) {\small $x$};
\node[popblue] at (4.75,1.2) {\small $y=f(x)$};
\node[align=center] at (1.1,2.6) {\small strip of area $y\,\delta x$};
\end{tikzpicture}
\end{center}
\legende{A vertical strip of the region under $y=f(x)$: its moment about the $y$-axis is $xy\,\delta x$ and about the $x$-axis is $\tfrac12 y^2\,\delta x$.}
\end{popfigure}

\begin{propriete}[Centroid of the region under a curve]
For the region under $y=f(x)$ from $x=a$ to $x=b$, with area $A=\displaystyle\int_a^b y\,dx$:
\[ \bar{x}=\frac{M_y}{A}=\frac{\displaystyle\int_a^b x\,y\,dx}{\displaystyle\int_a^b y\,dx},
\qquad
\bar{y}=\frac{M_x}{A}=\frac{\displaystyle\frac12\int_a^b y^2\,dx}{\displaystyle\int_a^b y\,dx}. \]
The derivation (strip argument above) is examinable and frequently asked.
\end{propriete}

\begin{attention}[$\tfrac12 y^2$ versus $\pi y^2$]
For the centroid coordinate $\bar{y}$ of a \emph{plane region}, the integrand is $\tfrac12 y^2$ (the factor $\tfrac12$ is the height of the strip's centre). For the \emph{volume} of revolution it is $\pi y^2$. Mixing these two is one of the most common errors at the GCE: $\bar{y}$ has \textbf{no} $\pi$ in it.
\end{attention}

\textbf{Region between two curves.}\par
If the region lies between $y_2=f(x)$ (top) and $y_1=g(x)$ (bottom), each strip has height $y_2-y_1$ and its centre is at height $\tfrac{y_1+y_2}{2}$. The same strip argument gives
\[ A=\int_a^b (y_2-y_1)\,dx,\quad
M_y=\int_a^b x(y_2-y_1)\,dx,\quad
M_x=\frac12\int_a^b \left(y_2^{\,2}-y_1^{\,2}\right)dx, \]
and again $(\bar{x},\bar{y})=(M_y/A,\,M_x/A)$.

\begin{methode}[Exploit symmetry first]
Before computing anything, look for an axis of symmetry of the region:
\begin{enumerate}
\item If the region is symmetric about the $y$-axis (or any vertical line $x=c$), then $\bar{x}$ lies on it: $\bar{x}=0$ (or $\bar{x}=c$) with \textbf{no calculation}.
\item If the region is symmetric about the $x$-axis, then $\bar{y}=0$.
\item Compute only the remaining coordinate by integration, and state clearly: ``by symmetry, \ldots''. This earns a mark and saves time.
\end{enumerate}
\end{methode}

\begin{exoresolu}[Centroid of the region under a parabola]
Find the centroid of the region bounded by $y=4-x^2$ and the $x$-axis.
\tcblower
\textbf{Solution.} The curve meets the $x$-axis at $x=\pm 2$. The region is symmetric about the $y$-axis, so $\bar{x}=0$.
\[ A=\int_{-2}^{2}(4-x^2)\,dx=\Big[4x-\tfrac{x^3}{3}\Big]_{-2}^{2}=2\left(8-\tfrac{8}{3}\right)=\frac{32}{3}. \]
\[ M_x=\frac12\int_{-2}^{2}(4-x^2)^2\,dx=\frac12\int_{-2}^{2}\left(16-8x^2+x^4\right)dx
=\Big[8x-\tfrac{4x^3}{3}+\tfrac{x^5}{10}\Big]_{-2}^{2}=\frac{256}{15}. \]
Hence $\displaystyle \bar{y}=\frac{256/15}{32/3}=\frac{256}{15}\times\frac{3}{32}=\frac{8}{5}$.
The centroid is $\left(0,\ \tfrac85\right)$ — slightly below the mid-height $2$, as expected since the region is wider at the bottom.
\end{exoresolu}

\courssub{Standard Results to Derive and Memorise}

\begin{exemplebox}[Centroid of a triangle]
Take a triangle of base $b$ on the $x$-axis and height $h$, with apex above the origin: the sloping side is $y=h\left(1-\tfrac{x}{b}\right)$. Then
\[ A=\frac12 bh,\qquad
M_x=\frac12\int_0^b h^2\left(1-\tfrac{x}{b}\right)^2 dx
=\frac{h^2}{2}\cdot\frac{b}{3}=\frac{bh^2}{6}, \]
so $\bar{y}=\dfrac{bh^2/6}{bh/2}=\dfrac{h}{3}$. \textbf{The centroid of a triangle is at $\tfrac{h}{3}$ from the base} (and $\tfrac{2h}{3}$ from the apex).
\end{exemplebox}

\begin{exemplebox}[Centroid of a semicircular lamina]
For the upper semicircle of radius $r$, $y=\sqrt{r^2-x^2}$, $x\in[-r,r]$. By symmetry $\bar{x}=0$. Using $A=\tfrac12\pi r^2$:
\[ M_x=\frac12\int_{-r}^{r}(r^2-x^2)\,dx=\frac12\Big[r^2x-\tfrac{x^3}{3}\Big]_{-r}^{r}=\frac{2r^3}{3},
\qquad
\bar{y}=\frac{2r^3/3}{\pi r^2/2}=\frac{4r}{3\pi}. \]
\textbf{The centroid of a semicircular lamina is at $\dfrac{4r}{3\pi}\approx 0.424\,r$ from the diameter.} For a \emph{quadrant} of a circle, the same computation on $[0,r]$ gives $\bar{x}=\bar{y}=\dfrac{4r}{3\pi}$.
\end{exemplebox}

\begin{popfigure}
\begin{center}
\begin{tikzpicture}[scale=1.5]
\fill[popblueL] (2,0) arc (0:180:2) -- cycle;
\draw[thick,popblue] (2,0) arc (0:180:2);
\draw[thick,popblue] (-2,0) -- (2,0);
\draw[->,thick] (-2.5,0) -- (2.7,0) node[below]{$x$};
\draw[->,thick] (0,-0.2) -- (0,2.5) node[left]{$y$};
\fill[poporange] (0,0.849) circle (0.06);
\node[poporange,right,align=left] at (0.12,0.849) {\small centroid $\left(0,\frac{4r}{3\pi}\right)$};
\draw[<->,popmuted] (0,-0.25) -- (2,-0.25) node[midway,below]{\small $r$};
\draw[dashed,popmuted] (0,0) -- (0,0.849);
\node[left] at (0,0.42) {\small $\frac{4r}{3\pi}$};
\end{tikzpicture}
\end{center}
\legende{Semicircular lamina of radius $r$: the centroid lies on the axis of symmetry at height $\frac{4r}{3\pi}$ above the diameter.}
\end{popfigure}

\begin{propriete}[Table of standard centroids (proofs examinable)]
\begin{center}
\begin{tabularx}{\linewidth}{|l|X|l|}
\hline
\rowcolor{popblueL} \textbf{Body} & \textbf{Position of centroid / centre of gravity} & \textbf{Value} \\
\hline
Triangle & from the base, on the median & $h/3$ \\
\hline
Semicircular lamina & from the diameter, on the axis of symmetry & $4r/3\pi$ \\
\hline
Quadrant of a circle & from each straight edge & $4r/3\pi$ \\
\hline
Solid cone & from the base, on the axis & $h/4$ \\
\hline
Solid hemisphere & from the flat face, on the axis & $3r/8$ \\
\hline
\end{tabularx}
\end{center}
These results may be quoted in examinations, but you must be able to \emph{derive} each one by integration when asked.
\end{propriete}

\courssub{Composite Figures: Adding and Subtracting Parts}

Many GCE questions involve a lamina built from rectangles, triangles and semicircles, possibly with a hole cut out. The key principle is that \textbf{moments add}: if a lamina is split into parts of areas $A_i$ with centroids $(\bar{x}_i,\bar{y}_i)$, then
\[ \bar{x}=\frac{\sum A_i\bar{x}_i}{\sum A_i},\qquad \bar{y}=\frac{\sum A_i\bar{y}_i}{\sum A_i}. \]
A hole is treated as a part with \emph{negative} area.

\begin{methode}[Composite bodies: use a table]
\begin{enumerate}
\item Choose axes (usually along edges of the figure) and split the lamina into standard parts.
\item Draw a table with columns: \emph{part}, \emph{area} $A_i$, \emph{centroid} $(\bar{x}_i,\bar{y}_i)$, \emph{moments} $A_i\bar{x}_i$, $A_i\bar{y}_i$. Give holes a negative area.
\item Sum the columns and divide: $\bar{x}=\dfrac{\sum A_i\bar{x}_i}{\sum A_i}$, $\bar{y}=\dfrac{\sum A_i\bar{y}_i}{\sum A_i}$.
\end{enumerate}
\textbf{Exam tip:} presenting the table earns method marks even if an arithmetic slip occurs later — examiners can follow every step.
\end{methode}

\begin{exoresolu}[Rectangle with a semicircular top]
A lamina consists of a rectangle of width $4$ and height $6$, surmounted by a semicircle of radius $2$ on its top edge. Find the height of its centroid above the base.
\tcblower
\textbf{Solution.} Take the $x$-axis along the base. By symmetry $\bar{x}$ is on the central vertical line; we need only $\bar{y}$.
\begin{center}
\begin{tabular}{|l|c|c|c|}
\hline
\rowcolor{popblueL} Part & $A_i$ & $\bar{y}_i$ & $A_i\bar{y}_i$ \\
\hline
Rectangle & $24$ & $3$ & $72$ \\
\hline
Semicircle & $2\pi$ & $6+\dfrac{8}{3\pi}$ & $12\pi+16$ \\
\hline
Total & $24+2\pi$ & -- & $88+12\pi$ \\
\hline
\end{tabular}
\end{center}
(For the semicircle, $\bar{y}_i=6+\dfrac{4r}{3\pi}=6+\dfrac{8}{3\pi}$, so $A_i\bar{y}_i=2\pi\left(6+\dfrac{8}{3\pi}\right)=12\pi+\dfrac{16}{3}\cdot 3=12\pi+16$.)
\[ \bar{y}=\frac{88+12\pi}{24+2\pi}=\frac{44+6\pi}{12+\pi}\approx \frac{62.85}{15.14}\approx 4.15. \]
The centroid is about $4.15$ units above the base — pulled above the rectangle's own centre ($3$) by the semicircular cap, as expected.
\end{exoresolu}

\courssub{Centre of Gravity of a Solid of Revolution}

Spin the region under $y=f(x)$ about the $x$-axis. The solid obtained is symmetric about the $x$-axis, so its centre of gravity lies \emph{on} the axis: $\bar{y}=0$ (and $\bar{z}=0$) by symmetry. To find $\bar{x}$, slice the solid into thin discs perpendicular to the axis: the disc at position $x$ has radius $y$, thickness $\delta x$, volume $\pi y^2\,\delta x$, and its own centre of gravity is on the axis at $x$. Its moment about the $y$-$z$ plane is $x\cdot\pi y^2\,\delta x$. Summing and dividing by the total volume:

\begin{propriete}[Centre of gravity of a solid of revolution]
For the solid formed by rotating $y=f(x)$, $a\leq x\leq b$, about the $x$-axis:
\[ \bar{x}=\frac{\displaystyle\int_a^b x\,y^2\,dx}{\displaystyle\int_a^b y^2\,dx},\qquad \bar{y}=0. \]
(The factor $\pi$ cancels between numerator and denominator.) The disc argument is examinable.
\end{propriete}

\begin{popfigure}
\begin{center}
\begin{tikzpicture}[scale=1.0]
\draw[->,thick] (-0.3,0) -- (7.6,0) node[below]{$x$};
\draw[thick,popblue] (0.5,0) -- (6.5,2.1);
\draw[thick,popblue] (0.5,0) -- (6.5,-2.1);
\draw[popblue,thick] (6.5,0) ellipse (0.28 and 2.1);
\fill[poporange,opacity=0.6] (3.4,0) ellipse (0.22 and 1.09);
\draw[popdark] (3.4,0) ellipse (0.22 and 1.09);
\fill[popgreenL,opacity=0.6] (3.75,0) ellipse (0.22 and 1.21);
\draw[popdark] (3.75,0) ellipse (0.22 and 1.21);
\draw[popdark] (3.4,1.09) -- (3.75,1.21);
\draw[popdark] (3.4,-1.09) -- (3.75,-1.21);
\node[align=center] at (3.6,-1.9) {\small disc: volume $\pi y^2\delta x$};
\fill[poporange] (4.9,0) circle (0.07);
\node[poporange,above] at (4.9,0.12) {\small $\bar{x}$};
\draw[dashed,popmuted] (3.4,0) -- (3.4,1.09);
\node[left,popmuted] at (3.4,0.6) {\small $y$};
\node[below,popmuted] at (3.57,0) {\small $x$};
\end{tikzpicture}
\end{center}
\legende{A solid of revolution sliced into discs of radius $y$ and thickness $\delta x$; the centre of gravity lies on the axis at $\bar{x}$.}
\end{popfigure}

\begin{exemplebox}[Centre of gravity of a solid cone]
A cone of height $h$ and base radius $r$ is generated by rotating the line $y=\dfrac{r}{h}x$ about the $x$-axis for $x\in[0,h]$ (apex at the origin, base at $x=h$):
\[ \bar{x}=\frac{\displaystyle\int_0^h x\left(\tfrac{r}{h}x\right)^2dx}{\displaystyle\int_0^h \left(\tfrac{r}{h}x\right)^2dx}
=\frac{\displaystyle\int_0^h x^3\,dx}{\displaystyle\int_0^h x^2\,dx}
=\frac{h^4/4}{h^3/3}=\frac{3h}{4}. \]
Measured from the apex, $\bar{x}=\tfrac{3h}{4}$; equivalently \textbf{the centre of gravity of a solid cone is at $\dfrac{h}{4}$ from its base}.
\end{exemplebox}

\begin{exemplebox}[Centre of gravity of a solid hemisphere]
Rotate the quarter circle $y=\sqrt{r^2-x^2}$, $x\in[0,r]$, about the $x$-axis (flat face at $x=0$):
\[ \bar{x}=\frac{\displaystyle\int_0^r x(r^2-x^2)\,dx}{\displaystyle\int_0^r (r^2-x^2)\,dx}
=\frac{\Big[\tfrac{r^2x^2}{2}-\tfrac{x^4}{4}\Big]_0^r}{\Big[r^2x-\tfrac{x^3}{3}\Big]_0^r}
=\frac{\tfrac{r^4}{2}-\tfrac{r^4}{4}}{r^3-\tfrac{r^3}{3}}
=\frac{r^4/4}{2r^3/3}=\frac{3r}{8}. \]
\textbf{The centre of gravity of a solid hemisphere is at $\dfrac{3r}{8}$ from the flat face.}
\end{exemplebox}

\begin{attention}[Common mistakes]
\begin{itemize}
\item Forgetting the factor $\tfrac12$ in $M_x=\tfrac12\int y^2\,dx$ for a plane region.
\item Using $\int x y\,dx$ instead of $\int x y^2\,dx$ in the numerator for a solid of revolution: for a solid, the disc volume involves $y^2$.
\item Quoting $\bar{y}=0$ for a solid spun about the $x$-axis without mentioning \emph{symmetry} — the justification carries a mark.
\item For composite bodies, forgetting that a hole counts with \emph{negative} area or volume.
\end{itemize}
\end{attention}

\courssub{Applied Problems: Balancing and Stability}

\begin{experience}[Balancing point of a manufactured funnel]
A workshop in Bamenda spins a thin-walled funnel: a conical solid of height $12\ \mathrm{cm}$ (centre of gravity at $\tfrac{h}{4}=3\ \mathrm{cm}$ from its base) is joined to a cylindrical stem of length $8\ \mathrm{cm}$ and the same uniform density, whose centre of gravity is at its midpoint. Taking the cone's base as origin along the common axis, with the stem extending from $x=0$ to $x=-8$:
\begin{center}
\begin{tabular}{|l|c|c|c|}
\hline
\rowcolor{popblueL} Part & Volume & $\bar{x}_i$ & $V_i\bar{x}_i$ \\
\hline
Cone (radius $6$) & $\tfrac13\pi(36)(12)=144\pi$ & $3$ & $432\pi$ \\
\hline
Cylinder (radius $1$) & $8\pi$ & $-4$ & $-32\pi$ \\
\hline
Total & $152\pi$ & -- & $400\pi$ \\
\hline
\end{tabular}
\end{center}
\[ \bar{x}=\frac{400\pi}{152\pi}=\frac{50}{19}\approx 2.63\ \mathrm{cm}. \]
The funnel balances on a finger placed about $2.6\ \mathrm{cm}$ from the cone's base. The same table method answers stability questions for tanks, bowls and machined parts: the object topples when the vertical through its centre of gravity falls outside the base of support, so locating $(\bar{x},\bar{y})$ precisely is a genuine engineering task.
\end{experience}

\begin{savaistu}[Why engineers care]
The centroid governs the bending of beams (the ``neutral axis'' passes through the centroid of the cross-section), the stability of dams and retaining walls, and the balance of aircraft and vehicles. The results you have just derived — $\tfrac{h}{3}$, $\tfrac{4r}{3\pi}$, $\tfrac{h}{4}$, $\tfrac{3r}{8}$ — appear in every engineering handbook, and they are all one-line integrals.
\end{savaistu}

\begin{exercice}[Practice]
\begin{enumerate}
\item Find the centroid of the region bounded by $y=\sqrt{x}$, the $x$-axis and the line $x=4$.
\item A lamina is a square of side $2r$ with a semicircle of radius $r$ removed from one edge (diameter on that edge). Find the distance of the centroid from the opposite edge.
\item Show by integration that the centre of gravity of a solid hemisphere of radius $r$ is at $\tfrac{3r}{8}$ from its flat face, and deduce the centre of gravity of a solid hemispherical bowl combined with a flat circular lid (treat the lid as a uniform lamina and the bowl as a solid hemisphere of the same material).
\end{enumerate}
\end{exercice}

\begin{corrige}[Exercise 1]
$A=\displaystyle\int_0^4 \sqrt{x}\,dx=\Big[\tfrac23 x^{3/2}\Big]_0^4=\tfrac{16}{3}$.
$M_y=\displaystyle\int_0^4 x^{3/2}\,dx=\Big[\tfrac25 x^{5/2}\Big]_0^4=\tfrac{64}{5}$, so $\bar{x}=\dfrac{64/5}{16/3}=\dfrac{12}{5}$.
$M_x=\dfrac12\displaystyle\int_0^4 x\,dx=4$, so $\bar{y}=\dfrac{4}{16/3}=\dfrac34$. Centroid: $\left(\tfrac{12}{5},\tfrac34\right)$.
\end{corrige}

\begin{corrige}[Exercise 2]
Place the square with the cut edge on the line $x=0$ and the opposite edge at $x=2r$. Square: $A_1=4r^2$, $\bar{x}_1=r$. Semicircular hole (diameter on $x=0$, bulging into the square): $A_2=-\tfrac12\pi r^2$, $\bar{x}_2=\tfrac{4r}{3\pi}$.
\[ \bar{x}=\frac{4r^2\cdot r-\tfrac12\pi r^2\cdot\tfrac{4r}{3\pi}}{4r^2-\tfrac12\pi r^2}
=\frac{4r^3-\tfrac23 r^3}{r^2\left(4-\tfrac{\pi}{2}\right)}
=\frac{\tfrac{10}{3}r}{4-\tfrac{\pi}{2}}=\frac{20r}{3(8-\pi)}. \]
Distance from the opposite edge: $2r-\bar{x}=2r-\dfrac{20r}{3(8-\pi)}=\dfrac{(28-6\pi)r}{3(8-\pi)}\approx 0.65\,r$.
\end{corrige}

\begin{corrige}[Exercise 3]
With $y=\sqrt{r^2-x^2}$ on $[0,r]$: numerator $\displaystyle\int_0^r x(r^2-x^2)\,dx=\tfrac{r^4}{4}$; denominator $\displaystyle\int_0^r (r^2-x^2)\,dx=\tfrac{2r^3}{3}$; hence $\bar{x}=\dfrac{r^4/4}{2r^3/3}=\dfrac{3r}{8}$. 

For the solid hemispherical bowl plus lid: the lid is a uniform circular lamina of mass proportional to its area $\pi r^2$ (placed at $x=0$); the bowl is a uniform solid hemisphere of mass proportional to its volume $\tfrac23\pi r^3$ with centre of gravity at $x=\tfrac{3r}{8}$. Taking moments about the lid:
\[ \bar{x}=\frac{\tfrac23\pi r^3\cdot\tfrac{3r}{8}+\pi r^2\cdot 0}{\tfrac23\pi r^3+\pi r^2}
=\frac{\tfrac14\pi r^4}{\pi r^2\left(\tfrac{2r}{3}+1\right)}=\frac{r^2}{4\left(\tfrac{2r}{3}+1\right)}=\frac{3r^2}{8r+12}, \]
measured from the lid along the axis.
\end{corrige}

\begin{aretenir}[Key points of Lessons 55 \& 56]
\begin{itemize}
\item Plane region under $y=f(x)$: $\bar{x}=\dfrac{\int x y\,dx}{\int y\,dx}$, $\bar{y}=\dfrac{\tfrac12\int y^2\,dx}{\int y\,dx}$.
\item Between two curves: replace $y$ by $y_2-y_1$ and $y^2$ by $y_2^2-y_1^2$.
\item Symmetry gives one coordinate for free — always check first.
\item Standard results: triangle $\tfrac{h}{3}$ from base; semicircle $\tfrac{4r}{3\pi}$ from diameter; solid cone $\tfrac{h}{4}$ from base; solid hemisphere $\tfrac{3r}{8}$ from flat face.
\item Solid of revolution about the $x$-axis: $\bar{x}=\dfrac{\int x y^2\,dx}{\int y^2\,dx}$, $\bar{y}=0$ by symmetry.
\item Composite bodies: tabulate area/volume, centroid, moments; holes count negatively.
\end{itemize}
\end{aretenir}

\coursec{Theorems of Pappus and Synthesis (Lesson 57)}

In the previous sections we computed centroids of plane regions and centres of gravity of solids of revolution by direct integration. That work was honest but sometimes long. In this final lesson we discover a beautiful shortcut, due to the Greek geometer \cle{Pappus of Alexandria} (4th century AD), which links the two families of quantities of this chapter: \textbf{areas and volumes on one side, arc lengths and surfaces on the other, with centroids in the middle}. Pappus's theorems say, roughly: \emph{when a region spins around an external axis, the volume it generates equals its area multiplied by the distance travelled by its centroid}. We shall state both theorems, apply them to the torus, the cone and the sphere, use them \emph{backwards} to recover centroids without any integration, and finish with a synthesis of the whole chapter and a GCE strategy.

\courssub{Pappus's First Theorem: Volumes}

\textbf{Intuition first.} Take a plane lamina of area $A$ and rotate it once around an external axis. Every point of the lamina travels along a circle, but different points travel different distances. The centroid $G$ is, by definition, the ``average position'' of the area, so it is natural to guess that the volume swept out equals the area times the distance travelled by $G$, namely $2\pi d$ where $d$ is the distance from $G$ to the axis. Pappus's first theorem says this guess is exactly right.

\begin{propriete}[Pappus's First Theorem (volume)]
Let $\mathcal{R}$ be a plane region of area $A$, and let $\ell$ be an axis in its plane \textbf{which does not cross} $\mathcal{R}$. If $\mathcal{R}$ is rotated through one complete turn about $\ell$, the volume of the solid generated is
\[ V = A \times 2\pi d, \]
where $d$ is the perpendicular distance from the \cle{centroid} $G$ of $\mathcal{R}$ to the axis $\ell$. (The derivation from the washer formula is instructive and sketched below; quoting the theorem is examinable.)
\end{propriete}

\textbf{Why it is true.} Rotate about the $x$-axis a region bounded by $y = f(x)$ (above) and $y = g(x)$ (below), $a \le x \le b$, lying entirely above the axis (so $g(x) \ge 0$). The washer method gives
\[ V = \pi\int_a^b \left[f(x)^2 - g(x)^2\right] dx. \]
But the centroid ordinate of the region satisfies (Lesson 55)
\[ \bar{y} = \frac{1}{A}\int_a^b \tfrac{1}{2}\left[f(x)^2 - g(x)^2\right] dx
\quad\Longrightarrow\quad
\int_a^b \left[f^2 - g^2\right] dx = 2A\bar{y}. \]
Substituting, $V = \pi \cdot 2A\bar{y} = A \cdot 2\pi \bar{y}$, and $\bar{y}$ is precisely the distance from $G$ to the $x$-axis. $\blacksquare$

\begin{popfigure}
\begin{tikzpicture}[scale=1.0]
% axis
\draw[thick, popdark, ->] (0,-0.4) -- (0,4.6) node[above] {axis $\ell$};
% region
\fill[popblueL] (3.2,1.2) .. controls (4.4,0.8) and (5.0,1.6) .. (4.6,2.6)
   .. controls (4.2,3.4) and (3.0,3.2) .. (2.8,2.2)
   .. controls (2.7,1.7) and (2.8,1.4) .. (3.2,1.2) -- cycle;
\draw[popblue, thick] (3.2,1.2) .. controls (4.4,0.8) and (5.0,1.6) .. (4.6,2.6)
   .. controls (4.2,3.4) and (3.0,3.2) .. (2.8,2.2)
   .. controls (2.7,1.7) and (2.8,1.4) .. (3.2,1.2) -- cycle;
\node[popblue] at (4.6,3.1) {$\mathcal{R}$, area $A$};
% centroid
\fill[poporange] (3.7,2.1) circle (0.07) node[right, popdark] {\ $G$};
% distance d
\draw[<->, poppink, thick] (0,2.1) -- (3.7,2.1) node[midway, above] {$d$};
\draw[dashed, gray] (3.7,2.1) -- (3.7,0.2);
% circular path of G (ellipse suggesting a circle in perspective)
\draw[popgreen, thick, dashed] (0,2.1) ellipse (3.7 and 0.85);
\node[popgreen] at (-2.6,3.15) {path of $G$: $2\pi d$};
% arrow of rotation
\draw[->, popgreen] (-3.4,2.6) arc (150:60:0.9);
\end{tikzpicture}
\legende{Pappus I: a region of area $A$ whose centroid $G$ is at distance $d$ from an external axis sweeps a volume $V = 2\pi d\,A$.}
\end{popfigure}

\courssub{Pappus's Second Theorem: Surfaces}

Exactly the same idea applied to a \emph{curve} instead of a region gives the surface area of revolution.

\begin{propriete}[Pappus's Second Theorem (surface)]
Let $\mathcal{C}$ be a plane curve of length $L$, and let $\ell$ be an axis in its plane \textbf{which does not cross} $\mathcal{C}$. Rotating $\mathcal{C}$ once about $\ell$ generates a surface of area
\[ S = L \times 2\pi d, \]
where $d$ is the distance from the \cle{centroid of the curve} (its ``centre of length'') to $\ell$. (Quoting the theorem is examinable; it follows from $S = 2\pi\int y\, ds$ exactly as above.)
\end{propriete}

\begin{attention}[The axis must not cross the region or the curve]
Both theorems require the axis of rotation \textbf{not to intersect} the generating region (Pappus I) or the generating curve (Pappus II). If the axis cuts the region, the solid self-overlaps and the formula counts some volume twice. Always check this condition before applying Pappus.
\end{attention}

\begin{attention}[$d$ is measured from the centroid to the axis]
The quantity $d$ is the \textbf{perpendicular distance from $G$ to the axis of rotation}, not automatically $\bar{x}$ or $\bar{y}$. For rotation about the vertical line $x = a$, the distance is $d = |\bar{x} - a|$. For rotation about $y = -2$, it is $d = |\bar{y} + 2|$. Misreading the axis is the most frequent error in Pappus questions.
\end{attention}

\courssub{Direct Applications}

\begin{exemplebox}[The torus (anchor ring)]
Rotate a disc of radius $r$ whose centre is at distance $R$ ($R > r$) from the axis.
\begin{itemize}
\item \textbf{Volume (Pappus I):} the area is $A = \pi r^2$, the centroid is the centre of the disc, so $d = R$:
\[ V = \pi r^2 \times 2\pi R = 2\pi^2 R r^2. \]
\item \textbf{Surface (Pappus II):} the generating curve is the circle, of length $L = 2\pi r$, whose centroid is again its centre:
\[ S = 2\pi r \times 2\pi R = 4\pi^2 R r. \]
\end{itemize}
Two integrations avoided — the whole computation fits on two lines.
\end{exemplebox}

\begin{popfigure}
\begin{tikzpicture}[scale=0.9]
% axis
\draw[thick, popdark, ->] (-4.2,-2.6) -- (-4.2,2.8) node[above] {axis};
% generating circle
\draw[popblue, thick, fill=popblueL] (-1.6,0) circle (0.9);
\fill[poporange] (-1.6,0) circle (0.06);
\node[popdark] at (-1.6,-0.35) {\scriptsize centre};
\draw[<->, poppink] (-4.2,0.55) -- (-1.6,0.55) node[midway, above] {$R$};
\draw[<->, popblue] (-1.6,0) -- (-0.85,0.53) node[midway, above right] {\scriptsize $r$};
% torus body (right side): outer and inner ellipses
\draw[popdark, thick] (2.2,0) ellipse (2.5 and 1.05);
\draw[popdark, thick] (2.2,0.13) ellipse (0.7 and 0.28);
\draw[popdark, thick, dashed] (2.2,-0.13) ellipse (0.7 and 0.28);
% path of centroid
\draw[popgreen, dashed] (2.2,0) ellipse (1.6 and 0.66);
\fill[poporange] (3.8,0) circle (0.06) node[right, popdark] {\scriptsize $G$};
% labels
\node[popblue, align=center] at (2.2,-1.7) {$V = 2\pi^2 R r^2$,\quad $S = 4\pi^2 R r$};
\draw[->, popmuted] (-0.4,0) -- (0.6,0);
\end{tikzpicture}
\legende{A circle of radius $r$ whose centre travels on a circle of radius $R$ generates a torus with $V = 2\pi^2 R r^2$ and $S = 4\pi^2 R r$.}
\end{popfigure}

\begin{exemplebox}[Cone and sphere in seconds]
\textbf{Cone.} Rotate a right triangle with legs $h$ (on the axis) and $r$ about the leg $h$. Its area is $A = \tfrac{1}{2}rh$ and its centroid is at distance $\bar{x} = \tfrac{r}{3}$ from the axis (centroid of a right triangle: one third of the way from the leg on the axis). Hence
\[ V = \tfrac{1}{2}rh \times 2\pi \cdot \tfrac{r}{3} = \tfrac{1}{3}\pi r^2 h. \]

\textbf{Sphere.} Rotate a \emph{semicircular disc} of radius $r$ about its diameter. Area $A = \tfrac{1}{2}\pi r^2$, centroid at distance $\dfrac{4r}{3\pi}$ from the diameter:
\[ V = \tfrac{1}{2}\pi r^2 \times 2\pi \cdot \frac{4r}{3\pi} = \frac{4}{3}\pi r^3. \]
\end{exemplebox}

\courssub{Using Pappus Backwards: Finding Centroids}

This is the examiners' favourite trick: if the volume (or surface) of revolution is \emph{already known} by geometry, Pappus's theorem becomes an equation for the unknown centroid.

\begin{methode}[Finding a centroid without integration]
\begin{enumerate}
\item Identify the region (or curve) and the solid (or surface) it generates by rotation.
\item Write the known volume (or surface area) of that solid from memory.
\item Apply Pappus: $V = 2\pi d\, A$ (or $S = 2\pi d\, L$) with $d$ unknown.
\item Solve for $d$: this is the distance of the centroid from the axis.
\end{enumerate}
\end{methode}

\begin{exoresolu}[Centroid of a semicircular disc — the classic GCE trick]
A semicircular disc of radius $r$, rotated about its diameter, generates a sphere of volume $V = \dfrac{4}{3}\pi r^3$. Deduce the distance of the centroid of the semicircular disc from its diameter.
\tcblower
\textbf{Solution.} The area of the semicircular disc is $A = \dfrac{1}{2}\pi r^2$. By Pappus's first theorem, with $d$ the distance from $G$ to the diameter (the axis does not cross the interior of the region, so the theorem applies):
\[ V = A \times 2\pi d
\quad\Longrightarrow\quad
\frac{4}{3}\pi r^3 = \frac{1}{2}\pi r^2 \times 2\pi d = \pi^2 r^2 d. \]
Solving,
\[ d = \frac{\tfrac{4}{3}\pi r^3}{\pi^2 r^2} = \frac{4r}{3\pi}. \]
Hence the centroid of a semicircular disc lies on its axis of symmetry, at distance $\boxed{\dfrac{4r}{3\pi}}$ from the diameter — a result obtained with \emph{no integration at all}. The same trick applied to a semicircular \emph{arc} (length $L = \pi r$, generating the sphere's surface $S = 4\pi r^2$) gives the centroid of the arc: $4\pi r^2 = \pi r \cdot 2\pi d \Rightarrow d = \dfrac{2r}{\pi}$.
\end{exoresolu}

\begin{exoresolu}[Rotation about a shifted axis]
The region bounded by $y = x^2$, the $x$-axis and $x = 2$ has area $A = \dfrac{8}{3}$ and centroid $\bar{x} = \dfrac{3}{2}$, $\bar{y} = \dfrac{6}{5}$ (computed in Lesson 55). Find the volume generated when this region is rotated once about the line $x = 3$.
\tcblower
\textbf{Solution.} The axis $x = 3$ is vertical and does not cross the region (which lies in $0 \le x \le 2$), so Pappus I applies. The distance from the centroid to the axis is
\[ d = |\bar{x} - 3| = \left|\tfrac{3}{2} - 3\right| = \tfrac{3}{2}. \]
Therefore
\[ V = A \times 2\pi d = \frac{8}{3} \times 2\pi \times \frac{3}{2} = 8\pi. \]
Note carefully: $d$ is $|\bar{x} - 3|$, \textbf{not} $\bar{x}$. Using $\bar{x}$ would give the wrong answer — always measure from the actual axis.
\end{exoresolu}

\begin{attention}[Common mistakes with Pappus]
\begin{itemize}
\item Using $\bar{x}$ instead of $|\bar{x} - a|$ when the axis is $x = a$.
\item Applying the theorem when the axis cuts through the region (e.g. rotating a full disc about a diameter and expecting Pappus to give the sphere — it fails because the axis crosses the region; use a \emph{semicircular} disc instead).
\item Confusing the centroid of a \emph{region} (for volumes) with the centroid of a \emph{curve} (for surfaces): for a semicircle these are $\dfrac{4r}{3\pi}$ and $\dfrac{2r}{\pi}$ respectively.
\end{itemize}
\end{attention}

\courssub{Synthesis: The Whole Chapter in One Problem}

\begin{exoresolu}[Complete study of a circular arc]
A quarter-circle of radius $r$ in the first quadrant is given by $y = \sqrt{r^2 - x^2}$, $0 \le x \le r$. 
(a) Find its arc length. 
(b) Rotating it about the $x$-axis, find the surface generated. 
(c) Deduce by Pappus II the centroid of the arc. 
(d) The region under the arc is rotated about the $x$-axis: find the volume and deduce the centroid of the quarter-disc.
\tcblower
\textbf{Solution.}
\textbf{(a)} The quarter-circle has length $L = \dfrac{2\pi r}{4} = \dfrac{\pi r}{2}$.

\textbf{(b)} The surface is a hemisphere: $S = 2\pi r^2$.

\textbf{(c)} By Pappus II, $S = L \times 2\pi d$, so
\[ d = \frac{S}{2\pi L} = \frac{2\pi r^2}{2\pi \cdot \tfrac{\pi r}{2}} = \frac{2r}{\pi}. \]
By symmetry $\bar{x} = \bar{y} = \dfrac{2r}{\pi}$.

\textbf{(d)} The volume is a hemisphere: $V = \dfrac{2}{3}\pi r^3$. The quarter-disc has area $A = \dfrac{\pi r^2}{4}$. By Pappus I,
\[ d = \frac{V}{2\pi A} = \frac{\tfrac{2}{3}\pi r^3}{2\pi \cdot \tfrac{\pi r^2}{4}} = \frac{4r}{3\pi}, \]
so the centroid of the quarter-disc is $\left(\dfrac{4r}{3\pi}, \dfrac{4r}{3\pi}\right)$. Every formula of the chapter — arc length, surface, volume, centroid, Pappus — appears in this one exercise.
\end{exoresolu}

\courssub{Chapter Revision Map and GCE Strategy}

\begin{center}
\begin{tabularx}{\linewidth}{|l|X|X|}
\hline
\rowcolor{popblueL} \textbf{Question wording} & \textbf{Tool} & \textbf{Key formula} \\
\hline
``mean value'', ``average value'' & Lesson 51 & $\bar{f} = \frac{1}{b-a}\int_a^b f\,dx$ \\
\hline
``RMS'', ``effective value'' & Lesson 51 & $\sqrt{\frac{1}{b-a}\int_a^b f^2\,dx}$ \\
\hline
``length of the curve/arc'' & Lesson 52 & $L = \int_a^b \sqrt{1 + (y')^2}\,dx$ \\
\hline
``surface generated by rotating'' & Lesson 53 & $S = 2\pi\int y\, ds$ \\
\hline
``volume of the solid generated'' & Lesson 54 & $V = \pi\int y^2\,dx$ \\
\hline
``centroid / centre of gravity'' & Lessons 55--56 & $\bar{x} = \frac{1}{A}\int x\, dA$ \\
\hline
volume \emph{and} centroid both involved & Lesson 57 & $V = 2\pi d\, A$,\quad $S = 2\pi d\, L$ \\
\hline
\end{tabularx}
\end{center}

\begin{aretenir}[Exam tip: think Pappus first]
When a GCE question gives you a volume (or surface) and asks for a centroid — or gives a centroid and asks for a volume — Pappus's theorem is almost always the intended short route. Before launching into a long integration, ask: \emph{``Is this Pappus in disguise?''} Check that the axis does not cross the region, measure $d$ from the \emph{actual} axis, and the answer often comes in three lines.
\end{aretenir}

\begin{savaistu}[Pappus of Alexandria]
Pappus lived in Alexandria around the 4th century AD, and his \emph{Collection} is one of the last great works of Greek geometry. His theorems on solids of revolution were rediscovered and republished by the Swiss mathematician Paul Guldin in the 17th century — which is why they are sometimes called the \emph{Pappus--Guldin theorems}. Engineers still use them daily: the volume of a toroidal water tank, the material needed for a doughnut-shaped pipe, or the mass of a turned metal piece on a lathe are all one-line Pappus computations.
\end{savaistu}

\begin{exercice}[Pappus practice]
\begin{enumerate}
\item An equilateral triangle of side $a$ is rotated about an axis parallel to one side, at distance $a$ from that side, with the axis lying on the opposite side of the triangle from the side (so the axis is outside the triangle). Find the volume generated.
\item Using the known volume of a cone, $V = \tfrac{1}{3}\pi r^2 h$, recover by Pappus the position of the centroid of a right triangle of legs $r$ and $h$ when rotated about the leg of length $h$.
\item A semicircular arc of radius $6$ cm is rotated about a line parallel to its diameter and $10$ cm from it, on the side opposite to the arc (so the axis is outside the arc). Find the surface area generated.
\end{enumerate}
\end{exercice}

\begin{corrige}[Exercise 1]
\textbf{1.} Area $A = \tfrac{\sqrt{3}}{4}a^2$; height $h = \tfrac{\sqrt{3}}{2}a$. The centroid of an equilateral triangle is at distance $\tfrac{h}{3} = \tfrac{\sqrt{3}}{6}a$ from each side. The axis is parallel to one side at distance $a$ from it, on the opposite side of the triangle. Hence the distance from the centroid to the axis is $d = a + \tfrac{\sqrt{3}}{6}a$. Then $V = 2\pi d A = 2\pi\left(a + \tfrac{\sqrt{3}a}{6}\right)\tfrac{\sqrt{3}a^2}{4} = \tfrac{\pi a^3}{2}\left(\sqrt{3} + \tfrac{1}{2}\right)$.

\textbf{2.} Rotating the right triangle (legs $r$, $h$) about the leg $h$: the area is $A = \tfrac{1}{2}rh$. By Pappus I, $\tfrac{1}{3}\pi r^2 h = \tfrac{1}{2}rh \times 2\pi d$, giving $d = \tfrac{r}{3}$: the centroid is one third of the way from the axis leg.

\textbf{3.} Arc length $L = \pi r = 6\pi$ cm. Centroid of a semicircular arc is at distance $\tfrac{2r}{\pi} = \tfrac{12}{\pi}$ cm from the diameter. The axis is on the opposite side of the diameter from the arc, at distance $10$ cm from the diameter. Hence $d = 10 + \tfrac{12}{\pi}$ cm. Surface area $S = 2\pi d L = 2\pi\left(10 + \tfrac{12}{\pi}\right) 6\pi = 120\pi^2 + 144\pi \approx 1638\ \mathrm{cm^2}$.
\end{corrige}

\newpage

\coursec{Integration activity}

\coursec{Integrating Activity: The Nkwen Cooler}

\courssub{Aims of the activity}

This integrating activity closes the chapter \emph{Applications of Definite Integrals} by mobilising, in one single authentic modelling task, all the notions studied:

\begin{itemize}
\item compute a \cle{volume of a solid of revolution} by the disc method and convert units ($1\ \mathrm{dm}^3 = 1$ litre);
\item compute the \cle{area of a surface of revolution} using the arc-length element $ds = \sqrt{1+(y')^2}\,dx$;
\item determine the \cle{centre of gravity of a solid of revolution} and use it to discuss stability;
\item determine the \cle{centroid of a plane region} and \emph{verify} the volume with \cle{Pappus's first theorem};
\item extend the work to the \cle{mean value} and \cle{root mean square (RMS) value} of a function;
\item develop teamwork, scientific communication (poster) and the habit of \emph{justifying every integral set-up}, as required at the GCE Advanced Level.
\end{itemize}

\begin{propriete}[Key formulae for rotation about the $x$-axis]
For a curve $y=f(x)\ge 0$, $a\le x\le b$, rotated about the $x$-axis:
\[ V=\pi\int_a^b y^2\,dx,\qquad S=2\pi\int_a^b y\sqrt{1+(y')^2}\,dx,\qquad \bar{x}=\dfrac{\int_a^b x\,y^2\,dx}{\int_a^b y^2\,dx}. \]
For the generating plane region under the curve: $A=\int_a^b y\,dx$, centroid ordinate $\bar{y}=\dfrac{1}{2A}\int_a^b y^2\,dx$, and Pappus's first theorem: $V=2\pi\,\bar{y}\,A$.
Mean value of a function: $\bar{f}=\dfrac{1}{b-a}\int_a^b f(x)\,dx$; RMS value: $f_{\mathrm{RMS}}=\sqrt{\dfrac{1}{b-a}\int_a^b f(x)^2\,dx}$.
\end{propriete}

\courssub{Situation}

\textbf{Designing a clay water pot for the Bamenda market.}\par
Mama Gladys runs a pottery workshop near Nkwen, in Bamenda (North-West Region). Her clay water pots keep water cool by evaporation through the porous wall, and they are highly appreciated at the Bamenda main market. She wants to launch a new model, the ``Nkwen cooler'', whose inside profile is obtained by rotating the curve
\[ y=\sqrt{x}+1 \quad (\text{in decimetres}), \qquad 0\le x\le 4, \]
about the $x$-axis. The bottom of the pot is the disc at $x=0$ (radius $1$ dm) and the opening is the circle at $x=4$ (radius $3$ dm).

\begin{popfigure}
\begin{center}
\begin{tikzpicture}[scale=0.85]
\draw[->,popmuted] (-0.3,0) -- (4.7,0) node[below left]{$x$ (dm)};
\draw[->,popmuted] (0,-3.4) -- (0,3.6) node[below left]{$y$ (dm)};
\draw[very thick,popblue,domain=0:4,samples=80] plot (\x,{sqrt(max(0,\x))+1});
\draw[very thick,popblue,domain=0:4,samples=80] plot (\x,{-sqrt(max(0,\x))-1});
\draw[poporange,thick] (0,1) -- (0,-1);
\draw[poporange,thick] (4,3) -- (4,-3);
\draw[popline,dashed] (4,3) -- (4,0) node[midway,right,popdark]{$3$};
\draw[popline,dashed] (0,1) -- (0,0) node[midway,left,popdark]{$1$};
\node[popblue] at (2.2,3.1) {$y=\sqrt{x}+1$};
\node[popdark] at (2,-0.35) {axis of rotation};
\end{tikzpicture}
\end{center}
\legende{Profile of the pot and the solid generated by rotation about the $x$-axis.}
\end{popfigure}

Before firing the first batch, Mama Gladys asks your Further Mathematics team four practical questions:

\begin{enumerate}
\item How many \emph{litres} of water does one pot hold? (She wants to advertise an honest capacity.)
\item What is the \emph{outer surface area} of the pot? Glaze for the rim zone is sold per square decimetre, so she needs an estimate of the area of the surface of revolution.
\item Where is the \emph{centre of gravity} of the (empty) solid pot? A pot whose centre of gravity lies close to the bottom is more stable on the head or on a stand --- an important selling point.
\item Can the volume be \emph{checked by an independent method} (Pappus's theorem), so that the team is confident in the result?
\end{enumerate}

\textbf{Extension (for the strongest teams).} When the pot is half full, estimate the water depth along the axis and the mean radius of the water region. Besides, the temperature inside the workshop varies as $T(t)=25+5\sin\!\left(\dfrac{\pi t}{12}\right)$ (in \textdegree C, $t$ in hours, $0\le t\le 24$). Compute the RMS temperature over one day, which governs the drying of the clay.

\courssub{Tasks}

Work in teams of three or four. Every integral must be \emph{set up and justified} before being evaluated.

\begin{methode}[General approach for each task]
\begin{enumerate}
\item Identify the relevant formula (volume, surface area, centroid, etc.) and write it down.
\item Substitute the given function $y=\sqrt{x}+1$ and its derivative where needed.
\item Simplify the integrand algebraically.
\item Choose an appropriate integration technique (expansion, substitution, given primitives, etc.).
\item Evaluate the definite integral exactly, then give a decimal approximation with correct units.
\item Interpret the result in the context of the problem.
\end{enumerate}
\end{methode}

\begin{enumerate}
\item \textbf{Capacity.} Justify that the volume of the pot is $V=\pi\displaystyle\int_0^4(\sqrt{x}+1)^2\,dx$. Expand the integrand, compute $V$ exactly (as a multiple of $\pi$), then give the capacity in litres, rounded to one decimal place.
\item \textbf{Glaze.} Compute $y'$ and show that the arc element is $ds=\sqrt{1+\dfrac{1}{4x}}\,dx=\dfrac{\sqrt{4x+1}}{2\sqrt{x}}\,dx$. Using the substitution $u=\sqrt{x}$, show that
\[ S=2\pi\int_0^2(1+u)\sqrt{4u^2+1}\,du. \]
Evaluate $S$ (you may use $\int\sqrt{4u^2+1}\,du=\dfrac{u}{2}\sqrt{4u^2+1}+\dfrac{1}{4}\operatorname{arsinh}(2u)$ and $\int u\sqrt{4u^2+1}\,du=\dfrac{1}{12}(4u^2+1)^{3/2}$). Give $S$ to the nearest $\mathrm{dm}^2$.
\item \textbf{Stability.} Compute $\displaystyle\int_0^4 x(\sqrt{x}+1)^2\,dx$, deduce the abscissa $\bar{x}$ of the centre of gravity of the solid pot, and explain why the pot is stable (where does $G$ lie relative to the middle of the pot?).
\item \textbf{Pappus check.} For the plane region under $y=\sqrt{x}+1$ on $[0,4]$, compute the area $A$ and the ordinate $\bar{y}$ of its centroid. Apply Pappus's first theorem $V=2\pi\bar{y}A$ and compare with question 1.
\item \textbf{Extension.} (a) Find $h$ such that the pot is half full, i.e. $\displaystyle\int_0^h(\sqrt{x}+1)^2dx=\dfrac{34}{3}$, by trial to two decimal places; then compute the mean value of $y$ on $[0,h]$. (b) Show that the RMS temperature satisfies $T_{\mathrm{RMS}}^2=25^2+\dfrac{5^2}{2}$ and evaluate it.
\item \textbf{Poster.} Present all results, with units, exact values and decimal approximations, plus one sentence of advice to Mama Gladys.
\end{enumerate}

\begin{attention}[Common traps]
\begin{itemize}
\item Do not confuse $V=\pi\int y^2dx$ with $\pi\int y\,dx$: the disc method \emph{squares} the radius.
\item $1\ \mathrm{dm}^3=1$ L exactly; $1\ \mathrm{dm}^2=100\ \mathrm{cm}^2$.
\item In the surface formula the integrand is $y\,ds$, not $y\,dx$: forgetting $\sqrt{1+(y')^2}$ is the classic error.
\item The centroid of the \emph{region} (Pappus) and the centre of gravity of the \emph{solid} are different points: here $\bar{x}_{\text{solid}}\approx 2.42$ dm but $\bar{x}_{\text{region}}\approx 2.23$ dm.
\item When computing a mean value, the denominator is the length of the interval ($h-0$), not the area or volume.
\end{itemize}
\end{attention}

\courssub{Model answer}

\begin{exoresolu}[1. Capacity]
\noindent\textbf{Statement.} Compute the volume of the solid generated by rotating $y=\sqrt{x}+1$ about the $x$-axis for $0\le x\le 4$, and express the capacity in litres.

\noindent\textbf{Solution.}
A cross-section perpendicular to the $x$-axis is a disc of radius $y$, hence of area $\pi y^2$. Summing these slices from $x=0$ to $x=4$ gives
\[ V = \pi\int_0^4 y^2\,dx = \pi\int_0^4 (\sqrt{x}+1)^2\,dx. \]
Expand the integrand:
\[ (\sqrt{x}+1)^2 = x + 2\sqrt{x} + 1 = x + 2x^{1/2} + 1. \]
Integrate term by term:
\begin{align*}
V &= \pi\int_0^4 \left(x + 2x^{1/2} + 1\right)dx \\
  &= \pi\left[ \frac{x^2}{2} + \frac{4}{3}x^{3/2} + x \right]_0^4 \\
  &= \pi\left( \frac{16}{2} + \frac{4}{3}\cdot 8 + 4 \right) \\
  &= \pi\left( 8 + \frac{32}{3} + 4 \right) = \pi\left( 12 + \frac{32}{3} \right) = \frac{68\pi}{3}\ \mathrm{dm}^3.
\end{align*}
Since $1\ \mathrm{dm}^3 = 1$ litre, the capacity is
\[ V \approx \frac{68\times 3.1416}{3} \approx 71.2\ \text{litres}. \]
\end{exoresolu}

\begin{exoresolu}[2. Outer surface area]
\noindent\textbf{Statement.} Compute the area of the surface of revolution (the outer surface of the pot).

\noindent\textbf{Solution.}
First, $y = \sqrt{x}+1$, so $y' = \dfrac{1}{2\sqrt{x}}$. Then
\[ 1+(y')^2 = 1 + \frac{1}{4x} = \frac{4x+1}{4x}, \qquad
ds = \sqrt{1+(y')^2}\,dx = \frac{\sqrt{4x+1}}{2\sqrt{x}}\,dx. \]
The surface area formula is $S = 2\pi\int_0^4 y\,ds$. Substituting:
\[ S = 2\pi\int_0^4 (\sqrt{x}+1)\,\frac{\sqrt{4x+1}}{2\sqrt{x}}\,dx. \]
Use the substitution $u = \sqrt{x}$, so $x = u^2$, $dx = 2u\,du$. Limits: $x=0 \Rightarrow u=0$; $x=4 \Rightarrow u=2$. Then
\[ \sqrt{x}+1 = u+1,\quad \sqrt{4x+1} = \sqrt{4u^2+1},\quad \frac{dx}{2\sqrt{x}} = \frac{2u\,du}{2u} = du. \]
Hence
\[ S = 2\pi\int_0^2 (1+u)\sqrt{4u^2+1}\,du. \]
Split the integral:
\[ \int_0^2 (1+u)\sqrt{4u^2+1}\,du = \int_0^2 \sqrt{4u^2+1}\,du + \int_0^2 u\sqrt{4u^2+1}\,du. \]
Using the given primitives:
\begin{align*}
\int \sqrt{4u^2+1}\,du &= \frac{u}{2}\sqrt{4u^2+1} + \frac{1}{4}\operatorname{arsinh}(2u), \\
\int u\sqrt{4u^2+1}\,du &= \frac{1}{12}(4u^2+1)^{3/2}.
\end{align*}
Evaluate from $0$ to $2$:
\begin{align*}
\left[\frac{u}{2}\sqrt{4u^2+1} + \frac{1}{4}\operatorname{arsinh}(2u)\right]_0^2
&= \frac{2}{2}\sqrt{17} + \frac{1}{4}\operatorname{arsinh}(4) - 0 \\
&= \sqrt{17} + \frac{1}{4}\ln(4+\sqrt{17}), \\[4pt]
\left[\frac{1}{12}(4u^2+1)^{3/2}\right]_0^2
&= \frac{1}{12}(17^{3/2} - 1) = \frac{17\sqrt{17}}{12} - \frac{1}{12}.
\end{align*}
Summing:
\[ \int_0^2 (1+u)\sqrt{4u^2+1}\,du = \sqrt{17} + \frac{17\sqrt{17}}{12} + \frac{1}{4}\ln(4+\sqrt{17}) - \frac{1}{12}
= \frac{29\sqrt{17}}{12} + \frac{1}{4}\ln(4+\sqrt{17}) - \frac{1}{12}. \]
Numerically, $\sqrt{17}\approx 4.1231$, $\ln(4+\sqrt{17})\approx 2.0947$, giving
\[ \int_0^2 \cdots \approx 9.964 + 0.5237 - 0.0833 = 10.4044. \]
Thus
\[ S = 2\pi \times 10.4044 \approx 65.37\ \mathrm{dm}^2 \approx \cle{65.4\ \mathrm{dm}^2} \quad (\text{about } 6540\ \mathrm{cm}^2). \]
\end{exoresolu}

\begin{exoresolu}[3. Centre of gravity of the solid]
\noindent\textbf{Statement.} Find the abscissa $\bar{x}$ of the centre of gravity of the solid pot and comment on stability.

\noindent\textbf{Solution.}
By symmetry, the centre of gravity $G$ lies on the $x$-axis. Its abscissa is
\[ \bar{x} = \frac{\int_0^4 x\,y^2\,dx}{\int_0^4 y^2\,dx}. \]
We already have $\int_0^4 y^2\,dx = \frac{68}{3}$. Compute the numerator:
\begin{align*}
\int_0^4 x(\sqrt{x}+1)^2\,dx &= \int_0^4 x(x + 2\sqrt{x} + 1)\,dx \\
&= \int_0^4 \left(x^2 + 2x^{3/2} + x\right)dx \\
&= \left[ \frac{x^3}{3} + \frac{4}{5}x^{5/2} + \frac{x^2}{2} \right]_0^4 \\
&= \frac{64}{3} + \frac{4}{5}\cdot 32 + 8 = \frac{64}{3} + \frac{128}{5} + 8.
\end{align*}
Put over common denominator $15$:
\[ \frac{320}{15} + \frac{384}{15} + \frac{120}{15} = \frac{824}{15}. \]
Therefore
\[ \bar{x} = \frac{824/15}{68/3} = \frac{824}{15}\times\frac{3}{68} = \frac{824}{340} = \frac{206}{85} \approx 2.4235\ \mathrm{dm}. \]
The pot is $4$ dm tall, so $G$ is at about $60.6\%$ of the height from the bottom. This is slightly above the geometric midpoint ($2$ dm), meaning the wide belly near the opening pulls the centre of gravity upward. The pot is reasonably stable but could be improved by thickening the base (adding mass lower down). Note that $G$ lies inside the cavity, which is normal for a hollow object.
\end{exoresolu}

\begin{exoresolu}[4. Verification by Pappus's theorem]
\noindent\textbf{Statement.} Compute the area $A$ and centroid ordinate $\bar{y}$ of the generating plane region, then apply Pappus's first theorem to verify the volume.

\noindent\textbf{Solution.}
The generating region is bounded by $y=\sqrt{x}+1$, the $x$-axis, $x=0$ and $x=4$. Its area is
\[ A = \int_0^4 (\sqrt{x}+1)\,dx = \left[ \frac{2}{3}x^{3/2} + x \right]_0^4 = \frac{16}{3} + 4 = \frac{28}{3}\ \mathrm{dm}^2. \]
The centroid ordinate of this region is
\[ \bar{y} = \frac{1}{2A}\int_0^4 y^2\,dx = \frac{1}{2\cdot\frac{28}{3}}\times\frac{68}{3} = \frac{3}{56}\times\frac{68}{3} = \frac{68}{56} = \frac{17}{14}\ \mathrm{dm} \approx 1.214\ \mathrm{dm}. \]
(For information, the centroid abscissa is $\bar{x}_{\text{region}} = \frac{1}{A}\int_0^4 x y\,dx = \frac{78}{35}\approx 2.23$ dm, which differs from $\bar{x}_{\text{solid}}$.)
Pappus's first theorem states $V = 2\pi\bar{y}A$. Substituting:
\[ V = 2\pi \times \frac{17}{14} \times \frac{28}{3} = 2\pi \times \frac{17\times 2}{3} = \frac{68\pi}{3}\ \mathrm{dm}^3, \]
which matches exactly the result from question 1. The team can advertise the capacity with confidence.
\end{exoresolu}

\begin{exoresolu}[5. Extension]
\noindent\textbf{(a) Half-full depth and mean radius.}
Half the volume corresponds to $\frac{34}{3}$ under the integral sign. Define
\[ F(h) = \int_0^h (\sqrt{x}+1)^2\,dx = \frac{h^2}{2} + \frac{4}{3}h^{3/2} + h. \]
We need $F(h) = \frac{34}{3} \approx 11.3333$. Trial values:
\begin{align*}
F(2.50) &= 3.125 + 5.270 + 2.50 = 10.895, \\
F(2.55) &= 3.251 + 5.412 + 2.55 = 11.213, \\
F(2.56) &= 3.277 + 5.461 + 2.56 = 11.298, \\
F(2.57) &= 3.302 + 5.511 + 2.57 = 11.383.
\end{align*}
Since $11.333$ lies between $F(2.56)$ and $F(2.57)$, linear interpolation gives $h \approx 2.565$ dm. To two decimal places, $h \approx 2.57$ dm (or $2.56$ dm depending on rounding convention). The water reaches about $64\%$ of the height, because the pot widens upward.

The mean radius of the water region (mean value of $y$ on $[0,h]$) is
\[ \bar{y}_{[0,h]} = \frac{1}{h}\int_0^h (\sqrt{x}+1)\,dx = \frac{1}{h}\left( \frac{2}{3}h^{3/2} + h \right) = \frac{2}{3}\sqrt{h} + 1. \]
For $h = 2.56$, $\sqrt{h}=1.6$, so $\bar{y} = \frac{2}{3}\times 1.6 + 1 = 2.067$ dm. For $h = 2.57$, $\sqrt{h}\approx 1.603$, $\bar{y}\approx 2.069$ dm. So the mean radius is about $\cle{2.07\ \mathrm{dm}}$.

\noindent\textbf{(b) RMS temperature.}
The temperature varies as $T(t)=25+5\sin(\pi t/12)$ over $0\le t\le 24$ (one full period). The RMS value is
\[ T_{\mathrm{RMS}}^2 = \frac{1}{24}\int_0^{24} \left(25+5\sin\frac{\pi t}{12}\right)^2 dt. \]
Expand:
\[ \left(25+5\sin\frac{\pi t}{12}\right)^2 = 625 + 250\sin\frac{\pi t}{12} + 25\sin^2\frac{\pi t}{12}. \]
Over a full period, $\int_0^{24}\sin(\pi t/12)\,dt = 0$ and $\int_0^{24}\sin^2(\pi t/12)\,dt = \frac{1}{2}\times 24 = 12$. Hence
\[ T_{\mathrm{RMS}}^2 = \frac{1}{24}\left( 625\times 24 + 0 + 25\times 12 \right) = 625 + \frac{25}{2} = 637.5. \]
\[ T_{\mathrm{RMS}} = \sqrt{637.5} \approx 25.25\ \text{\textdegree C}. \]
Note that $T_{\mathrm{RMS}} > 25\ \text{\textdegree C}$ (the mean temperature): the RMS value always penalises fluctuations.
\end{exoresolu}

\courssub{Poster advice}

\begin{center}
\begin{tabularx}{\linewidth}{|l|X|}\hline
\rowcolor{popblueL}
\textbf{Quantity} & \textbf{Result} \\ \hline
Capacity & $\dfrac{68\pi}{3}\ \mathrm{dm}^3 \approx 71.2\ \mathrm{L}$ \\ \hline
Outer surface area & $\approx 65.4\ \mathrm{dm}^2$ ($6540\ \mathrm{cm}^2$) \\ \hline
Centre of gravity (from bottom) & $\dfrac{206}{85}\ \mathrm{dm} \approx 2.42\ \mathrm{dm}$ \\ \hline
Pappus verification & Exact match: $V=2\pi\bar{y}A = \dfrac{68\pi}{3}$ \\ \hline
Half-full depth & $h \approx 2.57\ \mathrm{dm}$ \\ \hline
Mean water radius (half-full) & $\approx 2.07\ \mathrm{dm}$ \\ \hline
RMS workshop temperature & $\approx 25.3\ \text{\textdegree C}$ \\ \hline
\end{tabularx}
\end{center}

\textbf{Advice to Mama Gladys:} ``Your Nkwen cooler holds about 71 L and needs $\approx 65\ \mathrm{dm}^2$ of glaze. The centre of gravity is at $2.42$ dm from the base; for better stability on the head, consider thickening the bottom wall or adding a wider foot. All figures are cross-checked by Pappus's theorem.''

\begin{savaistu}[Did you know?]
Pappus of Alexandria (4th century) stated his centroid theorems more than thirteen centuries before the invention of integral calculus. Today the same theorems let engineers compute volumes of turned pottery, of water tanks and of the clay pots sold in Bamenda, with a single multiplication $V=2\pi\bar{y}A$.
\end{savaistu}

\newpage

\coursec{Methods and worked examples}

\coursec{Applications of Definite Integrals}

\courssub{Mean Value and RMS Value of a Function}

\begin{definition}[Mean value of a function]
Let $f$ be a continuous function on the closed interval $[a,b]$. The \cle{mean value} (or average value) of $f$ on $[a,b]$ is defined as
\[
\mu = \frac{1}{b-a}\int_a^b f(x)\,dx.
\]
Geometrically, $\mu$ is the height of the rectangle with base $[a,b]$ that has the same area as the region under the curve $y=f(x)$.
\end{definition}

\begin{definition}[Root mean square (RMS) value]
The \cle{root mean square (RMS) value} of a continuous function $f$ on $[a,b]$ is
\[
f_{\mathrm{rms}} = \sqrt{\frac{1}{b-a}\int_a^b \bigl(f(x)\bigr)^2\,dx}.
\]
In physics and engineering, the RMS value of a periodic current or voltage gives the equivalent steady (DC) value that would deliver the same average power to a resistive load.
\end{definition}

\begin{methode}[Computing the mean value and the RMS value]
To compute the mean value and the RMS value of a continuous function $f$ on $[a,b]$:
\begin{enumerate}
\item Identify the interval $[a,b]$ and its length $b-a$.
\item \textbf{Mean value:} compute $\mu=\dfrac{1}{b-a}\displaystyle\int_a^b f(x)\,dx$. Find a primitive of $f$, evaluate between the bounds, then divide by $b-a$.
\item \textbf{RMS value:} first square the function, then compute $f_{\mathrm{rms}}=\sqrt{\dfrac{1}{b-a}\displaystyle\int_a^b \bigl(f(x)\bigr)^2\,dx}$.
\item For trigonometric squares, linearise first using $\sin^2x=\dfrac{1-\cos 2x}{2}$, $\cos^2x=\dfrac{1+\cos 2x}{2}$.
\item Give exact values first, then a decimal approximation if required.
\end{enumerate}
\textbf{Reflexes:} the RMS value is never smaller than the absolute mean value in physical applications; for a sinusoid $A\sin(\omega t)$ or $A\cos(\omega t)$ over a whole number of periods, the RMS value is $\dfrac{A}{\sqrt{2}}$ — memorise this, it saves time in the examination.
\end{methode}

\begin{exoresolu}[Mean and RMS values of an alternating current]
The current in a circuit at Bamenda varies as $i(t)=3\sin(100\pi t)$ amperes, where $t$ is in seconds. 
(a) Find the mean value of $i$ over the interval $\left[0,\dfrac{1}{100}\right]$. 
(b) Find the RMS value of $i$ over one period $T=\dfrac{1}{50}$ s.
\tcblower
\textbf{Solution.} The period of $i(t)=3\sin(100\pi t)$ is $T=\dfrac{2\pi}{100\pi}=\dfrac{1}{50}$ s, and $\left[0,\dfrac{1}{100}\right]$ is a half-period on which $i(t)\geq 0$.

\textbf{(a) Mean value over $\left[0,\frac{1}{100}\right]$.}
\begin{align*}
\mu &= \frac{1}{\frac{1}{100}-0}\int_0^{1/100} 3\sin(100\pi t)\,dt
=100\left[-\frac{3}{100\pi}\cos(100\pi t)\right]_0^{1/100}\\
&=100\cdot\frac{3}{100\pi}\bigl(-\cos\pi+\cos 0\bigr)
=\frac{3}{\pi}(1+1)=\frac{6}{\pi}\approx 1.91\ \mathrm{A}.
\end{align*}

\textbf{(b) RMS value over one period.}
\begin{align*}
i_{\mathrm{rms}}^2 &= \frac{1}{T}\int_0^{T} 9\sin^2(100\pi t)\,dt
=50\int_0^{1/50} 9\cdot\frac{1-\cos(200\pi t)}{2}\,dt\\
&=\frac{450}{2}\int_0^{1/50}\bigl(1-\cos(200\pi t)\bigr)\,dt
=225\left[t-\frac{\sin(200\pi t)}{200\pi}\right]_0^{1/50}\\
&=225\left(\frac{1}{50}-0\right)=\frac{9}{2}.
\end{align*}
Hence
\[
i_{\mathrm{rms}}=\sqrt{\frac{9}{2}}=\frac{3}{\sqrt{2}}=\frac{3\sqrt{2}}{2}\approx 2.12\ \mathrm{A},
\]
confirming the general rule $A/\sqrt{2}$.
\end{exoresolu}

\begin{aretenir}[Key formulas]
\begin{itemize}
\item Mean value: $\displaystyle \mu = \frac{1}{b-a}\int_a^b f(x)\,dx$.
\item RMS value: $\displaystyle f_{\mathrm{rms}} = \sqrt{\frac{1}{b-a}\int_a^b [f(x)]^2\,dx}$.
\item For $f(t)=A\sin(\omega t)$ or $A\cos(\omega t)$ over a complete period: $f_{\mathrm{rms}} = \dfrac{A}{\sqrt{2}}$.
\end{itemize}
\end{aretenir}

\begin{savaistu}[Alternating current in Cameroon]
The electricity distributed by ENEO in Cameroon is alternating current with a frequency of $50\ \mathrm{Hz}$ (period $0.02\ \mathrm{s}$). The voltage quoted on appliances ($230\ \mathrm{V}$) is the RMS value; the peak voltage is $230\sqrt{2}\approx 325\ \mathrm{V}$.
\end{savaistu}

\courssub{Arc Length of a Curve}

\begin{definition}[Arc length]
Let $y=f(x)$ be a continuously differentiable function on $[a,b]$. The length $L$ of the arc of the curve from $x=a$ to $x=b$ is
\[
L = \int_a^b \sqrt{1+\bigl(f'(x)\bigr)^2}\,dx.
\]
If the curve is given parametrically by $x=x(t),\ y=y(t)$ for $t\in[t_1,t_2]$, then
\[
L = \int_{t_1}^{t_2} \sqrt{\dot{x}^2+\dot{y}^2}\,dt.
\]
\end{definition}

\begin{methode}[Arc length of $y=f(x)$]
To find the length of the arc of $y=f(x)$ between $x=a$ and $x=b$:
\begin{enumerate}
\item Differentiate: compute $f'(x)$ and simplify $1+\bigl(f'(x)\bigr)^2$. Look for a \cle{perfect square} — this is the key trick in GCE questions.
\item Write $L=\displaystyle\int_a^b\sqrt{1+\bigl(f'(x)\bigr)^2}\,dx$.
\item If $1+(f')^2$ is a perfect square, take its square root and integrate directly.
\item Otherwise try a substitution (often hyperbolic or trigonometric).
\item For a parametric curve $x=x(t),\ y=y(t)$, use $L=\displaystyle\int_{t_1}^{t_2}\sqrt{\dot{x}^2+\dot{y}^2}\,dt$.
\end{enumerate}
\textbf{Reflex:} before integrating, always check whether $1+(f')^2$ factorises as a square, e.g. $1+\sinh^2x=\cosh^2x$, or $1+\frac{x}{2}+\frac{x^2}{16}=\left(1+\frac{x}{4}\right)^2$.
\end{methode}

\begin{exoresolu}[Length of a catenary]
A cable hangs between two pylons of the SONEL grid in the shape of the curve $y=\cosh x$. Find the exact length of the cable from $x=-1$ to $x=1$.
\tcblower
\textbf{Solution.} Here $f(x)=\cosh x$, so $f'(x)=\sinh x$. Then
\[ 1+\bigl(f'(x)\bigr)^2 = 1+\sinh^2 x = \cosh^2 x \quad\text{(hyperbolic identity)}. \]
Since $\cosh x>0$ for all $x$, $\sqrt{\cosh^2 x}=\cosh x$. Hence
\begin{align*}
L&=\int_{-1}^{1}\sqrt{1+\sinh^2 x}\,dx=\int_{-1}^{1}\cosh x\,dx
=\Bigl[\sinh x\Bigr]_{-1}^{1}\\
&=\sinh 1-\sinh(-1)=2\sinh 1 = e-\frac{1}{e}\approx 2.35 \text{ units}.
\end{align*}
The whole computation collapses to two lines because $1+\sinh^2x$ is a perfect square — exactly the situation the examiner designs.
\end{exoresolu}

\begin{exoresolu}[Arc length requiring algebraic preparation]
Find the length of the arc of the curve $y=\dfrac{x^2}{4}-\dfrac{\ln x}{2}$ from $x=1$ to $x=e$.
\tcblower
\textbf{Solution.} Differentiate: $y'=\dfrac{x}{2}-\dfrac{1}{2x}$. Then
\begin{align*}
1+(y')^2 &= 1+\left(\frac{x}{2}-\frac{1}{2x}\right)^2
=1+\frac{x^2}{4}-\frac{1}{2}+\frac{1}{4x^2}\\
&=\frac{x^2}{4}+\frac{1}{2}+\frac{1}{4x^2}
=\left(\frac{x}{2}+\frac{1}{2x}\right)^2.
\end{align*}
For $x>0$, $\dfrac{x}{2}+\dfrac{1}{2x}>0$, so
\begin{align*}
L&=\int_1^e\left(\frac{x}{2}+\frac{1}{2x}\right)dx
=\left[\frac{x^2}{4}+\frac{\ln x}{2}\right]_1^e\\
&=\left(\frac{e^2}{4}+\frac{1}{2}\right)-\left(\frac{1}{4}+0\right)
=\frac{e^2+1}{4}\approx 2.10 \text{ units}.
\end{align*}
Note the sign change of the middle term: $-2\cdot\frac{x}{2}\cdot\frac{1}{2x}=-\frac12$, which turns $1$ into $+\frac12$ and rebuilds a perfect square.
\end{exoresolu}

\begin{aretenir}[Arc length formula]
\[
L = \int_a^b \sqrt{1+(y')^2}\,dx \quad\text{or}\quad L = \int_{t_1}^{t_2} \sqrt{\dot{x}^2+\dot{y}^2}\,dt.
\]
Always simplify $1+(y')^2$ before integrating; it is often a perfect square in examination questions.
\end{aretenir}

\courssub{Areas of Surfaces of Revolution}

\begin{definition}[Surface area of revolution]
When the arc of the curve $y=f(x)\ge 0$ from $x=a$ to $x=b$ is rotated through $2\pi$ radians about the $x$-axis, the area $S$ of the surface generated is
\[
S = 2\pi\int_a^b y\sqrt{1+\bigl(f'(x)\bigr)^2}\,dx = 2\pi\int_a^b y\,ds,
\]
where $ds=\sqrt{1+(y')^2}\,dx$ is the arc length element. For rotation about the $y$-axis, the formula becomes
\[
S = 2\pi\int_a^b x\sqrt{1+\bigl(f'(x)\bigr)^2}\,dx.
\]
\end{definition}

\begin{methode}[Surface area of revolution about the $x$-axis]
To find the area of the surface generated when the arc $y=f(x)$, $a\leq x\leq b$, is rotated about the $x$-axis:
\begin{enumerate}
\item Compute $f'(x)$ and simplify $1+(f')^2$ (again, look for a perfect square).
\item Write $S=2\pi\displaystyle\int_a^b y\sqrt{1+\bigl(f'(x)\bigr)^2}\,dx$.
\item The integrand is $y$ times the \cle{arc length element}; combine and simplify before integrating.
\item For rotation about the $y$-axis, replace $y$ by $x$: $S=2\pi\displaystyle\int_a^b x\sqrt{1+(f')^2}\,dx$.
\end{enumerate}
\textbf{Reflex:} the formula is $2\pi\int y\,ds$, not $2\pi\int y\,dx$ — forgetting the square-root factor is the most common error.
\end{methode}

\begin{popfigure}
\begin{tikzpicture}[scale=0.8]
\draw[->] (-0.5,0) -- (4,0) node[right] {$x$};
\draw[->] (0,-0.5) -- (0,3) node[above] {$y$};
\draw[thick, popblue] plot[domain=0:3, samples=100] (\x, {0.5*\x*\x});
\draw[dashed] (3,0) node[below] {$b$} -- (3,4.5);
\draw[dashed] (0,0) node[below left] {$0$} -- (0,0);
\draw[poporange, thick] (1.5,1.125) -- (1.5,0) node[below] {$x$};
\draw[poporange, thick] (0,1.125) -- (1.5,1.125) node[left] {$y$};
\draw[popgreen, decorate, decoration={brace, amplitude=5pt}] (3.2,0) -- (3.2,4.5) node[midway, right] {$y=f(x)$};
\end{tikzpicture}
\legende{Surface of revolution: the arc $y=f(x)$ rotated about the $x$-axis generates a surface.}
\end{popfigure}

\begin{exoresolu}[Surface generated by a parabola arc]
The arc of the curve $y=\dfrac{x^2}{2}$ between $x=0$ and $x=\sqrt{3}$ is rotated through $2\pi$ about the $x$-axis. Find the area of the surface generated.
\tcblower
\textbf{Solution.} We have $y' = x$, so $1+(y')^2 = 1+x^2$ and
\[ S = 2\pi\int_0^{\sqrt{3}} \frac{x^2}{2}\sqrt{1+x^2}\,dx = \pi\int_0^{\sqrt{3}} x^2\sqrt{1+x^2}\,dx. \]
Use the substitution $x=\sinh t$, $dx=\cosh t\,dt$, $\sqrt{1+x^2}=\cosh t$. When $x=0$, $t=0$; when $x=\sqrt{3}$, $t=\operatorname{arsinh}\sqrt{3}=\ln(2+\sqrt{3})$. Then
\begin{align*}
S&=\pi\int_0^{t_1}\sinh^2 t\cosh^2 t\,dt
=\frac{\pi}{4}\int_0^{t_1}\sinh^2(2t)\,dt
=\frac{\pi}{8}\int_0^{t_1}\bigl(\cosh 4t-1\bigr)dt\\
&=\frac{\pi}{8}\left[\frac{\sinh 4t}{4}-t\right]_0^{t_1}.
\end{align*}
With $t_1=\ln(2+\sqrt{3})$: $\sinh t_1=\sqrt{3}$, $\cosh t_1=2$, so $\sinh 2t_1=2\sinh t_1\cosh t_1=4\sqrt{3}$ and $\cosh 2t_1=1+2\sinh^2t_1=7$. Then $\sinh 4t_1=2\sinh 2t_1\cosh 2t_1=2\cdot 4\sqrt{3}\cdot 7=56\sqrt{3}$. Hence
\[ S=\frac{\pi}{8}\left(\frac{56\sqrt{3}}{4}-\ln(2+\sqrt{3})\right)
=\frac{\pi}{8}\Bigl(14\sqrt{3}-\ln(2+\sqrt{3})\Bigr)\approx 9.34\ \text{square units}. \]
\end{exoresolu}

\begin{aretenir}[Surface area formulas]
\begin{itemize}
\item About $x$-axis: $S = 2\pi\displaystyle\int_a^b y\sqrt{1+(y')^2}\,dx$.
\item About $y$-axis: $S = 2\pi\displaystyle\int_a^b x\sqrt{1+(y')^2}\,dx$.
\item Parametric: $S = 2\pi\displaystyle\int_{t_1}^{t_2} y(t)\sqrt{\dot{x}^2+\dot{y}^2}\,dt$ (about $x$-axis).
\end{itemize}
\end{aretenir}

\courssub{Volumes of Solids of Revolution}

\begin{definition}[Volume of revolution (disc/washer method)]
When the region under the curve $y=f(x)\ge 0$ from $x=a$ to $x=b$ is rotated through $2\pi$ radians about the $x$-axis, the volume $V$ of the solid generated is
\[
V = \pi\int_a^b \bigl(f(x)\bigr)^2\,dx.
\]
If the region lies between two curves $y_1(x)\ge y_2(x)\ge 0$, the washer method gives
\[
V = \pi\int_a^b \bigl(y_1^2 - y_2^2\bigr)\,dx.
\]
For rotation about the $y$-axis, either invert to $x=g(y)$ and use $V=\pi\displaystyle\int_c^d x^2\,dy$, or use cylindrical shells: $V=2\pi\displaystyle\int_a^b x\,f(x)\,dx$.
\end{definition}

\begin{methode}[Volume of revolution about the $x$-axis (disc/washer)]
To find the volume generated when the region under $y=f(x)$, from $x=a$ to $x=b$, is rotated through $2\pi$ radians about the $x$-axis:
\begin{enumerate}
\item Sketch the region and identify the axis of rotation.
\item Write the disc formula $V=\pi\displaystyle\int_a^b y^2\,dx$.
\item Square the function \emph{before} integrating; linearise trigonometric squares.
\item If the region lies \emph{between} two curves $y_1\geq y_2\geq 0$, use washers: $V=\pi\displaystyle\int_a^b\left(y_1^2-y_2^2\right)dx$.
\item For rotation about the $y$-axis, either invert to $x=g(y)$ and use $V=\pi\displaystyle\int_c^d x^2\,dy$, or use cylindrical shells.
\end{enumerate}
\textbf{Reflexes:} do not forget the factor $\pi$; the integrand is $y^2$, never $y$; state the units cubed.
\end{methode}

\begin{popfigure}
\begin{tikzpicture}[scale=0.7]
\draw[->] (-0.5,0) -- (4,0) node[right] {$x$};
\draw[->] (0,-0.5) -- (0,3) node[above] {$y$};
\draw[thick, popblue] plot[domain=0:3, samples=100] (\x, {0.5*\x*\x});
\draw[dashed] (3,0) node[below] {$b$} -- (3,4.5);
\draw[poporange, thick] (1.5,0) -- (1.5,1.125);
\draw[poporange, thick] (0,1.125) -- (1.5,1.125);
\draw[popgreen, decorate, decoration={brace, amplitude=5pt}] (3.2,0) -- (3.2,4.5) node[midway, right] {$y=f(x)$};
\draw[poppink, fill=poppink!30] (1.5,0) rectangle (1.6,1.125);
\end{tikzpicture}
\legende{Volume of revolution: a thin disc of radius $y$ and thickness $dx$ gives volume $\pi y^2\,dx$.}
\end{popfigure}

\begin{exoresolu}[Volume of a solid with exponential function]
The region bounded by the curve $y=\sqrt{x}\,e^{x}$, the $x$-axis and the lines $x=0$ and $x=1$ is rotated through $2\pi$ about the $x$-axis. Find the exact volume of the solid generated.
\tcblower
\textbf{Solution.} By the disc method with $y^2 = x e^{2x}$:
\[ V=\pi\int_0^1 x e^{2x}\,dx. \]
Integrate by parts with $u=x$, $dv=e^{2x}dx$, so $du=dx$, $v=\dfrac{e^{2x}}{2}$:
\begin{align*}
\int_0^1 x e^{2x}\,dx &= \left[\frac{x e^{2x}}{2}\right]_0^1-\int_0^1\frac{e^{2x}}{2}\,dx
=\frac{e^2}{2}-\left[\frac{e^{2x}}{4}\right]_0^1\\
&=\frac{e^2}{2}-\frac{e^2}{4}+\frac{1}{4}=\frac{e^2+1}{4}.
\end{align*}
Therefore
\[ V=\frac{\pi\left(e^2+1\right)}{4}\approx 6.61\ \text{cubic units}. \]
The squaring step removed the square root immediately — always square $y$ first and simplify before choosing an integration technique.
\end{exoresolu}

\begin{aretenir}[Volume formulas]
\begin{itemize}
\item Disc (about $x$-axis): $V = \pi\displaystyle\int_a^b y^2\,dx$.
\item Washer (about $x$-axis): $V = \pi\displaystyle\int_a^b (y_1^2-y_2^2)\,dx$.
\item Shell (about $y$-axis): $V = 2\pi\displaystyle\int_a^b x\,y\,dx$.
\end{itemize}
\end{aretenir}

\courssub{Centroids of Plane Figures}

\begin{definition}[Centroid of a plane region]
The \cle{centroid} (or geometric centre) $G(\bar{x},\bar{y})$ of a plane region of area $A$ is given by
\[
\bar{x} = \frac{1}{A}\iint_R x\,dA,\qquad \bar{y} = \frac{1}{A}\iint_R y\,dA.
\]
For a region bounded by $y=f(x)\ge 0$, the $x$-axis, and the lines $x=a$, $x=b$:
\[
A = \int_a^b f(x)\,dx,\quad
\bar{x} = \frac{1}{A}\int_a^b x\,f(x)\,dx,\quad
\bar{y} = \frac{1}{2A}\int_a^b \bigl(f(x)\bigr)^2\,dx.
\]
If the region is between two curves $y_1(x)\ge y_2(x)$:
\[
A = \int_a^b (y_1-y_2)\,dx,\quad
\bar{x} = \frac{1}{A}\int_a^b x(y_1-y_2)\,dx,\quad
\bar{y} = \frac{1}{2A}\int_a^b (y_1^2-y_2^2)\,dx.
\]
\end{definition}

\begin{methode}[Centroid of a region under a curve]
To locate the centroid $G(\bar{x},\bar{y})$ of the region under $y=f(x)$, $a\leq x\leq b$:
\begin{enumerate}
\item Compute the area $A=\displaystyle\int_a^b f(x)\,dx$ first — you will divide by it twice.
\item Compute $\bar{x}=\dfrac{1}{A}\displaystyle\int_a^b x\,f(x)\,dx$ (often by parts).
\item Compute $\bar{y}=\dfrac{1}{2A}\displaystyle\int_a^b \bigl(f(x)\bigr)^2dx$ — note the factor $\dfrac12$ and the square.
\item Exploit symmetry: if the region is symmetric about a line, the centroid lies on it.
\end{enumerate}
\textbf{Reflex:} for a region between two curves, replace $f$ by $y_1-y_2$ in $A$ and $\bar x$, and use $\bar y=\dfrac{1}{2A}\displaystyle\int_a^b\left(y_1^2-y_2^2\right)dx$.
\end{methode}

\begin{exoresolu}[Centroid under a parabola]
Find the coordinates of the centroid of the region bounded by the curve $y=4-x^2$ and the $x$-axis.
\tcblower
\textbf{Solution.} The curve meets the $x$-axis at $x=\pm 2$. By symmetry about the $y$-axis, $\bar{x}=0$.

\textbf{Area:}
\[ A=\int_{-2}^{2}\left(4-x^2\right)dx=\left[4x-\frac{x^3}{3}\right]_{-2}^{2}
=\left(8-\frac{8}{3}\right)-\left(-8+\frac{8}{3}\right)=\frac{32}{3}. \]

\textbf{Ordinate of the centroid:}
\begin{align*}
\int_{-2}^{2}\left(4-x^2\right)^2dx &=\int_{-2}^{2}\left(16-8x^2+x^4\right)dx
=\left[16x-\frac{8x^3}{3}+\frac{x^5}{5}\right]_{-2}^{2}\\
&=2\left(32-\frac{64}{3}+\frac{32}{5}\right)=\frac{512}{15}.
\end{align*}
Therefore
\[ \bar{y}=\frac{1}{2A}\cdot\frac{512}{15}=\frac{512}{15}\cdot\frac{3}{64}=\frac{8}{5}. \]
The centroid is $G\left(0,\dfrac{8}{5}\right)$. Check: $\bar y$ lies below the maximum height $4$ and above $0$, and for a parabolic spandrel-type region the value $\tfrac{2}{5}$ of the height is plausible.
\end{exoresolu}

\begin{aretenir}[Centroid formulas]
\begin{itemize}
\item Region under $y=f(x)$: $A=\int_a^b f\,dx$, $\bar{x}=\frac{1}{A}\int_a^b xf\,dx$, $\bar{y}=\frac{1}{2A}\int_a^b f^2\,dx$.
\item Region between $y_1$ and $y_2$: $A=\int_a^b (y_1-y_2)\,dx$, $\bar{x}=\frac{1}{A}\int_a^b x(y_1-y_2)\,dx$, $\bar{y}=\frac{1}{2A}\int_a^b (y_1^2-y_2^2)\,dx$.
\item Symmetry: if the region is symmetric about a line, the centroid lies on that line.
\end{itemize}
\end{aretenir}

\courssub{Centre of Gravity of a Solid of Revolution}

\begin{definition}[Centre of gravity of a solid of revolution]
For a solid of uniform density generated by rotating the region under $y=f(x)\ge 0$, $a\le x\le b$, about the $x$-axis, the centre of gravity lies on the axis of rotation by symmetry. Its $x$-coordinate $\bar{x}$ is given by
\[
\bar{x} = \frac{\int_a^b x\,y^2\,dx}{\int_a^b y^2\,dx}.
\]
The factor $\pi$ cancels between the moment and the volume.
\end{definition}

\begin{methode}[Centre of gravity of a solid of revolution about the $x$-axis]
For the solid generated by rotating the region under $y=f(x)$, $a\leq x\leq b$, about the $x$-axis (uniform density):
\begin{enumerate}
\item By symmetry, the centre of gravity lies on the axis of rotation: $\bar y = 0$.
\item Compute the volume $V=\pi\displaystyle\int_a^b y^2\,dx$.
\item Compute the moment $M=\pi\displaystyle\int_a^b x\,y^2\,dx$ (disc of radius $y$ at position $x$).
\item Then $\bar{x}=\dfrac{M}{V}=\dfrac{\displaystyle\int_a^b x y^2\,dx}{\displaystyle\int_a^b y^2\,dx}$ — the $\pi$ cancels.
\end{enumerate}
\textbf{Reflex:} same integrand $y^2$ as for the volume, with an extra factor $x$ in the numerator; expect integration by parts.
\end{methode}

\begin{exoresolu}[Centre of gravity of a hemisphere]
The region bounded by $y=\sqrt{r^2-x^2}$, the $x$-axis, for $0\leq x\leq r$, is rotated about the $x$-axis to generate a solid hemisphere of radius $r$. Show that its centre of gravity is at distance $\dfrac{3r}{8}$ from the flat face.
\tcblower
\textbf{Solution.} Here $y^2=r^2-x^2$. By symmetry the centre of gravity lies on the $x$-axis.

\textbf{Volume:}
\[ V=\pi\int_0^r\left(r^2-x^2\right)dx=\pi\left[r^2x-\frac{x^3}{3}\right]_0^r
=\pi\left(r^3-\frac{r^3}{3}\right)=\frac{2\pi r^3}{3}, \]
which is indeed half the volume of a sphere.

\textbf{Moment about the $y$-axis plane:}
\begin{align*}
M&=\pi\int_0^r x\left(r^2-x^2\right)dx=\pi\int_0^r\left(r^2x-x^3\right)dx\\
&=\pi\left[\frac{r^2x^2}{2}-\frac{x^4}{4}\right]_0^r
=\pi\left(\frac{r^4}{2}-\frac{r^4}{4}\right)=\frac{\pi r^4}{4}.
\end{align*}

\textbf{Centre of gravity:}
\[ \bar{x}=\frac{M}{V}=\frac{\pi r^4/4}{2\pi r^3/3}=\frac{r}{4}\cdot\frac{3}{2}=\frac{3r}{8}. \]
Hence the centre of gravity of a solid hemisphere is at $\dfrac{3r}{8}$ from its plane face, as required.
\end{exoresolu}

\begin{aretenir}[Centre of gravity of a solid of revolution]
For rotation about the $x$-axis: $\bar{x} = \dfrac{\int_a^b x y^2\,dx}{\int_a^b y^2\,dx}$, $\bar{y}=0$.
For rotation about the $y$-axis: $\bar{y} = \dfrac{\int_c^d y x^2\,dy}{\int_c^d x^2\,dy}$, $\bar{x}=0$.
\end{aretenir}

\courssub{Theorems of Pappus}

\begin{propriete}[Pappus's theorems]
Let a plane region (or arc) be rotated about an external axis in its plane, not intersecting the region (or arc).
\begin{enumerate}
\item \textbf{First theorem (volume):} The volume $V$ of the solid generated is $V = 2\pi\,\bar{d}\,A$, where $A$ is the area of the region and $\bar d$ the distance of its centroid from the axis of rotation.
\item \textbf{Second theorem (surface area):} The area $S$ of the surface generated is $S = 2\pi\,\bar{d}\,L$, where $L$ is the length of the arc and $\bar d$ the distance of the centroid \emph{of the arc} from the axis.
\end{enumerate}
\end{propriete}

\begin{methode}[Applying Pappus's theorems]
\begin{enumerate}
\item \textbf{Direct use:} given $A$ (or $L$) and $\bar d$, compute $V$ (or $S$) without any integration.
\item \textbf{Reverse use:} given a known volume or surface (sphere, torus, cone), \emph{deduce} the position of the centroid: $\bar d=\dfrac{V}{2\pi A}$ or $\bar d=\dfrac{S}{2\pi L}$.
\item \textbf{Check:} the region (or arc) must not cross the axis of rotation; otherwise the theorem does not apply.
\end{enumerate}
\textbf{Reflex:} Pappus is the fastest route for circles, rectangles and semicircles rotated about external axes; always check that the region does not intersect the axis.
\end{methode}

\begin{popfigure}
\begin{tikzpicture}[scale=0.7]
\draw[->] (-1,0) -- (6,0) node[right] {$x$};
\draw[->] (0,-1) -- (0,4) node[above] {$y$};
\draw[thick, popblue] (3,0) circle (1);
\draw[dashed] (0,0) -- (0,3);
\draw[poporange, thick] (3,0) -- (0,0) node[midway, below] {$d=3$};
\draw[popgreen, fill=popgreen!20] (3,0) circle (1);
\node at (3,0) {$\times$};
\node at (3,1.3) {$A=\pi r^2$};
\end{tikzpicture}
\legende{Pappus's first theorem: a disc of radius $r$ centred at distance $d$ from the $y$-axis generates a torus of volume $V=2\pi d\cdot \pi r^2$.}
\end{popfigure}

\begin{exoresolu}[Volume of a torus and centroid of a semicircular arc]
(a) A circular disc of radius $2$ cm, centred at distance $5$ cm from an axis in its plane, is rotated about that axis. Find the volume of the torus (anchor ring) generated.

(b) A semicircular arc of radius $r$ is rotated about its diameter to generate a sphere of surface area $4\pi r^2$. Use Pappus's second theorem to find the distance of the centroid of the semicircular arc from the diameter.
\tcblower
\textbf{Solution. (a)} The disc has area $A=\pi(2)^2=4\pi$ cm$^2$ and its centroid is its centre, at distance $\bar d=5$ cm from the axis. The disc does not cross the axis since $5>2$. By Pappus's first theorem:
\[ V=2\pi\bar d\,A=2\pi(5)(4\pi)=40\pi^2\ \text{cm}^3\approx 394.8\ \text{cm}^3. \]

\textbf{(b)} The semicircular arc has length $L=\pi r$. Rotated about its diameter, it sweeps the sphere of area $S=4\pi r^2$. If $\bar d$ is the distance of the centroid of the \emph{arc} from the diameter, Pappus's second theorem gives
\[ S=2\pi\bar d\,L \;\Longrightarrow\; 4\pi r^2 = 2\pi\bar d(\pi r)
\;\Longrightarrow\; \bar d=\frac{4\pi r^2}{2\pi^2 r}=\frac{2r}{\pi}. \]
Hence the centroid of a semicircular arc lies at distance $\dfrac{2r}{\pi}\approx 0.637\,r$ from the diameter. (Compare: the centroid of a semicircular \emph{lamina} is at $\dfrac{4r}{3\pi}\approx 0.424\,r$ — do not confuse the two.)
\end{exoresolu}

\begin{savaistu}[Pappus of Alexandria]
Pappus of Alexandria (c. 290–350 AD) was one of the last great Greek mathematicians. His theorems, also known as Guldin's rules (after Paul Guldin, 1577–1643), allow elegant computation of volumes and surfaces of revolution without integration. They are particularly useful for checking integration results.
\end{savaistu}

\begin{attention}[Frequent errors in the examination]
\begin{itemize}
\item Forgetting the factor $\pi$ in volumes, or the factor $2\pi$ in surface areas and Pappus's theorems.
\item Integrating $y$ instead of $y^2$ in the volume formula.
\item Omitting the arc-length element $\sqrt{1+(y')^2}$ in surface area integrals.
\item Confusing the centroid of an arc ($\tfrac{2r}{\pi}$ for a semicircle) with that of a lamina ($\tfrac{4r}{3\pi}$).
\item Applying Pappus when the region crosses the axis of rotation — the theorem then fails.
\item For the RMS value, forgetting the outer square root at the end.
\end{itemize}
\end{attention}

\begin{exercice}[Synthesis exercise]
\begin{enumerate}
\item The current in an AC circuit is given by $i(t)=5\sin(100\pi t)$ A. Find (a) the mean value over a half-cycle, (b) the RMS value over a full cycle.
\item Find the length of the curve $y=\frac{1}{3}(x^2+2)^{3/2}$ from $x=0$ to $x=1$.
\item The arc of the curve $y=\frac{x^3}{3}$ from $x=0$ to $x=1$ is rotated about the $x$-axis. Find the area of the surface generated.
\item The region bounded by $y=x^2$, $y=0$, $x=2$ is rotated about the $y$-axis. Find the volume of the solid generated (use shells).
\item Find the centroid of the region bounded by $y=\sin x$, $y=0$, $x=0$, $x=\pi$.
\item A semicircular lamina of radius $r$ is rotated about its diameter. Use Pappus's first theorem to find the volume of the sphere generated, and hence deduce the centroid of the semicircular lamina.
\end{enumerate}
\end{exercice}

\begin{corrige}[Exercise 1]
(a) Half-cycle interval: $[0, 1/100]$. $\mu = \frac{1}{1/100}\int_0^{1/100} 5\sin(100\pi t)\,dt = 100\cdot\frac{5}{100\pi}(-\cos\pi+\cos0) = \frac{10}{\pi}\approx 3.18\ \mathrm{A}$.
(b) Full cycle $T=1/50$. $i_{\mathrm{rms}}^2 = 50\int_0^{1/50} 25\sin^2(100\pi t)\,dt = 50\cdot\frac{25}{2}\cdot\frac{1}{50} = \frac{25}{2}$, so $i_{\mathrm{rms}} = \frac{5}{\sqrt{2}}\approx 3.54\ \mathrm{A}$.
\end{corrige}

\begin{corrige}[Exercise 2]
$y' = x\sqrt{x^2+2}$. $1+(y')^2 = 1+x^2(x^2+2) = x^4+2x^2+1 = (x^2+1)^2$. $L = \int_0^1 (x^2+1)\,dx = \left[\frac{x^3}{3}+x\right]_0^1 = \frac{4}{3}$.
\end{corrige}

\begin{corrige}[Exercise 3]
$y' = x^2$, $1+(y')^2 = 1+x^4$. $S = 2\pi\int_0^1 \frac{x^3}{3}\sqrt{1+x^4}\,dx$. Let $u=1+x^4$, $du=4x^3dx$. $S = \frac{2\pi}{3}\int_1^2 \frac{\sqrt{u}}{4}\,du = \frac{\pi}{6}\cdot\frac{2}{3}u^{3/2}\big|_1^2 = \frac{\pi}{9}(2\sqrt{2}-1)$.
\end{corrige}

\begin{corrige}[Exercise 4]
Shell method: $V = 2\pi\int_0^2 x\cdot x^2\,dx = 2\pi\int_0^2 x^3\,dx = 2\pi\left[\frac{x^4}{4}\right]_0^2 = 8\pi$.
\end{corrige}

\begin{corrige}[Exercise 5]
$A = \int_0^\pi \sin x\,dx = 2$. $\bar{x} = \frac{1}{2}\int_0^\pi x\sin x\,dx = \frac{1}{2}\bigl([-x\cos x]_0^\pi + \int_0^\pi \cos x\,dx\bigr) = \frac{1}{2}(\pi + 0) = \frac{\pi}{2}$. $\bar{y} = \frac{1}{4}\int_0^\pi \sin^2 x\,dx = \frac{1}{4}\int_0^\pi \frac{1-\cos 2x}{2}\,dx = \frac{1}{8}\bigl[x-\frac{\sin 2x}{2}\bigr]_0^\pi = \frac{\pi}{8}$. Centroid: $(\pi/2, \pi/8)$.
\end{corrige}

\begin{corrige}[Exercise 6]
Semicircular lamina area $A = \frac{1}{2}\pi r^2$. Rotated about diameter gives sphere volume $V = \frac{4}{3}\pi r^3$. By Pappus: $V = 2\pi\bar{d} A \Rightarrow \frac{4}{3}\pi r^3 = 2\pi\bar{d}\cdot\frac{1}{2}\pi r^2 \Rightarrow \bar{d} = \frac{4r}{3\pi}$. This is the centroid of the semicircular lamina.
\end{corrige}

\newpage

\coursec{Exercises}

\coursec{Applications of Definite Integrals}

\courssub{Mean Value and RMS Value of a Function}

\begin{definition}[Mean Value of a Function]
Let $f$ be a continuous function on the closed interval $[a,b]$. The \cle{mean value} (or average value) of $f$ over $[a,b]$ is defined by
\[
\overline{f} = \frac{1}{b-a}\int_{a}^{b} f(x)\;dx.
\]
\end{definition}

\begin{propriete}[Geometrical Interpretation]
The mean value $\overline{f}$ is the height of a rectangle with base $[a,b]$ that has the same area as the region under the curve $y=f(x)$ between $x=a$ and $x=b$. If $f(x)\ge 0$ on $[a,b]$, then
\[
\text{Area} = \int_{a}^{b} f(x)\;dx = \overline{f}\,(b-a).
\]
\end{propriete}

\begin{definition}[Root Mean Square (RMS) Value]
The \cle{root mean square} (RMS) value of a continuous function $f$ on $[a,b]$ is
\[
f_{\mathrm{rms}} = \sqrt{\frac{1}{b-a}\int_{a}^{b} \bigl[f(x)\bigr]^{2}\;dx}.
\]
\end{definition}

\begin{aretenir}[Why RMS for Alternating Currents?]
For an alternating current $i(t)$, the mean value over a full period is zero (equal positive and negative halves), yet it delivers power. The instantaneous power in a resistor $R$ is $P(t)=i(t)^{2}R$. The mean power is $\overline{P}=R\,\overline{i^{2}}$. The RMS current $I_{\mathrm{rms}}=\sqrt{\overline{i^{2}}}$ is the steady (DC) current that would dissipate the same mean power in the same resistor. Hence RMS is the \emph{effective} value.
\end{aretenir}

\begin{exemplebox}[Mean and RMS of a Sinusoidal Function]
Find the mean value and the RMS value of $f(t)=5\sin(100\pi t)$ over one period.
\begin{align*}
\text{Period } T &= \frac{2\pi}{100\pi} = 0.02\ \mathrm{s}. \\
\overline{f} &= \frac{1}{T}\int_{0}^{T} 5\sin(100\pi t)\;dt = 0. \\
f_{\mathrm{rms}} &= \sqrt{\frac{1}{T}\int_{0}^{T} 25\sin^{2}(100\pi t)\;dt}
= \sqrt{\frac{25}{T}\int_{0}^{T} \frac{1-\cos(200\pi t)}{2}\;dt}
= \sqrt{\frac{25}{2}} = \frac{5}{\sqrt{2}}\ \mathrm{A}.
\end{align*}
\end{exemplebox}

\begin{methode}[Mean Value of a Rectified Waveform]
To find the mean value of $|f(t)|$ over a period $T$ when $f(t)$ is sinusoidal:
\begin{enumerate}
\item Use symmetry: $\displaystyle \frac{1}{T}\int_{0}^{T}|f(t)|\;dt = \frac{2}{T}\int_{0}^{T/2}f(t)\;dt$ (if $f\ge 0$ on $[0,T/2]$).
\item For $f(t)=I_{0}\sin(\omega t)$, $\displaystyle \overline{|f|} = \frac{2I_{0}}{\pi}$.
\end{enumerate}
\end{methode}

\courssub{Arc Length of a Curve}

\begin{propriete}[Arc Length Formula -- Cartesian Form]
If $y=f(x)$ has a continuous derivative on $[a,b]$, the length $L$ of the curve from $x=a$ to $x=b$ is
\[
L = \int_{a}^{b} \sqrt{1+\left(\frac{dy}{dx}\right)^{2}}\;dx.
\]
If the curve is given as $x=g(y)$ on $[c,d]$, then
\[
L = \int_{c}^{d} \sqrt{1+\left(\frac{dx}{dy}\right)^{2}}\;dy.
\]
\end{propriete}

\begin{propriete}[Arc Length -- Parametric Form]
If a curve is given parametrically by $x=x(t)$, $y=y(t)$ for $\alpha\le t\le \beta$, with continuous derivatives and not simultaneously zero, then
\[
L = \int_{\alpha}^{\beta} \sqrt{\left(\frac{dx}{dt}\right)^{2}+\left(\frac{dy}{dt}\right)^{2}}\;dt.
\]
\end{propriete}

\begin{exemplebox}[Length of a Parabolic Arc]
Find the length of the curve $y=\dfrac{x^{2}}{2}$ from $x=0$ to $x=1$.
\[
\frac{dy}{dx}=x,\quad ds=\sqrt{1+x^{2}}\;dx,\quad L=\int_{0}^{1}\sqrt{1+x^{2}}\;dx.
\]
Use substitution $x=\sinh u$, $dx=\cosh u\,du$, $\sqrt{1+x^{2}}=\cosh u$. When $x=0$, $u=0$; when $x=1$, $u=\operatorname{arsinh}1=\ln(1+\sqrt{2})$.
\[
L = \int_{0}^{\ln(1+\sqrt{2})}\cosh^{2}u\;du = \int_{0}^{\ln(1+\sqrt{2})}\frac{\cosh 2u+1}{2}\;du
= \frac{1}{2}\left[\frac{\sinh 2u}{2}+u\right]_{0}^{\ln(1+\sqrt{2})}.
\]
Since $\sinh 2u = 2\sinh u\cosh u = 2x\sqrt{1+x^{2}}$, at $x=1$: $\sinh 2u = 2\sqrt{2}$. Thus
\[
L = \frac{1}{2}\left(\frac{2\sqrt{2}}{2}+\ln(1+\sqrt{2})\right) = \frac{\sqrt{2}}{2}+\frac{1}{2}\ln(1+\sqrt{2}).
\]
\end{exemplebox}

\begin{exemplebox}[Length of a Catenary]
A cable hangs in the shape $y=\cosh x$ between $x=-1$ and $x=1$. Find its length.
\[
\frac{dy}{dx}=\sinh x,\quad ds=\sqrt{1+\sinh^{2}x}\;dx = \cosh x\;dx.
\]
\[
L = \int_{-1}^{1}\cosh x\;dx = \bigl[\sinh x\bigr]_{-1}^{1} = 2\sinh 1 \approx 2.35\ \mathrm{(3\ s.f.)}.
\]
\end{exemplebox}

\courssub{Surfaces and Volumes of Revolution}

\begin{propriete}[Volume of Revolution -- Disc/Washer Method (about $x$-axis)]
The volume generated by rotating the region bounded by $y=f(x)$, the $x$-axis, $x=a$ and $x=b$ ($f(x)\ge 0$) about the $x$-axis is
\[
V = \pi\int_{a}^{b} \bigl[f(x)\bigr]^{2}\;dx.
\]
If the region is between two curves $y=f(x)$ (upper) and $y=g(x)$ (lower), $f(x)\ge g(x)\ge 0$, the washer method gives
\[
V = \pi\int_{a}^{b} \bigl([f(x)]^{2}-[g(x)]^{2}\bigr)\;dx.
\]
\end{propriete}

\begin{propriete}[Volume of Revolution -- Shell Method (about $y$-axis)]
The volume generated by rotating the region under $y=f(x)$ ($f(x)\ge 0$) from $x=a$ to $x=b$ about the $y$-axis is
\[
V = 2\pi\int_{a}^{b} x\,f(x)\;dx.
\]
\end{propriete}

\begin{propriete}[Surface Area of Revolution (about $x$-axis)]
The area of the curved surface generated by rotating the arc $y=f(x)$ ($f(x)\ge 0$) from $x=a$ to $x=b$ about the $x$-axis is
\[
S = 2\pi\int_{a}^{b} f(x)\sqrt{1+\left(\frac{dy}{dx}\right)^{2}}\;dx = 2\pi\int_{a}^{b} y\;ds.
\]
If the rotation is about the $y$-axis, $S = 2\pi\int_{a}^{b} x\;ds$.
\end{propriete}

\begin{attention}[Common Mistakes]
\begin{itemize}
\item Forgetting the factor $2\pi$ in the surface area formula (it is $2\pi y\,ds$, not $\pi y\,ds$).
\item Using $dx$ instead of $ds$ in the surface area integral.
\item Confusing the disc method (slices perpendicular to axis) with the shell method (slices parallel to axis).
\item When using the washer method, always subtract the inner radius squared from the outer radius squared: $\pi(R^{2}-r^{2})$, not $\pi(R-r)^{2}$.
\end{itemize}
\end{attention}

\begin{exemplebox}[Volume and Surface of a Paraboloid]
The region under $y=x^{2}$, $0\le x\le 2$, is rotated about the $x$-axis.
\begin{enumerate}
\item Volume: $V = \pi\int_{0}^{2} (x^{2})^{2}\;dx = \pi\left[\frac{x^{5}}{5}\right]_{0}^{2} = \dfrac{32\pi}{5}$.
\item Surface area: $\dfrac{dy}{dx}=2x$, $ds=\sqrt{1+4x^{2}}\;dx$.
$S = 2\pi\int_{0}^{2} x^{2}\sqrt{1+4x^{2}}\;dx$.
Use $2x=\sinh u \Rightarrow x=\frac{1}{2}\sinh u$, $dx=\frac{1}{2}\cosh u\,du$, $\sqrt{1+4x^{2}}=\cosh u$.
When $x=0$, $u=0$; when $x=2$, $u=\operatorname{arsinh}4=\ln(4+\sqrt{17})$.
\[
S = 2\pi\int_{0}^{\ln(4+\sqrt{17})} \frac{1}{4}\sinh^{2}u \cdot \cosh u \cdot \frac{1}{2}\cosh u\;du
= \frac{\pi}{4}\int_{0}^{\ln(4+\sqrt{17})} \sinh^{2}u\cosh^{2}u\;du.
\]
Using $\sinh u\cosh u = \frac{1}{2}\sinh 2u$, the integrand becomes $\frac{1}{4}\sinh^{2}2u = \frac{1}{4}\cdot\frac{\cosh 4u-1}{2} = \frac{\cosh 4u-1}{8}$.
\[
S = \frac{\pi}{32}\int_{0}^{\ln(4+\sqrt{17})} (\cosh 4u-1)\;du
= \frac{\pi}{32}\left[\frac{\sinh 4u}{4}-u\right]_{0}^{\ln(4+\sqrt{17})}.
\]
With $\sinh 4u = 2\sinh 2u\cosh 2u = 4\sinh u\cosh u(2\cosh^{2}u-1)$ and $\sinh u=4$, $\cosh u=\sqrt{17}$, we obtain the exact value.
\end{enumerate}
\end{exemplebox}

\courssub{Centroids of Plane Figures and Centres of Gravity of Solids of Revolution}

\begin{definition}[Centroid of a Plane Region]
The \cle{centroid} $(\bar{x},\bar{y})$ of a plane region $R$ of area $A$ is given by
\[
\bar{x} = \frac{1}{A}\iint_{R} x\;dA,\qquad \bar{y} = \frac{1}{A}\iint_{R} y\;dA.
\]
For a region between curves $y=f(x)$ (top) and $y=g(x)$ (bottom) from $x=a$ to $x=b$:
\[
A = \int_{a}^{b} \bigl(f(x)-g(x)\bigr)\;dx,\quad
\bar{x} = \frac{1}{A}\int_{a}^{b} x\bigl(f(x)-g(x)\bigr)\;dx,\quad
\bar{y} = \frac{1}{A}\int_{a}^{b} \frac{f(x)+g(x)}{2}\bigl(f(x)-g(x)\bigr)\;dx = \frac{1}{2A}\int_{a}^{b} \bigl(f(x)^{2}-g(x)^{2}\bigr)\;dx.
\]
\end{definition}

\begin{definition}[Centre of Gravity of a Solid of Revolution]
For a solid of revolution of volume $V$ generated by rotating a region about the $x$-axis, the centre of gravity lies on the axis of symmetry. Its $x$-coordinate is
\[
\bar{x} = \frac{1}{V}\int x\;dV = \frac{\int_{a}^{b} x\,\pi y^{2}\;dx}{\int_{a}^{b} \pi y^{2}\;dx}.
\]
For a thin uniform shell (surface of revolution), the centroid height $\bar{y}$ (for rotation about $y$-axis) is
\[
\bar{y} = \frac{\int y\;dS}{\int dS} = \frac{\int 2\pi x\,y\;ds}{\int 2\pi x\;ds}.
\]
\end{definition}

\begin{exemplebox}[Centroid of a Semicircular Lamina]
Find the centroid of a uniform semicircular lamina of radius $r$ (bounded by $y=\sqrt{r^{2}-x^{2}}$ and $y=0$).
By symmetry, $\bar{x}=0$. Area $A = \frac{1}{2}\pi r^{2}$.
\[
\bar{y} = \frac{1}{2A}\int_{-r}^{r} (r^{2}-x^{2})\;dx = \frac{1}{\pi r^{2}}\left[r^{2}x-\frac{x^{3}}{3}\right]_{-r}^{r} = \frac{1}{\pi r^{2}}\cdot\frac{4r^{3}}{3} = \frac{4r}{3\pi}.
\]
The centroid is at distance $\dfrac{4r}{3\pi}$ from the diameter.
\end{exemplebox}

\begin{exemplebox}[Centre of Gravity of a Solid Hemisphere]
A solid hemisphere is generated by rotating the quarter circle $x^{2}+y^{2}=r^{2}$ ($x\ge 0, y\ge 0$) about the $x$-axis.
Volume $V = \pi\int_{0}^{r} (r^{2}-x^{2})\;dx = \frac{2}{3}\pi r^{3}$.
Moment about $x=0$: $M = \pi\int_{0}^{r} x(r^{2}-x^{2})\;dx = \pi\left[\frac{r^{2}x^{2}}{2}-\frac{x^{4}}{4}\right]_{0}^{r} = \frac{\pi r^{4}}{4}$.
Centre of gravity: $\bar{x} = \dfrac{M}{V} = \dfrac{\pi r^{4}/4}{2\pi r^{3}/3} = \dfrac{3r}{8}$ from the flat face.
\end{exemplebox}

\courssub{Theorems of Pappus and Synthesis}

\begin{propriete}[Pappus's First Theorem (Surface Area)]
If a plane curve (which does not cross the axis) is rotated about an external axis in its plane, the area of the surface of revolution generated equals the product of the arc length of the curve and the distance travelled by its centroid.
\[
S = (\text{arc length}) \times (2\pi \times \text{distance from centroid to axis}).
\]
\end{propriete}

\begin{propriete}[Pappus's Second Theorem (Volume)]
If a plane region (which does not cross the axis) is rotated about an external axis in its plane, the volume of the solid of revolution generated equals the product of the area of the region and the distance travelled by its centroid.
\[
V = (\text{area of region}) \times (2\pi \times \text{distance from centroid to axis}).
\]
\end{propriete}

\begin{attention}[Hypotheses for Pappus's Theorems]
The axis of rotation must \emph{not intersect} the interior of the curve (first theorem) or the region (second theorem). The curve/region must lie entirely on one side of the axis. If the axis cuts through the region, the theorems do not apply directly because the centroid's path would not be a simple circle, and the generated solid would self-intersect.
\end{attention}

\begin{exemplebox}[Torus via Pappus]
The circle $x^{2}+(y-3)^{2}=1$ is rotated about the $x$-axis.
\begin{itemize}
\item Area of circle $= \pi(1)^{2} = \pi$. Centroid at $(0,3)$, distance to axis $= 3$.
\item Volume (2nd theorem): $V = \pi \times (2\pi\times 3) = 6\pi^{2}$.
\item Arc length (circumference) $= 2\pi(1) = 2\pi$. Centroid of the curve is also at $(0,3)$.
\item Surface area (1st theorem): $S = 2\pi \times (2\pi\times 3) = 12\pi^{2}$.
\end{itemize}
\end{exemplebox}

\begin{savaistu}[Pappus and the Sphere]
Rotating a semicircle of radius $r$ about its diameter generates a sphere. Area of semicircle $= \frac{1}{2}\pi r^{2}$. Volume of sphere $= \frac{4}{3}\pi r^{3}$. By Pappus's second theorem:
\[
\frac{4}{3}\pi r^{3} = \frac{1}{2}\pi r^{2} \times 2\pi d \quad\Rightarrow\quad d = \frac{4r}{3\pi},
\]
where $d$ is the distance of the semicircle's centroid from the diameter. This is a classic way to find centroids without integration!
\end{savaistu}

\courssub{I check my knowledge}

\begin{exercice}[Multiple choice questions]
For each question, choose the single correct answer, justifying your choice briefly.
\begin{enumerate}
\item The mean value of $f(x)=x^{2}$ over $[0,3]$ is:
\begin{enumerate}
\item $3$ \quad (b) $9$ \quad (c) $\dfrac{9}{2}$ \quad (d) $27$
\end{enumerate}
\item The length of the arc of $y=x$ from $x=0$ to $x=1$ is:
\begin{enumerate}
\item $1$ \quad (b) $\sqrt{2}$ \quad (c) $\dfrac{\sqrt{2}}{2}$ \quad (d) $2$
\end{enumerate}
\item The volume generated by rotating $y=x^{2}$ from $x=0$ to $x=1$ about the $x$-axis is:
\begin{enumerate}
\item $\pi/5$ \quad (b) $\pi/3$ \quad (c) $\pi/2$ \quad (d) $\pi$
\end{enumerate}
\item The centroid of a semicircular lamina of radius $r$ lies
\begin{enumerate}
\item at the centre of the full circle
\item at a distance $r/2$ from the diameter
\item at a distance $4r/(3\pi)$ from the diameter
\item at a distance $2r/\pi$ from the diameter
\end{enumerate}
\end{enumerate}
\end{exercice}

\newpage

\coursec{Solutions to the exercises}

\begin{corrige}[Exercise 1 — Multiple choice questions]
\textbf{Question 1.} Mean value of $f(x)=x^{2}$ over $[0,3]$.

By definition,
\[
\overline{f} = \frac{1}{3-0}\int_{0}^{3} x^{2}\;dx = \frac{1}{3}\left[\frac{x^{3}}{3}\right]_{0}^{3} = \frac{1}{3}\times\frac{27}{3} = \frac{27}{9} = 3.
\]
\textbf{Correct answer: (a) $3$.}

\textbf{Examiner's expectation:} the candidate must write the definition $\overline{f}=\frac{1}{b-a}\int_{a}^{b}f(x)\,dx$ explicitly before substituting; a bare answer without the integral earns no method mark. A frequent error is to compute only $\int_{0}^{3}x^{2}\,dx=9$ and forget the division by $(b-a)$ — this leads to the distractor (b).

\textbf{Question 2.} Length of the arc of $y=x$ from $x=0$ to $x=1$.

Here $\dfrac{dy}{dx}=1$, so
\[
L = \int_{0}^{1}\sqrt{1+1^{2}}\;dx = \int_{0}^{1}\sqrt{2}\;dx = \sqrt{2}\,\bigl[x\bigr]_{0}^{1} = \sqrt{2}.
\]
\textbf{Correct answer: (b) $\sqrt{2}$.}

\textbf{Check (geometrical):} the arc is the straight segment from $(0,0)$ to $(1,1)$, whose Euclidean length is $\sqrt{1^{2}+1^{2}}=\sqrt{2}$ — the formula is consistent. The distractor (a) corresponds to forgetting the square root (integrating $dx$ only).

\textbf{Question 3.} Volume generated by rotating $y=x^{2}$ from $x=0$ to $x=1$ about the $x$-axis.

By the disc method,
\[
V = \pi\int_{0}^{1} \bigl(x^{2}\bigr)^{2}\;dx = \pi\int_{0}^{1} x^{4}\;dx = \pi\left[\frac{x^{5}}{5}\right]_{0}^{1} = \frac{\pi}{5}.
\]
\textbf{Correct answer: (a) $\pi/5$.}

\textbf{Examiner's expectation:} the candidate must square the function \emph{before} integrating: $y^{2}=x^{4}$, not $x^{2}$. Integrating $\pi\int_{0}^{1}x^{2}\,dx$ gives the distractor (b) $\pi/3$, the classic trap of this question.

\textbf{Question 4.} Centroid of a semicircular lamina of radius $r$.

As shown in the worked example of the course, by symmetry $\bar{x}=0$ and
\[
\bar{y} = \frac{1}{2A}\int_{-r}^{r}\bigl(r^{2}-x^{2}\bigr)\;dx
= \frac{1}{\pi r^{2}}\cdot\frac{4r^{3}}{3} = \frac{4r}{3\pi},
\]
measured from the diameter. (Recall $A=\tfrac12\pi r^{2}$, so $2A=\pi r^{2}$.)
\textbf{Correct answer: (c) at a distance $4r/(3\pi)$ from the diameter.}

\textbf{Remark:} the result can also be recovered \emph{without integration} by Pappus's second theorem applied to the rotation of the semicircle about its diameter, which generates a sphere:
\[
\frac{4}{3}\pi r^{3} = \frac{1}{2}\pi r^{2}\times 2\pi d \;\Longrightarrow\; d=\frac{4r}{3\pi}.
\]
The distractor (d) $2r/\pi$ is the distance of the centroid of a semicircular \emph{arc} (a wire), not of the lamina — candidates must distinguish carefully between the two situations.
\end{corrige}

\newpage

\coursec{Summary sheet}

\begin{popfigure}
\begin{tikzpicture}[
  mindmap,
  grow cyclic,
  every node/.style={concept, font=\small},
  concept color=popblue,
  root concept/.append style={concept color=popdark, font=\bfseries\small, text=white},
  level 1/.append style={level distance=3.4cm, sibling angle=51.4},
  level 1 concept/.append style={font=\footnotesize}
]
\node[root concept]{Applications of Definite Integrals}
  child[concept color=popblue]{ node {Mean value $\bar{f}$ and RMS value} }
  child[concept color=popgreen]{ node {Arc length $L=\int_a^b\sqrt{1+(y')^2}\,dx$} }
  child[concept color=poporange]{ node {Surface of revolution $S=2\pi\int y\,ds$} }
  child[concept color=popgold]{ node {Volume of revolution $V=\pi\int y^2\,dx$} }
  child[concept color=poppurple]{ node {Centroid of a plane region} }
  child[concept color=poppink]{ node {Centre of gravity of a solid of revolution} }
  child[concept color=popgreenL]{ node {Theorems of Pappus} };
\end{tikzpicture}
\legende{Mind map of the chapter: seven key ideas built on the definite integral.}
\end{popfigure}

\begin{bilan}
\textbf{Key definitions.}
\begin{itemize}
\item \cle{Mean value} of $f$ on $[a,b]$: $\bar{f}=\dfrac{1}{b-a}\displaystyle\int_a^b f(x)\,dx$. It is the height of the rectangle on $[a,b]$ having the same area as the region under the curve.
\item \cle{Root mean square (RMS) value}: $f_{\mathrm{rms}}=\sqrt{\dfrac{1}{b-a}\displaystyle\int_a^b [f(x)]^2\,dx}$. Always $f_{\mathrm{rms}}\geq |\bar{f}|$; essential for alternating currents ($I_{\mathrm{rms}}=I_0/\sqrt{2}$ for $I=I_0\sin\omega t$).
\item \cle{Arc length element}: $ds=\sqrt{dx^2+dy^2}$, so $ds=\sqrt{1+(y')^2}\,dx$ (Cartesian), $ds=\sqrt{\dot{x}^2+\dot{y}^2}\,dt$ (parametric), $ds=\sqrt{r^2+(dr/d\theta)^2}\,d\theta$ (polar).
\item \cle{Centroid} $G(\bar{x},\bar{y})$ of a plane region: the point where the region would balance; for a uniform lamina it coincides with the centre of mass.
\end{itemize}

\textbf{Formulas to know by heart.}
\begin{itemize}
\item Arc length: $L=\displaystyle\int_a^b\sqrt{1+\left(\tfrac{dy}{dx}\right)^2}\,dx=\int_{t_1}^{t_2}\sqrt{\dot{x}^2+\dot{y}^2}\,dt$.
\item Volume about the $x$-axis (discs): $V=\pi\displaystyle\int_a^b y^2\,dx$; about the $y$-axis: $V=\pi\displaystyle\int_c^d x^2\,dy$; between two curves: $V=\pi\displaystyle\int_a^b\big(y_1^2-y_2^2\big)\,dx$ (washers). Shells about the $y$-axis: $V=2\pi\displaystyle\int_a^b x\,y\,dx$.
\item Surface of revolution about the $x$-axis: $S=2\pi\displaystyle\int_a^b y\,ds$; about the $y$-axis: $S=2\pi\displaystyle\int_a^b x\,ds$, with the appropriate $ds$.
\item Centroid of the region under $y=f(x)$, $a\leq x\leq b$, of area $A$:
\[\bar{x}=\frac{1}{A}\int_a^b x\,y\,dx,\qquad \bar{y}=\frac{1}{2A}\int_a^b y^2\,dx.\]
Between curves $y_1\geq y_2$: $\bar{x}=\dfrac{1}{A}\displaystyle\int_a^b x(y_1-y_2)\,dx$, $\bar{y}=\dfrac{1}{2A}\displaystyle\int_a^b (y_1^2-y_2^2)\,dx$.
\item Centre of gravity of a solid of revolution about the $x$-axis (on the axis by symmetry):
\[\bar{x}=\frac{\displaystyle\int_a^b x\,y^2\,dx}{\displaystyle\int_a^b y^2\,dx}=\frac{1}{V}\int_a^b \pi x y^2\,dx.\]
\item \cle{Pappus I} (volumes): $V=2\pi\,\bar{r}\,A$ = (area) $\times$ (distance travelled by its centroid).
\item \cle{Pappus II} (surfaces): $S=2\pi\,\bar{r}\,L$ = (arc length) $\times$ (distance travelled by the centroid of the arc).
\end{itemize}

\textbf{Methods.}
\begin{itemize}
\item For arc length and surface integrals, first compute $y'$, then simplify $1+(y')^2$ \emph{before} integrating; look for a perfect square (e.g.\ with $y=\cosh x$, $1+\sinh^2x=\cosh^2x$).
\item For volumes, sketch the region, identify the axis of rotation, choose discs/washers (integrate perpendicular to the axis) or shells (parallel to the axis).
\item For centroids, exploit \cle{symmetry} first: if the region is symmetric about an axis, the centroid lies on it. Use composite-body algebra: moments add, so $\bar{x}=\dfrac{\sum A_i\bar{x}_i}{\sum A_i}$ (holes count negatively).
\item Pappus works both ways: use it to find $V$ or $S$ from a known centroid, or to \emph{find a centroid} from a known volume (e.g.\ semicircular lamina: $\bar{y}=\dfrac{4r}{3\pi}$).
\end{itemize}

\textbf{Common mistakes.}
\begin{itemize}
\item Forgetting the $\pi$ in $V=\pi\int y^2dx$, or squaring wrongly: it is $y^2$, not $y$.
\item Using $dx$ instead of $ds$ in the surface formula $S=2\pi\int y\,ds$.
\item Mixing the axes: rotating about the $y$-axis requires $x$ (not $y$) as the radius.
\item Confusing mean value and RMS: the RMS involves the \emph{square inside} the integral, then a square root.
\item In Pappus, the generating region or arc must \emph{not cross} the axis of rotation.
\item For $\bar{y}$ of a lamina, forgetting the factor $\tfrac12$: $\bar{y}=\frac{1}{2A}\int y^2dx$.
\end{itemize}

\textbf{GCE tips.}
\begin{itemize}
\item State the formula before substituting; examiners award method marks for a correctly quoted formula and correct limits.
\item Keep answers exact (in terms of $\pi$, $\ln 2$, etc.) unless a decimal is requested; give units cubed for volumes, squared for areas.
\item Standard results save time: semicircle arc centroid $\frac{2r}{\pi}$ from the diameter; semicircular lamina $\frac{4r}{3\pi}$; hemisphere $\bar{x}=\frac{3r}{8}$ from the flat face; cone $\bar{x}=\frac{h}{4}$ from the base ($\frac{3h}{4}$ from the apex).
\item Check plausibility: a centroid must lie \emph{inside} the region; a volume of revolution must be positive and smaller than the circumscribing cylinder.
\end{itemize}
\end{bilan}

\findecours

\end{document}
